\documentclass[11pt]{article}
\usepackage[T1]{fontenc}
\usepackage{lmodern,amsmath,amssymb,amsthm,booktabs,longtable,array,float,microtype}
\usepackage[margin=1in]{geometry}
\usepackage[colorlinks=true,linkcolor=blue,citecolor=blue,urlcolor=blue]{hyperref}
\hypersetup{pdftitle={Locally Irregular Two-Edge-Colourings of Doubled Graphs: Constructions, Extensions, and Obstructions},pdfauthor={Qiyuan Ma},pdfsubject={Subcubic proof and further conditional constructions; unrestricted conjecture remains open}}
\newtheorem{theorem}{Theorem}[section]
\newtheorem{lemma}[theorem]{Lemma}
\newtheorem{proposition}[theorem]{Proposition}
\newtheorem{corollary}[theorem]{Corollary}
\theoremstyle{definition}
\newtheorem{definition}[theorem]{Definition}
\newtheorem{remark}[theorem]{Remark}
\newcommand{\doub}[1]{{}^{2}\!#1}
\newcommand{\lir}{\operatorname{lir}}
\newcommand{\safe}{\mathcal R}
\newcommand{\SC}{\mathsf S}
\newcommand{\MC}{\mathsf M}
\newcommand{\F}{\mathbb F_3}
\newcommand{\code}[1]{\nolinkurl{#1}}
\newcommand{\per}{\operatorname{per}}
\newcommand{\supp}{\operatorname{supp}}
\makeatletter
\newcommand{\compactcontents}{\begingroup\small
\renewcommand{\l@section}{\@dottedtocline{1}{0pt}{2em}}
\tableofcontents\endgroup}
\makeatother
\setlength{\emergencystretch}{3em}
\renewcommand{\labelitemi}{\(\bullet\)}
\title{Locally Irregular Two-Edge-Colourings of Doubled Graphs:\\Constructions, Extensions, and Obstructions}
\author{Qiyuan Ma\\Shanghai Jiao Tong University\\{\small\href{mailto:mqyqwerty@sjtu.edu.cn}{\texttt{mqyqwerty@sjtu.edu.cn}}}\\{\small ORCID: \href{https://orcid.org/0009-0000-9532-0821}{\texttt{0009-0000-9532-0821}}}}
\date{28 October 2026}
\begin{document}
\maketitle
\begin{abstract}
We study the conjecture that replacing every edge of a connected simple
graph other than $K_2$ by two parallel copies produces a multigraph whose
edges admit a locally irregular two-colouring. We present a
computer-assisted proof for all subcubic underlying graphs and develop
constructions for several families of unbounded maximum degree. These
include sparse direct branch cores, degree-homogeneous three-colourable
cores joined by paths, balanced partial subdivisions, parity-separated
degree patterns, and clique block trees with explicit interface capacity.
Two general sufficient conditions are
$\delta(G)\ge8\chi(G)(2\chi(G)+1)-7$ and
$8\delta(G)\ge7|V(G)|+6$; the latter admits red degrees distinct at all
vertices and blue degrees distinct at all vertices. We extend the dense
construction to prescribed external contributions, forbidden final
degrees, and sequential absorption of multiple regions. Complement-path
lifts give further infinite families. Exact fixed-target factor criteria,
exchange and polynomial identities, and explicit obstructions clarify
the limits of several stronger strategies. Complete finite interface
certificates and executable checkers accompany the proofs. Classical
factor, matching and interpolation tools are explicitly attributed.
We also prove a signed-gap lemma with at most two occurrences per
variable, use it to colour degree-three/four path--cycle partitions,
and derive exact fixed-interface criteria and arbitrary-degree
root-port constructions for odd components.
The conjecture for arbitrary connected graphs is not proved.
\end{abstract}
\clearpage
\compactcontents
\clearpage
\section{Problem, scope, and provenance}\label{sec:overview}

For a finite simple graph $G$, let $\doub G$ denote the multigraph in
which every original edge has two parallel copies. A colour class is
\emph{locally irregular} if the degrees of the endpoints of each of its
edges are different in that colour class. Parallel copies may have
different colours. Isolated vertices cause no difficulty. The conjecture
of Grzelec and Wo\'zniak~\cite{GW23} asks whether
\[
 G\text{ connected},\quad G\not\cong K_2
 \quad\Longrightarrow\quad \lir(\doub G)\le2.
\]
All statements below concern finite graphs. Original graphs are simple;
auxiliary connection graphs may have labelled parallel edges and loops.
A loop in a connection graph represents a return path in the original
simple graph, rather than an original loop.

The basic language is an edge weight $x_e\in\{0,1,2\}$ equal to the
number of red copies. The actual colour degrees are
$R(v)=\sum_{e\ni v}x_e$ and $B(v)=2d_G(v)-R(v)$. A comparison is required
in a colour precisely when that colour is present on the edge. Requiring
both comparisons on every edge is a stronger sufficient condition that
some, but not all, of our constructions satisfy.

The original conjecture and results for bipartite graphs appear in
\cite{GW23}; cacti are treated in~\cite{GW24}. The regular-graph result,
a three-colour bound for subcubic graphs, and special subcubic families
are established in~\cite{GOW26}. In its arXiv version, Theorems~14--15
cover independent degree-three vertices and subdivisions of selected
cubic edges into paths of length at least five. We give a full subcubic
proof, retaining its complete finite interfaces. The later sections
address different questions: unbounded-degree constructions, prescribed
boundaries, target feasibility, and limitations of proposed proof methods.

The identifiers A1--A9, B1--B4 and C1--C5 refer to a prior inventory of
research notes; they are preserved for traceability in
Table~\ref{tab:coverage}. They are groups of related assertions, not a
count of mutually independent or historically new theorems. We fold
ordinary existence results for smaller subcubic families into A1 and
retain stronger boundary statements separately. D1--D5 identify completed
extensions documented after 17:00 China Standard Time on 27 October
2026, with their source snapshot fixed at 20:31:04 that evening.
E1--E4 contain four completed notes from 28 October 2026, frozen at
10:25:33 China Standard Time. They add integer rather than merely
modular separation, path/cycle completions, and parity-port certificates.
Work completed after that cutoff, including the later component-dynamic-programming
note, is outside this fixed snapshot.
The original notes, snapshot hashes and a detailed coverage ledger
accompany the manuscript.

\begin{longtable}{>{\raggedright\arraybackslash}p{0.10\linewidth}>{\raggedright\arraybackslash}p{0.53\linewidth}>{\raggedright\arraybackslash}p{0.27\linewidth}}
\caption{Coverage of the retained inventory.}\label{tab:coverage}\\
\toprule Identifier&Assertion group&Location\\\midrule\endfirsthead
\toprule Identifier&Assertion group&Location\\\midrule\endhead
A1&All subcubic graphs; safe interfaces and the two-mark kernel&
Sections~\ref{sec:subcubic}--\ref{sec:complete}, Appendix~\ref{app:cert}\\
A2&Three sparse direct-branch hypotheses&Section~\ref{sec:branches}\\
A3&Degree-homogeneous, three-colourable direct cores&Section~\ref{sec:homogeneous}\\
A4--A5&Balanced half-edges and simultaneous preserved-colour extensions&Section~\ref{sec:balanced}\\
A6&Near-centre parity realisation and degree-gap separation&Section~\ref{sec:parity}\\
A7--A8&Arbitrary clique root sums, block trees, and capacity extensions&Section~\ref{sec:cliques}\\
A9&Tree gluing of given coloured pieces&Section~\ref{sec:pieceglue}\\
B1&Trees, paths, decorated saturation, preserved subdivisions, offsets&Sections~\ref{sec:generalports}, \ref{sec:trees}, \ref{sec:preserve}, \ref{sec:oddforest}\\
B2&Fixed targets, matching, factors, exchanges, defect bounds&Section~\ref{sec:targets}\\
B3&Exact polynomial, free parameters, permanent and parity formula&Section~\ref{sec:algebra}\\
B4&Normal forms, peripheral elimination and minimal counterexamples&Section~\ref{sec:reductions}\\
C1&Cut, parity, modulo-three and modulo-two limitations&Sections~\ref{sec:mod2limits}, \ref{sec:templates}\\
C2&Nonincreasing single-edge and two-unit descent barriers&Section~\ref{sec:descent}\\
C3&Target obstructions, nonconvexity and unbounded parity radius&Sections~\ref{sec:targetbarriers}, \ref{sec:radius}\\
C4&Sharp fixed-interface path and extension thresholds&Sections~\ref{sec:treecapacity}, \ref{sec:cliquecapacity}, \ref{sec:interfacebarriers}\\
C5&Gluing, prescribed root sums and multiple marked ports&Sections~\ref{sec:piecebarrier}, \ref{sec:interfacebarriers}\\
D1&High minimum degree relative to chromatic number&Section~\ref{sec:highdegree}\\
D2&Multiple CRT regions, low-side ports, matched tree attachments&Sections~\ref{sec:crtregions}, \ref{sec:irregularcore}\\
D3&Dense antitone targets: $8\delta\ge7n+6$&Section~\ref{sec:dense}\\
D4&Complement lifts, nonregular families, structural restrictions&Section~\ref{sec:complements}\\
D5&Dense forbidden targets and sequential region absorption&Section~\ref{sec:denselists}\\
E1&Signed gaps with two occurrences; paired partitions and fixed-mixed integer repairs&Section~\ref{sec:signedgap}\\
E2&Degree-three/four path--cycle side against a covering matching&Section~\ref{sec:pathcyclepart}\\
E3&Protected parity joins and exact independent-side completions&Section~\ref{sec:isolatepart}\\
E4&Arbitrary-degree odd-component root ports and capacity-aware parity test&Section~\ref{sec:rootports}\\
\bottomrule
\end{longtable}

The matching reductions, Tutte's factor criterion, jump-system exchange
principles and polynomial interpolation are classical tools, not claims
of invention. We give the required applications and elementary arguments
so that the boundary between a feasible target and a legal colouring is
explicit. The degree-constrained subgraph theorem used for D1 is quoted
with its hypotheses and applied only to a simple host graph. The
subcubic proof uses the degree-choosability theorem recorded
in~\cite{CR17}. Other background dependencies are indicated where used.

This is a research preprint with extensive computer assistance. Its
infinite-family assertions are justified by the displayed arguments and,
where indicated, complete finite kernels plus lifting inductions.
Random and bounded-order checks test implementations; they do not prove
the conjecture. No independent expert review or proof-assistant
verification is claimed, and no absolute historical-priority claim is
made. In particular, none of the strategy obstructions below is a
counterexample to the original conjecture.

\section{The subcubic theorem}\label{sec:subcubic}
\begin{theorem}\label{thm:main}
Let $G$ be a connected simple graph with $G\not\cong K_2$ and
$\Delta(G)\le3$.  Then $\lir(\doub G)\le2$.
\end{theorem}

The proof uses finite calculations for three auxiliary state systems.
Their complete domains, recurrences and checking procedures are specified
below and in Appendix~\ref{app:cert}.  The finite calculations are not
enumerations of all graphs up to an order bound.  Inductions on paths,
trees and blocks extend the verified relations to graphs of arbitrary
order.  The only external colouring theorem used is the classical
degree-list theorem, stated as Theorem~A in \cite{CR17}.

\section{Edge weights and path interfaces}\label{sec:path}

Assign to every original edge $e$ a weight $x_e\in\{0,1,2\}$, the number
of its red copies.  Define
\begin{equation}\label{eq:actual}
 R(v)=\sum_{e\ni v}x_e,\qquad B(v)=2d_G(v)-R(v).
\end{equation}
The required condition on an edge $uv$ of weight $x$ is
\begin{equation}\label{eq:rule}
 x>0\Longrightarrow R(u)\ne R(v),\qquad
 x<2\Longrightarrow B(u)\ne B(v).
\end{equation}
When both endpoints have degree three, it is sufficient and necessary
that their red degrees differ.  The same observation applies to two
degree-two vertices, with $B=4-R$ instead of $6-R$.

Let $W=\{v:d_G(v)=3\}$ and $H=G[W]$.  Vertices of $W$ are called
\emph{roots}; a connected component of $H$ is called a \emph{core}.
Every edge outside $H$ belongs to a maximal path with internal vertices of
degree two.  Such a path is an \emph{ordinary connection} if its endpoints
are distinct roots, a \emph{return path} if both endpoints are the same
root, and a \emph{pendant path} if one endpoint is a root and the other is
a leaf.  An ordinary connection has length at least two, and a return
path has length at least three.  A pendant path can have any positive
length.  Distinct roots of an ordinary connection may belong to the same
core.  A length-one edge between roots is already an edge of $H$.

The weight of the first edge at a root is the corresponding \emph{port}.
A port is \emph{pure} if its weight is 0 or 2, and \emph{mixed} if its
weight is 1.  A root red degree is physically compatible with a port $a$
when $a\le R\le a+4$.  At a return path its two port contributions must
both be included in the same root degree.

\begin{lemma}[Path interfaces]\label{lem:path}
The following statements hold for physically compatible root degrees.
\begin{enumerate}
\item A two-edge ordinary connection with pure end ports $a,b$ is
admissible exactly when $R(u)\ne2+b$ and $R(v)\ne2+a$.
\item A three-edge ordinary connection with arbitrary ports $a,b$ is
admissible exactly when
\begin{equation}\label{eq:three}
 a\ne b,\qquad
 \neg\bigl((a=1\text{ or }b=1)\text{ and }R(u)=R(v)=3\bigr).
\end{equation}
\item Opposite pure ports are admissible on every ordinary connection of
length at least three.  Any two pure ports are admissible at length at
least four.
\item A return path accepts ports $(0,2)$ at every legal length.  If its
root red degree is nonzero modulo three, each of $(0,1)$ and $(1,2)$ is
also admissible.
\item A pendant path of length at least three accepts either pure root
port.  At length one, ports 0 and 2 exclude root red degrees 4 and 2,
respectively; at length two they exclude 2 and 4, respectively.
\end{enumerate}
\end{lemma}
\begin{proof}
For a path with weights $x_1,\ldots,x_\ell$, all comparisons between its
internal degree-two vertices are equivalent to
$x_i\ne x_{i+2}$.  At a degree-three root of red degree $r$, its incident
port $a$ permits the next weight to lie in
\begin{equation}\label{eq:P}
 P_r(a)=\{0,1,2\}\setminus
 \bigl(\{r-a:a>0\}\cup\{r-a-2:a<2\}\bigr),
\end{equation}
where values outside $\{0,1,2\}$ have no effect.
Part~1 follows by direct substitution.  For three edges, the middle
comparison is $a\ne b$ and the remaining weight must lie in
$P_{R(u)}(a)\cap P_{R(v)}(b)$.  A pure port has at least two choices.
A mixed port has at least two choices except at red degree 3, where its
set is $\{1\}$.  A pure-port set excludes 1 exactly at red degree 3.
This proves part~2.

For longer paths, separate the odd and even edge positions.  They form
two three-colour paths under the relations $x_i\ne x_{i+2}$.
At even length at least four, each position path has one fixed endpoint
and an endpoint list of size at least two.  At odd length at least five,
one position path has two fixed endpoints at distance at least two and
the other has two lists of size at least two.  Both can always be coloured.
Together with part~2 this proves part~3.  If one port is mixed but its
root red degree is not 3, its set in \eqref{eq:P} again has at least two
elements, so the same argument applies to longer paths.

For a return path set both endpoint red degrees equal to the actual
root degree.  Its nonzero residue excludes red degree 3; the preceding
argument gives part~4.  At a pure leaf port 0 or 2, the preceding weight
has allowed set $\{0,1\}$ or $\{1,2\}$, respectively.  The position-path
argument proves part~5 for length at least four.  At length three choose
the leaf port opposite to the root port and intersect two sets of size at
least two.  The length-one and length-two exclusions follow directly
from \eqref{eq:rule}.
\end{proof}

We will use the safe set
\begin{equation}\label{eq:safe}
 \safe=\{0,1,3,5,6\}.
\end{equation}
In particular, pure ports at safe roots automatically satisfy either
end of every two-edge connection.  To attach a pendant path at a modular
root of target residue $\phi$, take port 2 exactly when its length is one
and $\phi=1$, or its length is two and $\phi=2$; take port 0 otherwise.
Both choices are allowed at any length when $\phi=0$.

If $W$ is empty, $G$ is a single vertex, a path, or a cycle.
A two-edge path takes weights $(0,0)$; longer paths take first weight 0
and last weight 2 and use the position-path argument with leaf endpoint
conditions.  For a cycle, give its edge positions a proper three-colouring
in the auxiliary graph joining positions $i$ and $i+2$ modulo its length.
That auxiliary graph has maximum degree two.  The resulting weights
satisfy all comparisons.  Also, for $K_4$, on edges
$(01,02,03,12,13,23)$ take weights $(1,2,0,0,0,2)$; the red degrees are
$(3,1,4,2)$.  We henceforth exclude these cases.

\section{Three finite auxiliary lemmas}\label{sec:finite}

A core $C$ has $3-d_C(v)$ private ports at each vertex $v$.
In an auxiliary port assignment its red degree is the sum of its core
edge weights and its private port weights.  Private ports here are free
variables, not prescribed boundary weights.

\begin{lemma}[Free safe ports]\label{lem:free}
Suppose $C$ is a tree, a cactus all of whose cycle blocks are odd, or a
connected bipartite subcubic graph with at most two vertices of core
degree three.  There is an assignment for which all private ports are
pure, every red degree belongs to $\safe$, and adjacent core vertices
have distinct red degrees.
\end{lemma}

\begin{lemma}[One mixed port]\label{lem:one}
Let $C$ be a bipartite subcubic cactus.  Any specified private port can
be set to 1 so that its root red degree is 1 or 5, all other private
ports are pure, every root red degree is safe, and adjacent core vertices
have distinct red degrees.
\end{lemma}

We specify the finite calculations and their infinite extensions before
stating the third lemma.  A rooted state $(x,r)$ records the incoming
core-edge weight (or one actual free port at the final root) and the red
degree of that vertex.  Its domain is
\begin{equation}\label{eq:Omega}
 \Omega=\{(x,r):x\in\{0,1,2\},\ r\in\safe,\ x\le r\le x+4\},
 \qquad |\Omega|=9.
\end{equation}
The six rooted message sets are
\begin{align*}
 \mathcal L&=\{(0,0),(1,1),(1,3),(1,5),(2,6)\},\\
 \mathcal O&=\{(0,1),(0,3),(1,1),(1,3),(1,5),(2,3),(2,5)\},\\
 \mathcal A&=\Omega,\\
 \mathcal C&=\{(0,0),(0,1),(1,3),(2,5),(2,6)\},\\
 \mathcal D&=\{(0,0),(0,1),(0,3),(1,3),(2,3),(2,5),(2,6)\},\\
 \mathcal E&=\{(0,0),(0,1),(1,1),(1,3),(1,5),(2,5),(2,6)\}.
\end{align*}
For $c\in\{0,1,2\}$ child sets $X_1,\ldots,X_c$, define
\begin{equation}\label{eq:node}
 N(X_1,\ldots,X_c)=
 \left\{\left(x,x+\sum_{i=1}^c y_i+\sum_{j=1}^{2-c}p_j\right):
 \begin{array}{l}
 (y_i,s_i)\in X_i,\ p_j\in\{0,2\},\\
 r=x+\sum_i y_i+\sum_jp_j\in\safe,\ r\ne s_i\ (1\le i\le c)
 \end{array}\right\}.
\end{equation}
Here $x\in\{0,1,2\}$.  Every one of the $1+6+36=43$ ordered
choices has one of the six message sets as its exact result.

For a non-anchor vertex of a cycle, its third incidence is a free pure
port, denoted $F$, or one child message set $X$.  Define a Boolean matrix
$M_X$ on $\Omega$ by
\begin{equation}\label{eq:letter}
 (M_X)_{(x,r),(y,s)}=1
 \quad\Longleftrightarrow\quad
 s\ne r\text{ and }
 \begin{cases}
 s-x-y\in\{0,2\},&X=F,\\
 \exists(p,t)\in X:\ s=x+y+p,\ s\ne t,&X\ne F.
 \end{cases}
\end{equation}
The rooted interface of a cycle product $M$ consists of the states
$(z,r)$ for which there are $a,b\in\{0,1,2\}$ and $(b,s)\in\Omega$
with $r=z+a+b$, $s\ne r$, and $M_{(a,r),(b,s)}=1$.

\begin{proposition}[Finite closure identities]\label{prop:closure}
The node recurrence \eqref{eq:node} is closed on the six message sets.
Let $\mathcal P_k$ be the set of all products of $k$ letter matrices
from \eqref{eq:letter}.  Then
\begin{align*}
 (|\mathcal P_2|,|\mathcal P_4|,|\mathcal P_6|)&=(14,11,2),&
 \mathcal P_8&=\mathcal P_6,\\
 (|\mathcal P_3|,|\mathcal P_5|,|\mathcal P_7|)&=(23,6,2),&
 \mathcal P_9&=\mathcal P_7.
\end{align*}
Their rooted interface types have the following multiplicities:
\begin{center}
\begin{tabular}{ccl}
\toprule
Parity of cycle length&Number of letters&Rooted types\\
\midrule
odd&2&$\mathcal C$ once, $\mathcal D$ thirteen times\\
odd&4&$\mathcal A$ ten times, $\mathcal D$ once\\
odd&6&$\mathcal A$ twice\\
even&3&$\mathcal A$ twenty-one times, $\mathcal O$ twice\\
even&5&$\mathcal A$ six times\\
even&7&$\mathcal A$ twice\\
\bottomrule
\end{tabular}
\end{center}
Moreover the node recurrence restricted to
$\mathcal L,\mathcal O,\mathcal A$ stays within these three types.
\end{proposition}
\begin{proof}[Computer-assisted finite calculation]
Evaluate \eqref{eq:node} for the 43 choices.  Construct the seven
matrices by \eqref{eq:letter} and multiply them as Boolean relations.
The files \code{safe_cactus_certificate.py} and
\code{safe_all_cactus_certificate.py} independently recompute the
complete layers, compare sets rather than just their cardinalities,
and compute every rooted interface.  They use Python integer and set
operations only.  The complete layer sets are supplied in
\code{safe_cactus_interfaces.json} and
\code{safe_all_cactus_interfaces.json}.  Appendix~\ref{app:cert}
gives the common execution command.
\end{proof}

\begin{proof}[Proof of the tree and cactus parts of Lemmas~\ref{lem:free}
and~\ref{lem:one}]
For a tree, root at a vertex with an available private port, use that
port as the incoming incidence, and apply \eqref{eq:node} upward.
Only $\mathcal L,\mathcal O,\mathcal A$ occur.  All accept incoming
weight 0, which gives the free-port assertion, and incoming weight 1
with red degree 1 or 5, which gives the marked assertion.

For a cactus of maximum degree three, different cycle blocks cannot
share a vertex.  Contract each cycle block to a unit and induct on the
resulting tree of units.  A non-anchor cycle vertex has precisely the
third-incidence choices in \eqref{eq:letter}; two children there would
violate the degree bound.  A cycle of length $\ell$ uses $\ell-1$
letters.  The identities in Proposition~\ref{prop:closure} imply, by
induction under multiplication by any two letters, that every longer
product of the same parity has one of the listed rooted types.
Thus the argument covers arbitrary cycle lengths and numbers of blocks.

Every one of the six types accepts incoming weight 0.  This proves the
odd-cactus assertion.  In a bipartite cactus every cycle is even, and
the induction stays in $\mathcal L,\mathcal O,\mathcal A$.  Root at the
specified private port and choose $(1,1)$ or $(1,5)$ to prove
Lemma~\ref{lem:one}.  The incoming port is an actual private incidence,
not an extra edge added to the original degree.
\end{proof}

For core chains with one pure private port at each internal vertex,
write
\begin{equation}\label{eq:T}
 (x,r)\ T\ (y,s)\quad\Longleftrightarrow\quad
 r\ne s,\quad s=x+y+p\in\safe\text{ for some }p\in\{0,2\}.
\end{equation}
Both states lie in $\Omega$.  The supplied certificate gives a six-step
walk for every ordered pair of states, so
\begin{equation}\label{eq:sat}
 T^6=\Omega\times\Omega.
\end{equation}
Every state has an incoming neighbour, and therefore
$T^k=\Omega\times\Omega$ for every $k\ge6$.
Here internal chain vertices have original degree three, unlike the
degree-two internal vertices in Section~\ref{sec:path}.

\begin{proof}[Proof of the remaining part of Lemma~\ref{lem:free}]
If there are no core-degree-three vertices, the core is a single vertex,
a path, or an even cycle, already covered above.  Otherwise compress its
core-degree-two chains.  If $n_1,n_3$ count core degrees one and three,
and $\beta=|E(C)|-|V(C)|+1$, then
\begin{equation}\label{eq:euler}
 n_1=n_3+2-2\beta.
\end{equation}
With one degree-three vertex the only shapes are a three-arm tree and
an even cycle with an arm.  With two, the only shapes are a double-star
tree, an even cycle attached to a fork with two arms, an even cycle with
an arm at each of its two forks, a theta, or two even cycles joined by
a chain.  These follow by counting loops and paths in the compressed
shape, whose two forks have degree three.

There are 163 complete finite shape and length representatives, with
their domain specified in Table~\ref{tab:sparse}.  Every representative
has a witness assigning odd safe red degree to each fork, consistent
with its three incident weights.  The checker
\code{sparse_bip_core_certificate.py} verifies the complete domain,
all fork sums and all chain relations, as well as \eqref{eq:sat}.

Keep a length $n\le6$, replace a longer odd length by 7, and a longer
even length by 8.  Any chain of length at least seven has at least six
internal transitions, so \eqref{eq:sat} lifts its witness with unchanged
end weights and fork degrees.  Parities and simplicity restrictions
are preserved.  An arm can have arbitrary positive length: an odd-red
fork can be completed immediately at a leaf by two pure ports; from
every state one step reaches a state that admits such a completion,
and every state has an outgoing step.  These finite assertions are also
checked.  Thus all arm lengths and all chain lengths are covered.
\end{proof}

\begin{definition}[Four-port tree]\label{def:four}
A four-port tree is a subcubic tree core with two distinguished vertices
$a,b$ at odd distance.  Neither has a short port.  On each side of the
$a$--$b$ path there are exactly two positions, each either a direct long
private port at the distinguished vertex or an arm ending at a leaf
with one long and one short port.  Every other internal vertex of such
an arm or of the $a$--$b$ path has one short private port.  Thus there
are precisely four long ports.  The distinguished vertices have no
other incidences.
\end{definition}

Figure~\ref{fig:four} shows one possible topology.  A position at a
distinguished vertex may instead be a direct long port.
\begin{figure}[htbp]
\centering
\setlength{\unitlength}{0.9mm}
\begin{picture}(140,54)
\put(35,25){\circle*{2.4}}\put(105,25){\circle*{2.4}}
\put(35,25){\line(1,0){70}}
\put(55,25){\circle*{1.5}}\put(85,25){\circle*{1.5}}
\put(55,25){\line(0,1){8}}\put(85,25){\line(0,1){8}}
\put(55,36){\makebox(0,0){$s$}}\put(85,36){\makebox(0,0){$s$}}
\put(35,25){\line(-4,3){20}}\put(35,25){\line(-4,-3){20}}
\put(105,25){\line(4,3){20}}\put(105,25){\line(4,-3){20}}
\put(15,40){\circle*{1.5}}\put(15,10){\circle*{1.5}}
\put(125,40){\circle*{1.5}}\put(125,10){\circle*{1.5}}
\put(15,40){\line(-1,0){13}}\put(15,10){\line(-1,0){13}}
\put(125,40){\line(1,0){13}}\put(125,10){\line(1,0){13}}
\put(15,40){\line(0,1){8}}\put(15,10){\line(0,-1){8}}
\put(125,40){\line(0,1){8}}\put(125,10){\line(0,-1){8}}
\put(15,51){\makebox(0,0){$s$}}\put(15,0){\makebox(0,0){$s$}}
\put(125,51){\makebox(0,0){$s$}}\put(125,0){\makebox(0,0){$s$}}
\put(2,44){\makebox(0,0){$\ell_1$}}\put(2,6){\makebox(0,0){$\ell_2$}}
\put(138,44){\makebox(0,0){$\ell_3$}}\put(138,6){\makebox(0,0){$\ell_4$}}
\put(35,18){\makebox(0,0){$a$}}\put(105,18){\makebox(0,0){$b$}}
\end{picture}
\caption{A four-port tree with two distinguished vertices $a,b$, four
long ports $\ell_i$, and short ports labelled $s$.  The middle path
is odd; the arms and middle path may be longer than shown, with one
short port at every added internal core vertex.  All vertices drawn
have original degree three.  Only $a,b$ have no short incidence and
may use red degrees 2 or 4.}
\label{fig:four}
\end{figure}

\begin{lemma}[Two mixed ports]\label{lem:two}
In any four-port tree, any two distinct long ports can be set to 1 so
that their roots have red degree different from 3.  All other private
ports can be pure, adjacent core vertices can have distinct red
degrees, and every vertex carrying a short port can have safe red
degree.  The two distinguished vertices may have red degree 2 or 4.
\end{lemma}
\begin{proof}
For an arm, its incoming message at the leaf is $\mathcal L$ if the
long port is pure, and
$Q=\{(0,1),(2,5)\}$ if the long port is marked.  At an internal
vertex the message map is
\begin{equation}\label{eq:U}
 U(X)=\{(x,x+y+p):(y,s)\in X,\ p\in\{0,2\},\
 x+y+p\in\safe,\ x+y+p\ne s\}.
\end{equation}
Exact iteration gives cardinalities $5,7,9,9$ from $\mathcal L$, and
$2,3,5,7,7,9,9$ from $Q$.  In each case the last two sets are equal
to $\Omega$, not merely equinumerous.  Thus an unmarked position has
four types, direct or arm lengths $1,2,\ge3$; a marked position has
seven types, direct or arm lengths $1,2,3,4,5,\ge6$.

Up to exchanging the two ends, the mark distribution is $(2,0)$ or
$(1,1)$.  The first has
$\binom82\binom52=280$ choices of unordered pairs of position types;
the second has $(7\cdot4)^2=784$ choices.  Two direct marked ports at
the same distinguished vertex are included.  The connecting path is
odd, and its representatives have lengths $1,3,5,7$.  Hence the complete
domain contains $(280+784)\cdot4=4256$ cases.

For each case, \code{weak_two_marked_interfaces.json} gives the end
states, the two position choices at each end, and compatibility with
the connecting path.  The checker recomputes the domain, arm sets,
actual sums and connecting relations; every case passes.  A second
verifier recomputes existence from these definitions without using the
stored witnesses.  Details appear in Appendix~\ref{app:cert}.

Arm fixed points prove lifting at every larger arm length.  For the
middle path, \eqref{eq:sat} proves lifting at every length at least
seven if the two end degrees are safe.  If an end red degree is 2 or 4,
it differs from every internal safe red degree; a word valid with the
same end-edge weight and virtual end red degree 3 is still valid with
the actual degree.  The virtual state is physical for every possible
end weight.  Thus larger odd lengths preserve both actual end degrees
and end weights.  These observations prove the lemma for arbitrary
four-port trees, rather than only for the finite representatives.
\end{proof}

\section{Classifying cores and counting long ports}\label{sec:class}

Call an ordinary connection \emph{short} if its length is two, and
\emph{long} if its length is at least three.  For a bipartite core $C$
with parts $A,B$, count the following objects on each part:
\begin{enumerate}
\item every pendant port rooted in that part;
\item every endpoint in that part of a short connection whose two roots
both belong to $C$;
\item every vertex in that part incident with neither a pendant path
nor any short connection, including short connections to other cores.
\end{enumerate}
Let the resulting counts be $k_A,k_B$.  In item~3 a core-degree-three
vertex is always counted, even though it has no private port.

Assign cores to classes $\SC$ and $\MC$ by the following fixed rules,
before choosing any weights.  Put in $\MC$ every core with a return
path, every nonbipartite core incident with an ordinary three-edge
connection, and every bipartite core incident with an ordinary
three-edge connection and satisfying $\max(k_A,k_B)\ge2$.
Among the remaining cores, put in $\SC$ those covered by
Lemma~\ref{lem:free}, and put the rest in $\MC$.

The letters specify how a core will be handled, not that all degrees
in its final colouring must have a common residue.  In particular,
Lemma~\ref{lem:two} may give red degree 2 or 4 at vertices of an
$\SC$ core that carry no short port.

\begin{lemma}\label{lem:count}
An $\SC$ core incident with an ordinary three-edge connection is
bipartite and has at most four long ports.  If it has four ports that
belong to three-edge connections between $\SC$ cores, it is a four-port
tree as in Definition~\ref{def:four}.
\end{lemma}
\begin{proof}
Nonbipartite cores incident with such a connection have been assigned
to $\MC$.  For a remaining bipartite core, $k_A,k_B\le1$.
Let $p$ be the number of pendant ports, $s$ the number of short
connections internal to $C$, and $t$ the number of vertices in item~3.
Then
\begin{equation}\label{eq:budget}
 t+p+2s=k_A+k_B\le2.
\end{equation}
There is no return path.  A core leaf has at most one long port, plus
one more if it is counted in $t$.  A core-degree-two vertex with a long
port is counted in $t$.  A core-degree-three vertex is counted in $t$
but has no long port.  Therefore, writing $L$ for the total number of
long ports, \eqref{eq:euler} gives, for a nontrivial core,
\begin{equation}\label{eq:Lbound}
 L\le n_1+t-n_3=2-2\beta+t
 \le4-2\beta-p-2s\le4.
\end{equation}
A one-vertex core has only three ports.

If four of these ports belong to the stated three-edge connections,
all inequalities are equalities.  Thus $\beta=p=s=0$, $t=2$, and the
two vertices counted in $t$ lie in opposite parts.  All core-degree-three
vertices are among them.  Their connecting path has odd length.  Each
other internal core vertex has its one short port; each other core
leaf has one long and one short port.  Removing the middle path leaves
exactly two positions at each end, either direct ports or unbranched
arms.  There are no other long ports, pendant paths or internal short
connections.  This is precisely Definition~\ref{def:four}.
\end{proof}

\begin{corollary}\label{cor:cactus}
Every $\SC$ core incident with an ordinary three-edge connection is a
bipartite cactus.
\end{corollary}
\begin{proof}
Its $L$ is positive, so \eqref{eq:Lbound} gives $\beta\le1$.
A connected graph with at most one cycle is a cactus.
\end{proof}

\section{Making all safe three-edge connections compatible}\label{sec:glue}

Form an indexed multigraph $F$ whose vertices are the $\SC$ cores.
Every ordinary three-edge connection with both ends in $\SC$ gives an
edge of $F$.  Roots in the same core give a loop, and different
connections may give parallel edges.  Loops contribute two to degree.
By Lemma~\ref{lem:count}, $\Delta(F)\le4$.

\begin{lemma}\label{lem:orient}
There is a spanning forest $T$ of $F$ and an orientation of
$D=E(F)\setminus E(T)$ with outdegree at most two at each vertex.
Outdegree two is possible only at a vertex of $F$ of degree four.
\end{lemma}
\begin{proof}
Take any spanning forest.  Orient $D$ so that its outdegree at each
vertex lies between the floor and ceiling of half its degree.  Such
an orientation is obtained by joining each odd-degree vertex to a
temporary vertex, orienting Euler tours in the resulting even-degree
multigraph, and deleting the temporary edges.  A loop has one tail.

In a component of $F$ with at least two vertices, every vertex has an
incident forest edge.  Thus its $D$-degree is at most three, and at
most two if its $F$-degree is at most three.  In a singleton component
all edges are loops, and at most two loops are possible.  These facts
prove both assertions.
\end{proof}

For every $\SC$ core of outdegree zero use Lemma~\ref{lem:free}.
For outdegree one, mark its outgoing half-port by 1 and use
Corollary~\ref{cor:cactus} and Lemma~\ref{lem:one}.
For outdegree two, Lemma~\ref{lem:count} gives a four-port tree, to
which Lemma~\ref{lem:two} applies.  Only long ports are marked.
Every short-incident root remains safe.

Each $D$ edge now has exactly one mixed and one pure port, with red
degree different from 3 at its mixed end.  It is compatible by
Lemma~\ref{lem:path}.  A loop in $F$ uses two distinct actual roots of
the same core; it is not an original return path.

Every $T$ edge still has pure ports $a,b$.  Choose a colour-complement
bit $\eta_C$ for each core.  Complementation changes every core weight
and private port from $x$ to $2-x$ and changes each root red degree
from $r$ to $6-r$.  To make the two ports of a $T$ edge opposite impose
\begin{equation}\label{eq:xor}
 \eta_C\oplus\eta_{C'}=1\oplus(a/2)\oplus(b/2).
\end{equation}
These equations are simultaneously solvable on a forest.
Complementation preserves safety at short ports, unequal red degrees
on core edges, and the condition that a marked root is not of red
degree 3.  Thus all $D$ connections remain compatible.

All ordinary three-edge connections between $\SC$ cores are now
compatible.  At an $\SC$--$\MC$ three-edge connection the $\SC$ port
is pure, since marking was confined to $D$.  Set the $\MC$ port
opposite to it.  At a three-edge connection between two $\MC$ roots
use ports 0 and 2.  At longer ordinary connections any $\MC$ port
can be set to 0.  These choices put no restriction on the eventual
physical red degrees of $\MC$ roots.

\section{Targets and realizability modulo three}\label{sec:mod}

All arithmetic on target values in this section is in $\F$.

\begin{lemma}[Realizing sums]\label{lem:linear}
Let $C$ be a connected simple graph and $q:V(C)\to\F$.
There are weights in $\{0,1,2\}$ whose incident sums equal $q$ modulo
three if and only if, in the bipartite case with parts $A,B$,
\begin{equation}\label{eq:balance}
 \sum_Aq-\sum_Bq=0.
\end{equation}
There is no constraint if $C$ is nonbipartite.
\end{lemma}
\begin{proof}
In a bipartite graph, each edge contributes once to each part, proving
necessity.  Fix a spanning tree.  Setting non-tree weights to zero,
remove leaves successively and set their parent edge to the remaining
target.  The final equation is exactly \eqref{eq:balance}.
For a nonbipartite graph, sign vertices alternately along the tree.
Some non-tree edge has equal signs at its ends.  Its coefficient in
the signed defect is $2$ or $-2$, invertible in $\F$.  Choose its
weight to cancel the defect and then eliminate along the tree.
Choose representatives 0,1,2 for the resulting field elements.
\end{proof}

For an $\MC$ core, let $c(v)$ be the sum of its private ports and seek
a proper target colouring $\phi(v)\in\F$ of its core edges, then
realize the core sums $q(v)=\phi(v)-c(v)$.  The final actual degree
satisfies $R(v)\equiv\phi(v)$, so adjacent core degrees differ.

If a short connection joins $v$ to an $\SC$ root with port $a$, delete
$(2+a)\bmod3$ from the list at $v$.  Denote the resulting list by
$L(v)$.  Since $a$ is pure, only colours 1 or 2 are deleted; hence
\begin{equation}\label{eq:lists}
 0\in L(v),\qquad |L(v)|\ge d_C(v).
\end{equation}
Indeed at most $3-d_C(v)$ private ports can impose deletions.
A return path or a three-edge port at a vertex leaves a strict
surplus $|L(v)|>d_C(v)$ there.

We use the following established theorem \cite[Theorem~A]{CR17}:
a connected graph that is not a Gallai tree is colourable from any
lists of sizes at least the vertex degrees.  A Gallai tree is a
connected graph whose blocks are complete graphs or odd cycles.
Also, a single strict-surplus vertex suffices in any connected graph:
root a spanning tree there and colour in reverse tree order.

\subsection{Nonbipartite cores}

If a nonbipartite $\MC$ core has a return path or an ordinary
three-edge connection, it has a strict-surplus vertex.  Otherwise,
if it were a Gallai tree, its blocks would be bridges, odd cycles,
or $K_4$.  A $K_4$ uses all degree-three incidences, forcing
$G=K_4$, already treated.  Without it the core is an odd-cycle
cactus and would have been assigned to $\SC$.  Thus the remaining
core is not a Gallai tree.  In all cases a proper list target exists,
and Lemma~\ref{lem:linear} realizes it for any subsequent private
offsets.  No bipartite defect needs to be balanced.

\subsection{Bipartite cores with at least two safe short exits}

Let $f$ count short connections from $C$ to $\SC$ cores, with
multiplicity.  If $f\ge2$, choose a proper target from the lists.
This is justified as follows.  If $C$ contains a cycle, its bipartite
structure makes it a non-Gallai graph.  If it is a tree, its assignment
to $\MC$ implies the presence of a return path or a three-edge
connection, giving a strict-surplus vertex.

Reserve this core's own port on each of two safe short exits.  Each
can be set independently to 0 or 2: the opposite root is safe and
therefore accepts either value.  Each changes the defect in
\eqref{eq:balance} by a nonzero element of $\F$, without changing
any target.  Two such binary changes cover all residues and will
balance the defect after other offsets have been fixed.

\subsection{Bipartite cores with at most one safe short exit}

Suppose $f\le1$.  If $C$ has a return path, choose its root in part
$B$, set $\phi=0$ on $A$ and $\phi=t$ on $B$, and choose
$t\in\{1,2\}$ to respect the possible single list restriction.
The return path accepts all three port pairs $(0,1),(0,2),(1,2)$,
because its root target is nonzero.  Their total contributions
$1,2,3$ cover the three defects.

Now suppose there is no return path.  If $C$ was assigned to $\MC$
by the $k$-count rule, some side has at least two objects counted in
$k$.  Any other bipartite $\MC$ core not already covered has at
least three core-degree-three vertices; otherwise
Lemma~\ref{lem:free} would have placed it in $\SC$.  One side
therefore has two such vertices, each counted by item~3.  Again one
side has at least two counted objects.

Take this side as $A$, put $\phi=0$ on $A$, and put
$\phi=t\in\{1,2\}$ on $B$.  Since all lists contain 0 and there
is at most one deletion in total, a choice of $t$ respects all lists.
Two objects on $A$ will provide independent defect adjustments.

\section{Balancing all cores without disturbing short connections}
\label{sec:controls}

First fix the target at every root incident with a short connection.
For a short connection between two $\MC$ roots, give either endpoint
the port
\begin{equation}\label{eq:shortchoice}
 p=\begin{cases}2,&\text{the other root has target }2,\\
                  0,&\text{otherwise}.
   \end{cases}
\end{equation}
Then the other root target differs from $(2+p)\bmod3$, as required
by Lemma~\ref{lem:path}.  Give each $\MC$ pendant its port from
Section~\ref{sec:path}; give each return path the default pair
$(0,2)$ until a selected return adjustment is made.

Consider the remaining case from Section~\ref{sec:mod}: a bipartite
core with targets 0 on $A$, $t$ on $B$ and two counted objects on
$A$.  Define
\begin{equation}\label{eq:defect}
 \delta_C=\sum_{v\in A}(\phi(v)-c(v))-
           \sum_{v\in B}(\phi(v)-c(v)).
\end{equation}
Its counted objects have the following interpretations.
\begin{enumerate}
\item A pendant port at a zero-target root may take either pure
value, at every length, by Lemma~\ref{lem:path}.
\item For an internal short connection whose endpoint $u\in A$
has target 0, the \emph{other endpoint's own port} is free to take
0 or 2.  The actual red degree at $u$ is 0 modulo three, hence
avoids both 2 and 4.  The comparison at the other endpoint is
protected by the already chosen port at $u$ in \eqref{eq:shortchoice}.
\item A zero-target root with no pendant and no short incidence may
change its target from 0 to $3-t$.  This remains different from
all neighbouring core targets $t$.  No short port or pendant choice
depends on this target.
\end{enumerate}

Each object changes \eqref{eq:defect} by a nonzero element of $\F$.
The variable freed by item~2 is at the opposite endpoint, so its
sign in the defect can be either sign; nonzeroness is all that is
needed.  Different short endpoints free different actual ports.
If both ends are on $A$, both actual ports are independently free.
An item~3 vertex has no short or pendant incidence, so it cannot
overlap variables of the first two kinds.  Thus two counted objects
really give independent binary choices.

For any nonzero $\alpha,\beta\in\F$,
\begin{equation}\label{eq:binary}
 \{\alpha\epsilon+\beta\epsilon':
          \epsilon,\epsilon'\in\{0,1\}\}=\F.
\end{equation}
Consequently the two chosen objects can cancel any initial defect.
The two-safe-exit case uses the same calculation.  The return case
uses its three total contributions instead.

All cores can be balanced separately.  Changing a target in item~3
never changes a short endpoint or pendant target.  Changing a port
uses only a pendant, a short connection internal to the same core,
or this core's end of a connection to an $\SC$ root.  It never
changes another $\MC$ core's offsets.  On a safe short exit the
opposite root accepts either pure port, and its own degree does not
change.  Long path interfaces accept arbitrary actual root degrees.
There is therefore no remaining coupled system between different
modular cores.

Apply Lemma~\ref{lem:linear} to each core with its balanced targets
and final offsets.  All resulting core weights belong to $\{0,1,2\}$.
The actual root degrees are the sums of these incident core weights
and the fixed ports, so they are physical and satisfy
$R(v)\equiv\phi(v)$.  Every direct core edge has distinct red
degrees at its ends and hence distinct blue degrees as well.

\section{Completion of the proof}\label{sec:complete}

\begin{proof}[Proof of Theorem~\ref{thm:main}]
The cases with no degree-three root and the case $K_4$ were treated
in Section~\ref{sec:path}.  Otherwise decompose the graph into cores
and paths and classify its cores as in Section~\ref{sec:class}.
Sections~\ref{sec:glue}--\ref{sec:controls} assign all core weights
and all root ports, with unequal degrees along every direct core
edge.  It remains to fill the path interiors.

For a short connection between two $\SC$ roots, both roots carrying
short ports are safe.  Any exceptional red degree 2 or 4 allowed by
Lemma~\ref{lem:two} occurs only at a root without short ports.
For an $\SC$--$\MC$ short connection, the safe root's comparison
is automatic and the modular root's comparison follows from its list.
For an $\MC$--$\MC$ short connection, \eqref{eq:shortchoice} and
the frozen short targets ensure both comparisons; Section~\ref{sec:controls}
preserves them.

Every ordinary three-edge connection either has opposite pure ports
or is a $D$ edge with one mixed port whose root degree is not 3.
Every longer ordinary connection has pure ports.  Return and pendant
paths have their admissible pairs or ports.  Lemma~\ref{lem:path}
therefore fills every path.

The interiors are edge-disjoint, and their first and last weights
were already counted in the root degrees.  Filling them does not
alter a root degree.  This verifies core edges, internal path edges,
root-end edges and leaf-end edges, exhausting $E(G)$.  Replacing a
weight $x$ by $x$ red and $2-x$ blue copies yields \eqref{eq:rule}
on every edge of the doubled graph.
\end{proof}

The argument uses the common identity $B=6-R$ at the roots, their
three-incidence budget and the resulting bound $\Delta(F)\le4$.
No theorem for graphs of larger maximum degree is asserted here.


\section{Interfaces with arbitrary endpoint degrees}\label{sec:generalports}

Here a \emph{root} means a vertex of original degree at least three,
without an upper bound. We use the same decomposition into direct root
edges, ordinary paths, return paths and pendant paths as before. A
non-direct ordinary path has two different actual end edges; their
weights may be chosen independently. A direct edge has one shared weight.
This distinction is essential throughout.

For a path $x_1=a,\ldots,x_\ell=b$ with internal vertices of degree two,
prescribed final endpoint counts $(r,s)$ give the allowed next weights
\begin{equation}\label{eq:generalP}
 P_{r,s}(a)=\{0,1,2\}\setminus
 \bigl(\{r-a:a>0\}\cup\{4-a-s:a<2\}\bigr).
\end{equation}
Values outside $\{0,1,2\}$ have no effect. Every internal comparison is
equivalent to $x_i\ne x_{i+2}$. Thus the odd and even positions form
independent proper three-colour paths with endpoint constraints.

\begin{lemma}\label{lem:generalpath}
If the two allowed sets in~\eqref{eq:generalP} each have at least two
elements, any end weights extend at length at least four; at length
three they extend if and only if $a\ne b$. Opposite pure end weights
extend at every length at least three for all physical endpoint counts.
If all four endpoint colour degrees are at least four, distinct end
weights extend already at length two.
\end{lemma}
\begin{proof}
At length three choose the middle value in the intersection of two
two-element subsets of a three-element set. The other condition is
$a\ne b$. At length four choose $x_2\in P_L\setminus\{b\}$ and
$x_3\in P_R\setminus\{a\}$. At length five choose
$x_3\notin\{a,b\}$ and different $x_2\in P_L,x_4\in P_R$.
For greater lengths each position path has at least two edges between
its endpoint constraints. A three-colour path of at least two edges
connects arbitrary prescribed colours, by induction starting with a
middle colour different from both ends. For a pure end weight, at most
one value is forbidden, proving the opposite-pure assertion. If
$r,s\ge4$, a mixed port forbids no value, and a pure port can only
forbid its own value, when the relevant colour degree is four. At
length two distinct ports therefore satisfy both endpoint checks.
\end{proof}

We will also use the following small extension of the first part: at
length at least three, different end weights extend if one allowed set
is $\{1\}$, the other contains 1 and has at least two elements, and the
opposite fixed end weight is pure. At lengths three and four this
follows from the same explicit choices; at length five the two allowed
even-position values can be made different; greater lengths follow
from the position-path induction.

\subsection{Exact preservation under subdivision}\label{sec:preserve}

\begin{theorem}\label{thm:preserve}
Let a legal colouring of $\doub Q$ be given. Replace a set of original
edges $e=uv$ by internally disjoint paths of lengths $L_e\ge4$. Set
\[
 \epsilon_e(v)=\mathbb{1}_{\{x_e=1,\ R_Q(v)=B_Q(v)\in\{1,3\}\}}.
\]
If $L_e\ge4+\epsilon_e(u)+\epsilon_e(v)$ for every replaced edge, the
colouring extends preserving every old vertex colour degree and every
unreplaced edge weight. At lengths four and five this condition is
also necessary when both new end weights must equal $x_e$.
\end{theorem}
\begin{proof}
Keep each end weight equal to $a=x_e$. Pure ports have at least two
allowed next values. For $a=1$, the forbidden next values are $R-1$
and $3-B$. If both belong to $\{0,1,2\}$, then $R,B\in[1,3]$.
Since $R+B$ is even, the allowed set has size one exactly when
$(R,B)=(1,1)$ or $(3,3)$, and that set is $\{1\}$. At length four
each next value must differ from the opposite fixed end value $a$;
one singleton is therefore fatal. At length five the two next values
must differ from one another, so precisely two singletons are fatal.
At length at least six each position path can connect arbitrary
nonempty boundary sets. These conditions give the stated threshold.
All old incident contributions remain identical, so all unreplaced
comparisons persist simultaneously.
\end{proof}

For cubic initial data there are 60 ordered legal end states: for each
$a$, each endpoint red degree lies in $\{a,\ldots,a+4\}$ and the two
degrees differ. All 40 pure-port states extend at length four. Of the
20 mixed-port states, exactly eight fail: one end red degree is three
and the other is one of $1,2,4,5$, in either order. These are fixed-data
obstructions, not counterexamples to Theorem~\ref{thm:main}.

In particular selected edges whose endpoints all have degree at least
four can always be replaced by paths of length at least four while
preserving a given colouring. Known regular-graph colourability
\cite{GOW26} makes this unconditional for every connected regular
mother graph of degree at least four. Known bipartite colourability
\cite{GW23} gives the same conclusion for a bipartite mother graph
when every selected endpoint has degree at least four. The initial
colouring is supplied by those theorems; it is not assumed for a
general unknown mother graph.

If a vertex-minimal counterexample to the general conjecture exists,
it has no ordinary degree-two thread of length at least five between
distinct nonadjacent roots; if both roots have degree at least four,
length at least four is excluded. Replace that thread by a single edge.
The smaller graph is simple and is not $K_2$. Its legal colouring
exists by minimality. On a mixed new edge at most one endpoint can have
$(R,B)=(3,3)$, since two such endpoints would already violate the edge.
Theorem~\ref{thm:preserve} restores the thread, a contradiction. The
nonadjacency assumption cannot be dropped in this argument.

\section{One-bridge trees and decorated paths}\label{sec:trees}

For a rooted tree $(T,r)$ and an external bridge of weight $c$, let
$P_c(T,r)$ be the achievable \emph{internal} root red sums when all
internal edges are legal after the root receives external contribution
$(c,2-c)$. The external bridge itself is not yet checked.

\begin{lemma}\label{lem:treeports}
For a one-vertex tree all three sets are $\{0\}$. For a two-vertex tree,
$P_0=\{0\},P_1=\{0,1,2\},P_2=\{2\}$. Every tree of order at least two
has $|P_1|\ge3$; every tree of order at least three has
$|P_0|,|P_2|\ge2$.
\end{lemma}
\begin{proof}
Use simultaneous induction on tree order. Delete the root, leaving
$k$ rooted subtrees with root internal degrees $d_i$. If the root edge
to subtree $i$ has weight $a_i$ and its internal root sum is $p$, the
only new restrictions, for final parent counts $A,B$, are
\[
 a_i>0\Rightarrow p+a_i\ne A,\qquad
 a_i<2\Rightarrow2d_i-p+2-a_i\ne B.
\]
Thus a mixed edge forbids at most two child candidates, a pure edge at
most one. Colour complementation gives
$P_{2-c}=\{2k-p:p\in P_c\}$.

For $c=0$, put all $a_i=1$. The internal sum $k$ is feasible: larger
children have three mixed candidates, and singleton children have
counts $(1,1)$, different from parent counts $(k,k+2)$ unless $k=1$,
which would be the two-vertex base case. Change one $a_i$ to 2 to seek
sum $k+1$ and parent counts $(k+1,k+1)$. That child must avoid
$p=k-1$ in $P_2$. A child of order at least three has two candidates;
a singleton fails only at the excluded $k=1$; a two-vertex child fails
only at $k=3$. The only remaining configuration is three two-vertex
children. Take root weights $(0,1,1)$ and all child internal weights
zero instead. The parent is $(2,6)$; the pure-blue child root has blue
degree four and the mixed child roots have counts $(1,3)$. This gives
a second sum distinct from $k=3$. Complementation proves the $P_2$ bound.

For $c=1$, all $a_i=1$ yield sum $k$ and counts $(k+1,k+1)$.
Changing one $a_i$ to zero seeks sum $k-1$ and counts $(k,k+2)$;
the child must avoid $p=2d_i-k$ in $P_0$. The only exception not
handled by the induction is two two-vertex children. In that case set
both root weights and both child internal weights to zero, giving
internal sum zero and parent counts $(1,5)$, with child blue degrees
four. Consequently a feasible sum $a\ne k$ always exists. Its colour
complement $2k-a$, together with $k$, gives three distinct sums.
\end{proof}

For an outside vertex of final counts $(A,B)$ the bridge imposes at
most one or two further exclusions. Lemma~\ref{lem:treeports} gives the
complete compatibility table
\begin{center}
\begin{tabular}{lccc}\toprule
Tree order&$c=0$&$c=1$&$c=2$\\\midrule
1&$B\ne2$&$A,B\ne1$&$A\ne2$\\
2&$B\ne4$&always&$A\ne4$\\
at least 3&always&always&always\\\bottomrule
\end{tabular}
\end{center}
Here ``always'' assumes physical counts containing the bridge copies.
This is an interface statement; adding the bridge changes outside
degrees and may invalidate other outside edges.

\subsection{Saturation of decorated paths}\label{sec:decorated}

\begin{theorem}\label{thm:decorated}
Let every internal vertex of a path have at most one bridge to a
disjoint rooted tree. A path of at least six edges accepts arbitrary
physical final endpoint counts and arbitrary prescribed first and last
edge weights, while colouring all attached trees. The endpoints may be
the same vertex if both incident contributions are counted.
\end{theorem}
\begin{proof}
Classify an internal vertex as $A$ (bare), $L$ (a singleton tree),
$S$ (a two-vertex tree), or $F$ (a larger tree). With incoming and
outgoing weights $a,p$, bridge weight $c$, and degree $D=2$ or 3, its
red sum is $r=a+p+c$. For $A$, put $c=0$. The compatibility table
above says: for $L$, $a+p=0$ forces $c=0$, $a+p=4$ forces $c=2$,
and otherwise any $c$ works; for $S$, $a+p=2$ forces $c=1$ and
otherwise any $c$ works; for $F$ every $c$ works.

Record a state $(p,D,r)$. There are 24 possible states,
\[
 \mathcal U=\{(p,2,r):p\le r\le p+2\}
 \cup\{(p,3,r):p\le r\le p+4\},\quad p\in\{0,1,2\}.
\]
An old state with outgoing weight $a$ can precede a new state
$(p,D',r')$ exactly when the indicated bridge is allowed and
$a>0\Rightarrow r\ne r'$, $a<2\Rightarrow2D-r\ne2D'-r'$.
Let $T_X$ be the resulting monotone map on state subsets. Complete
integer tables are given in Appendix~\ref{app:decorated}, including the
four lower-bound families $F_1,F_2,F_3,F_4$ of sizes $30,19,7,3$.

The initial endpoint has just 93 comparison types: each colour count
is $0,\ldots,6$ or greater than six, with physical and parity
restrictions. For each type and first letter the reachable set contains
an element of $F_1$. For every $M\in F_i$ and letter $X$, $T_X(M)$
contains an element of $F_{i+1}$, for $i=1,2,3$.
For each of the three elements of $F_4$, applying $T_X$ gives the
entire locally feasible state set $\Omega_X$. These are complete finite
identities, recomputed without the constructor by
\code{decorated_path_certificate_check.py}.

It follows by monotonicity and induction that after five internal
vertices the reachable set is $\Omega_X$, and remains saturated after
every further vertex. For a prescribed final weight, $\Omega_A$ has
three possible red sums, $\Omega_L,\Omega_F$ have five and $\Omega_S$
has at least four. The right boundary forbids at most two, so it can
always be met. Restore the attached trees using the bridge lemma.
The clipping of large endpoint counts is exact, since internal colour
degrees never exceed six. This proves arbitrary length and arbitrary
tree shapes from the finite identities.
\end{proof}

\section{Sparse direct branch cores}\label{sec:branches}

\begin{theorem}[A2]\label{thm:branches}
Let $W=\{v:d_G(v)\ge3\}$ for a connected simple graph $G\ne K_2$.
Each of the following hypotheses implies $\lir(\doub G)\le2$:
\begin{enumerate}
\item $G[W]$ is independent;
\item there is no degree-three vertex and $\Delta(G[W])\le1$;
\item $\Delta(G[W])\le1$ and every non-direct ordinary branch path
has length at least four.
\end{enumerate}
Return paths of length at least three and pendant paths of arbitrary
positive length are allowed in all three statements.
\end{theorem}

\subsection{Independent branch vertices}
\begin{proof}[Proof of the first statement]
For a root let $l,h,p$ count return paths, ordinary incidences and
pendant paths, respectively; its original degree is $D=2l+h+p$.
On $W$ form a simple graph joining distinct roots connected by an
ordinary path of length at least three. Choose a maximal independent
set $A$, then a maximal independent set $B$ in the remaining graph,
and put $C=W\setminus(A\cup B)$. Every $C$ root has a long neighbour
in both $A$ and $B$; every isolated root belongs to $A$.

Give every return path ports $(0,2)$. Give pendant root ports
0, 2, 1 on $A,B,C$, respectively. On ordinary paths assign port pairs
$AA:(0,0)$, $BB:(2,2)$, $AB:(0,2)$, $AC:(0,1)$, $BC:(2,1)$.
An $AA$ or $BB$ path necessarily has length two. Balance-orient the
multigraph of $CC$ paths and give each tail port 0 and head port 2.
For a length-two $AC$ path whose $C$ root has degree three, change
its $C$ port to 2; for such a $BC$ path change it to zero.

At a $C$ root of degree at least four the only imbalance between
its red sum and $D$ is at most one from the balanced $CC$ orientation.
At degree three the two long $A,B$ incidences leave at most one other
incidence, so either one $CC$ imbalance or one short-path correction
occurs, again of magnitude at most one. If the red sum is $D-1$,
change its long $A$ port from 1 to 2. If it is $D+1$, change its
long $B$ port from 1 to zero. These independent end ports affect no
other root. Now every $C$ root has $(R,B)=(D,D)$.

At an $A$ root the counts are $(2l,2(D-l))$, and at a $B$ root they
are reversed. The larger count is an even integer at least four.
Equality to four at a root with an external path occurs only when
$D=3,l=1,h+p=1$. Such a root is exceptional only in the following
two cases. If its unique external path has length two and joins
another $A$ root, change its own port from 0 to 2; if both ends
are exceptional change just one. Its counts become $(4,2)$ and the
midpoint has $(2,2)$, legal on the red and blue sides. If its unique
external path is a two-edge pendant, change that port to 1 and use
the word $(1,1)$; root, middle and leaf counts are $(3,3),(2,2),(1,1)$.
These exceptional roots have no long ordinary incidence, so none of
the previous adjustments is disturbed.

The remaining two-edge ordinary checks are listed below. The displayed
root counts refer to the colour present on the corresponding port.
\begin{center}\small
\begin{tabular}{llll}\toprule
Type&Ports&Middle $(R,B)$&Root comparisons\\\midrule
$AA$&$0,0$&$0,4$&both blue at least 6\\
$BB$&$2,2$&$4,0$&both red at least 6\\
$AB$&$0,2$&$2,2$&relevant counts at least 4\\
$AC$, $D_C=3$&$0,2$&$2,2$&blue at least 4; red 3\\
$AC$, $D_C\ge4$&$0,1$&$1,3$&blue at least 4; both at least 4\\
$BC$&complement&complement&same inequalities\\
$CC$&$0,2$ or $2,0$&$2,2$&both root counts at least 3\\\bottomrule
\end{tabular}
\end{center}
For long ordinary paths the end ports differ. An $A$ or $B$ allowed
set has at least two elements and contains 1. At a $C$ root the only
singleton occurs at $D=3$ with port 1, when it is $\{1\}$; its
other end is an $A$ or $B$ pure port. The general path lemma and its
singleton extension therefore apply. Every return path has opposite
pure ports. Pendant pure ports use the constant word at lengths one
or two, with the unique bad two-edge case already repaired; at longer
lengths use the opposite pure leaf port. A $C$ pendant uses $(1)$ or
$(1,1)$ at lengths one or two; otherwise give its leaf port zero and
apply the singleton extension. All root counts were fixed first.
Consequently the independent path fillings are simultaneously legal.
\end{proof}

\subsection{A direct matching with no degree-three roots}
\begin{proof}[Proof of the second statement]
Compress all paths, using a private placeholder leaf for each pendant
path only for the purpose of orientation. A balanced orientation of
this multigraph gives outdegrees $o(v)\in\{\lfloor D/2\rfloor,
\lceil D/2\rceil\}$. Give non-direct tail ports 1 and head ports 0;
give a pendant root port 1 or 0 according to its direction; give each
direct edge the shared weight 1. If $i_H(v)$ indicates an incoming
direct edge, the initial counts are
\[
 R_0(v)=o(v)+i_H(v)\ge2,\qquad
 B_0(v)=2D-R_0(v)\ge5.
\]
Repair each direct matching edge independently. If only its red
degrees coincide, change its weight to zero. The head initially had
red degree at least three, so both new red degrees remain at least
two. If only its blue degrees coincide, change its weight to 2.
The tail initially had blue degree at least six, so both new blue
degrees remain at least five. If neither coincides, retain weight 1.

If both coincide, orient the edge $u\to v$. Equal original degrees
and $o(u)=o(v)+1$ force $D=2k+1$, $k\ge2$, with common initial
counts $(k+1,3k+1)$. There are $k$ other outgoing incidences at $u$.
If one is an ordinary or return path, change its tail port 1 to 2,
giving $u$ counts $(k+2,3k)$. Otherwise choose an outgoing pendant
port and change 1 to zero, giving $(k,3k+2)$. The mixed direct edge
is now legal. Only tail ports were modified; head ports stay zero,
so even repairs selecting the same connection cannot interfere.
Different direct edges have disjoint endpoint sets.

Final root counts satisfy $R\ge2,B\ge5$, and a non-direct port 2
occurs only at $R\ge4$. Every non-direct port pair is $(0,1)$ or
$(0,2)$ up to reversal. Its two allowed sets have size at least two,
so every length at least three extends. At length two the midpoint
is respectively $(1,3)$ or $(2,2)$ and the same inequalities verify
both ends. This includes return paths at their legal lengths.
Pendant port zero uses $(0)$, $(0,0)$, or opposite leaf port 2;
pendant port one uses $(1)$, $(1,0)$, or leaf port zero, according
as its length is one, two, or at least three. Direct checking or
Lemma~\ref{lem:generalpath} verifies every case.
\end{proof}

\subsection{A direct matching with ordinary paths of length four}
\begin{proof}[Proof of the third statement]
Let $h\in\{0,1\}$ indicate a direct incidence, let $l$ count return
paths, and let $z$ count all other root incidences. Thus $D=h+2l+z$.
Give direct edges weight 1, return ports $(0,2)$, and choose the
$z$ free ports from $\{0,2\}$. For a root with $h=1$, if $a$ free
ports are 2 its counts are
\[
 (R,B)=(1+2l+2a,\ 1+2l+2(z-a)).
\]
For a direct edge, if both ends have free ports and one has $z\ge2$,
fix the other end and try $a=0,1,2$. At most two candidates are
forbidden. If both have $z=1$, their candidates are
$(D-1,D+1)$ and $(D+1,D-1)$: use the same choice at different degrees
and opposite choices at equal degrees. If one end has $z=0$, it has
a return path. Change one such port pair from $(0,2)$ to $(0,1)$.
That end now has even counts $(D-1,D+1)$ and the other keeps odd
counts. If both have $z=0$, change just one end. Thus every direct
edge is legal, independently of all others.

At roots with $h=0$, leave free ports zero; counts are
$(2l,2(l+z))$. The sole exceptional pendant is $l=z=1$ with length
two. Change its port to 1, making the root $(3,3)$, and use $(1,1)$.
Every ordinary non-direct path has pure ports and length at least
four, so the general path lemma applies even when its ports are equal.
Unmodified returns have ports $(0,2)$. A modified return has ports
$(0,1)$ at counts $(D-1,D+1)$; both allowed sets have at least two
elements, so it also extends at length three. At a directly matched
root, each colour used by a pendant pure port has odd degree at least
three, different from the short pendant values two and four. At an
unmatched root its blue degree is at least four, and equality to four
on a two-edge pendant was exactly the exception repaired above.
Long pendants use opposite pure leaf ports. These cases exhaust the
construction.
\end{proof}

In all three proofs the case $W=\varnothing$ is separate: a single
vertex has no edges; a two-edge path has word $(0,0)$; a longer path
uses opposite pure leaf ports. For a cycle, properly three-colour the
position graph joining positions $i$ and $i+2$ cyclically. Its maximum
degree is at most two and its resulting word satisfies every internal
comparison. The excluded one-edge path is $K_2$.

\section{Homogeneous direct cores with long connections}\label{sec:homogeneous}

\begin{lemma}[Even signed forms]\label{lem:evenforms}
Suppose each of finitely many homogeneous linear forms has an even
number, at least four, of distinct variables, all coefficients are
$\pm1$, and each variable occurs in at most two forms. There is a
choice of signs $\pm1$ for the variables making every form nonzero.
\end{lemma}
\begin{proof}
Make a signed multigraph whose vertices are forms and whose shared
variables are edges; variables occurring once are boundary variables.
In a component with a boundary variable, retain a spanning tree and
one root boundary variable, set all other variables to 1, and process
leaves upwards. The one remaining sign at each step avoids at most
one forbidden value.

In a component without boundary variables every vertex has even
degree at least four. Give each edge the product of its endpoint
coefficients as sign. There is a positive cycle: if no fundamental
cycle is positive, two negative fundamental cycles sharing a tree
edge have a positive cycle as their symmetric difference; if none
share tree edges the graph is a cactus and has a vertex of degree at
most two, a contradiction. Keep such a cycle and a forest rooted on
it. Eliminate off-cycle vertices as before. On the cycle the relations
are $a_i y_i+b_i y_{i-1}+C_i\ne0$ with even $C_i$.
For $C_i=0$ the allowed relation is a bijection on two signs; for
$C_i=\pm2$ it has three of the four pairs; for $|C_i|\ge4$ all four
pairs. If every relation is a bijection, positivity makes their
composition have a fixed point. Otherwise their composition has at
least three pairs: bijections preserve size, and composition of
relations with at least three pairs has at least three pairs. Such
a relation contains a diagonal pair. Hence the signs close around
the cycle.
\end{proof}

\begin{theorem}[A3]\label{thm:homogeneous}
Let $G$ have minimum degree at least two. In each component of
$H=G[\{v:d_G(v)\ge3\}]$, suppose the original degree is constant and
the component is three-colourable. Suppose all ordinary non-direct
connections have length at least three and there are no return paths.
Then $\lir(\doub G)\le2$.
\end{theorem}
\begin{proof}
Orient each long connection and give it opposite pure ports. Its
contribution $c(v)$ is twice an indegree. For an isolated direct-edge
component $uv$ of original degree $D$, the two long-incidence counts
are both $D-1$. Reversals of the long paths are sign variables, and
$c(u)-c(v)$ is an even form with $2(D-1)\ge4$ incidences. If a long
path joins $u$ and $v$ its sign occurs twice with combined coefficient
$\pm2$; it is private to this equation. Temporarily omit equations
with such variables. All remaining forms satisfy
Lemma~\ref{lem:evenforms}; then restore each omitted equation by its
private variable, one of whose two signs makes it nonzero. Give each
isolated direct edge weight 1. It is now legal in both colours.

For every larger direct component choose proper targets $\phi$ over
$\mathbb F_3$ and solve the internal incidence equations
$\sum_{e\ni v}x_e=\phi(v)-c(v)$. A nonbipartite component realises
arbitrary demands: a spanning tree and an edge joining two vertices
of the same tree parity allow adjustment of the signed total by an
invertible coefficient $\pm2$, followed by leaf elimination. In a
bipartite component solvability is precisely the alternating-sum
condition in Lemma~\ref{lem:linear}. Choose the larger side, of size
at least two, initially all zero and the other side all one. Changing
zero, one or two vertices of the first side to 2 changes that sum by
$0,2,4$, covering all residues. The targets stay proper. Representatives
$0,1,2$ of the solved edge weights give different red sums on every
direct edge. The original degrees agree there, so the blue sums also
differ. Isolated roots require no direct comparison. Finally fill
all long paths with their opposite pure ports.
\end{proof}

For an $r$-regular simple mother graph, replacing selected edges by
paths of length at least three is therefore allowed whenever the
retained graph is three-colourable. At $r=3$, a nonempty selected set
leaves every retained component 2-degenerate, since a cubic retained
component would be a whole connected component of the original graph.
This recovers the earlier cubic subdivision constructions, already
contained in A1. At unbounded degree the three-colourability hypothesis
is substantive and has not been removed.

\section{Balanced half-edges and simultaneous extensions}\label{sec:balanced}

\begin{lemma}\label{lem:halfedge}
Let $J$ be a multigraph, with loops allowed and no vertex of degree
one. Its incidences admit values $0,1,2$ summing to $d_J(v)$ at every
vertex, different at the two ends of every edge. Odd-degree vertices
have exactly one incidence of value 1 and even-degree vertices none.
For disjoint pairs of positive-degree vertices of the same parity,
one may instead require exactly one vertex of each pair to have sum
$d_J(v)\pm1$, leaving every other sum unchanged.
\end{lemma}
\begin{proof}
Balance-orient $J$, assigning 0 at tails and 2 at heads. At an odd
vertex choose one majority-valued incidence to change to 1. There
are $\lceil d/2\rceil\ge2$ different candidate edge identifiers;
a loop cannot supply two majority-valued candidates. The request-edge
bipartite graph has each request adjacent to at least two edges, each
edge to at most two requests. Hall's condition follows by counting,
so all changes can use distinct edges. The sums become exact and
end values stay different.

For an odd-degree pair skip the correction at one member. For an
even-degree pair create one request selecting any incidence at either
member and changing it to 1; the union contains at least two edge
identifiers. Each original vertex belongs to at most one request,
so the same Hall argument applies. These changes give precisely the
claimed one-unit deviations.
\end{proof}

\begin{theorem}[A4]\label{thm:partialbalanced}
Let a connected multigraph $M$ have original degrees $D(v)\ge4$.
Subdivide an edge set $J$ into paths of length at least two, with all
loops in $J$ and loop paths of length at least three. Suppose $J$ has
no vertex of degree one, the retained graph $H=M-J$ is simple, and
every retained edge has endpoints of different original degrees.
Then the simple expanded graph has a legal doubled two-colouring
with $R(v)=B(v)=D(v)$ at every old vertex.
\end{theorem}
\begin{proof}
Use Lemma~\ref{lem:halfedge} on $J$ and give every retained edge
weight 1. Each old vertex receives $d_J+d_H=D$ in both colours.
Retained edges compare different $D$ values. Each subdivided path
has distinct end weights and all four endpoint colour degrees at
least four, so Lemma~\ref{lem:generalpath} fills it at every allowed
length without changing the old sums.
\end{proof}

\begin{theorem}[A5]\label{thm:simultaneous}
Let $H$ be a simple graph whose components of order at least three
have given legal doubled two-colourings. On the same old vertex set
let $J$ be a connection multigraph with $k(v)=d_J(v)$ either zero
or at least four, constant on each component of $H$. Require $k\ge4$
on every $K_2$ component. Replace $J$ edges by internally disjoint
paths of length at least two, or at least three for loops. The
colouring extends to the simple expanded graph preserving every
given old edge weight. No initial colouring of $K_2$ is assumed.
\end{theorem}
\begin{proof}
Pair the endpoints of the $K_2$ components in
Lemma~\ref{lem:halfedge}. All other old vertices receive exactly
$(k,k)$ from new ports. Every given old colour-degree difference is
unchanged, since $k$ is equal at its two ends. Give each old $K_2$
edge weight 1. Its endpoints then have counts
$(k+1+\epsilon,k+1-\epsilon)$ and $(k+1,k+1)$ for
$\epsilon\in\{-1,1\}$, so that edge is legal. All vertices incident
with new paths have both colour degrees at least four. Different
new end ports can therefore be filled at length two by
Lemma~\ref{lem:generalpath}. All old and new checks hold together.
\end{proof}

A vertex-minimal counterexample cannot admit the latter decomposition
with $J\ne\varnothing$: deleting its new internal vertices gives
smaller old components, while the $K_2$ and isolated cases are already
handled. This does not assert that every graph has such a decomposition.

\section{Near-centre parity assignments}\label{sec:parity}

\begin{lemma}\label{lem:nearparity}
A parity prescription $p:V(G)\to\{0,1\}$ can be realised by weights
$0,1,2$ with $R(v)\equiv B(v)\equiv p(v)\pmod2$ and
$|R(v)-d_G(v)|\le1$ if and only if $\sum_{v\in C}p(v)$ is even
in every component $C$. The edges changed from weight 1 may be
restricted to a forest.
\end{lemma}
\begin{proof}
Necessity is the red handshake identity. Start with all weights 1.
In each spanning tree choose a $T$-join for
$T=\{v:p(v)\ne d_G(v)\pmod2\}$ by leaf elimination; $|T|$ is even.
Pair incidences of the chosen forest at each vertex, leaving one
unpaired at each odd-degree vertex. Its edges decompose into open
trails with those unpaired ends; there is no closed trail in a forest.
Along each trail change weights alternately to 2 and zero. Contributions
cancel at every paired passage, leaving a change of $\pm1$ exactly
at the vertices of $T$. All claims follow, including the exact centred
assignment when $T$ is empty.
\end{proof}

\begin{theorem}[A6]\label{thm:gapparity}
Let $F$ consist of the original edges $uv$ with
$|d_G(u)-d_G(v)|\le2$. If $F$ is bipartite and its bipartition
phases can be chosen to give an even total number of odd labels,
then $\doub G$ has a legal two-colouring. Both colours distinguish
the ends of every original edge.
\end{theorem}
\begin{proof}
Realise those labels by Lemma~\ref{lem:nearparity}. On $F$, both
colour degrees have different parities. Off $F$, the intervals
$[d_G(v)-1,d_G(v)+1]$ at the two ends are disjoint, and contain
both colour degrees. This verifies every comparison.
\end{proof}

The phase condition is exact: an odd-order component of $F$ allows
the total parity to be changed by a flip. If all components have even
order, it holds exactly when an even number of components have two
odd-sized sides. For example, connect the roots of three disjoint
$K_{r_i,r_i}$ modules in a triangle, where the $r_i\ge4$ are even
and pairwise differ by at least four. The graph has no degree-two
threads and is nonbipartite. Its $F$ consists precisely of the three
modules with even sides, so the theorem applies. This condition does
not cover general nonbipartite regular cores.

\section{Clique interfaces, block trees and capacity}\label{sec:cliques}

Fix a root $z$ of $K_n$, external contributions $A,B\ge0$ with
$A+B$ even, and a desired internal root red sum $S$. Only internal
edges are to be checked, using final root counts
$(A+S,B+2(n-1)-S)$.

\begin{theorem}[Complete clique root interface]\label{thm:cliqueport}
For every $n\ge4$ and $0\le S\le2(n-1)$ there is a legal internal
colouring with the prescribed root sum. For $n\ge5$ every nonroot
vertex may additionally have both colour degrees in $[1,2n-3]$.
For $n=4$ the same strengthening holds except at
$(A,B,S)=(1,3,4),(3,1,2)$; the ordinary interface still holds there.
\end{theorem}
\begin{proof}
For $n=4,5$, put $r=n-1$. External counts exceeding $2r$ are
comparison-equivalent to $2r+1$ for odd counts and $2r+2$ for even
counts, since every nonroot count is at most $2r$. Consequently the
entire base domain is
\[
 A,B\in\{0,\ldots,2r+2\},\quad A+B\equiv0\pmod2,
 \quad S\in\{0,\ldots,2r\}.
\]
It contains $287+549=836$ entries. Actual edge-weight witnesses are
in \code{rooted_clique_interfaces.json}. The independent standard-library
checker \code{rooted_clique_certificate.py} verifies equality with this
whole domain, root sums, every actual colour comparison, and every
nonroot bound outside the two exceptions. Thus the bases include
arbitrarily large physical external counts by exact clipping, not by
an unproved extrapolation.

For $n\ge6$ use strong induction by two vertices. If
$2\le S\le2n-4$, take a rooted $K_{n-2}$ with the same external
counts and internal sum $S-2$. Add $u,v$, set $wu=0,wv=2$ for
every old $w$, and $uv=1$. All old colour degrees increase by two;
the new counts are
\[
 u:(1,2n-3),\qquad v:(2n-3,1).
\]
Old nonroot counts are within $[2,2n-4]$, so all their new edges
are legal. Write the final root counts as $p,q$. If neither equals
$2n-3$, its two new pure edges are also legal. Otherwise both
$p,q$ are odd because $p+q$ is even. Choose an old nonroot $w_0$
and change its two new weights from $(0,2)$ to $(1,1)$. All old
counts stay fixed, while the new counts become
\[
 u:(2,2n-4),\qquad v:(2n-4,2).
\]
Inductively strengthened old nonroot counts, after adding two, lie
in $[3,2n-5]$ and are strictly separated from these boundary values.
The root counts are odd, so its comparisons with the new even values
hold. When $n=6$ and the old $K_4$ is one of the two exceptional
bases, the final root counts are instead $p=q=7$, different from
$2n-3=9$; the unmodified construction already works. Thus no missing
strengthened base is needed.

For $S=0,1$, take the old $K_{n-2}$ with the same internal sum but
formal external counts $A'=A-2,B'=B+2$. If $A\ge2$, these are
physical. If $A=0,1$, obtain the old witness using the nonnegative
replacement $\widehat A=A\bmod2$ and then reinterpret its root red
count with $A'$. This keeps the old internal edges legal: for $S=0$
the root has no red edge; for $S=1$ its sole red incidence meets a
vertex of red degree at least one, whereas the reinterpreted root
red count is nonpositive. Its blue count is unchanged. These formal
negative offsets are only an intermediate comparison argument.

Add $u,v$ with $wu=0,wv=2$ for old nonroots, $zu=zv=0$, and
$uv=2$. Relative to the reinterpreted old counts every old vertex
gains two in each colour, so the root now has exactly the originally
specified physical counts and internal sum $S$. The new counts are
$(2,2n-4),(2n-4,2)$. The old nonroot counts lie in $[3,2n-5]$;
the root blue count is at least $2n-3$. All new comparisons follow.
The old base parameter $S=0$ or 1 is never an exceptional $K_4$
parameter. Finally obtain $S=2n-3,2n-2$ by exchanging the colours
and $A,B$. These cases cover every $S$ and preserve the asserted
nonroot bounds.
\end{proof}

\subsection{Bouquets and arbitrary cores}\label{sec:cliquecapacity}

\begin{theorem}[A8]\label{thm:bouquet}
At a vertex of old degree $h$ in a given legal colouring, attach
cliques of orders $n_i\ge4$ sharing only that vertex, with total
new degree $L=\sum_i(n_i-1)\ge h$. Every old edge colour can be
preserved. Alternatively, for an arbitrary uncoloured simple core
$C$, attach such a bouquet at every vertex $v$ with
$L_v\ge d_C(v)$. The entire expanded graph can then be coloured.
\end{theorem}
\begin{proof}
For the first assertion, let the old counts at the root be $a,b$.
An added red contribution $s\in[0,2L]$ gives final counts
$(a+s,b+2L-s)$. Each old pure edge forbids at most one value of
$s$ and each mixed edge at most two. There are at most $2h$
forbidden values and $2L+1>2h$ candidates. Choose an allowed $s$
and split it into $s_i\in[0,2(n_i-1)]$. Clique $i$ receives external
counts
\[
 A_i=a+s-s_i,\quad
 B_i=b+2L-s-\bigl(2(n_i-1)-s_i\bigr).
\]
These are nonnegative with even sum. Theorem~\ref{thm:cliqueport}
colours it with internal sum $s_i$. Every clique uses the same
final root counts and all old neighbours keep their counts.

For the second assertion give every core edge weight 1. Put
$h_v=d_C(v)$ and $D_v=h_v+L_v$. In any vertex order choose
$R(v)\in[h_v,h_v+2L_v]$. A previously processed neighbour forbids
at most $R(w)$ and $2D_v-B(w)$, so at most $2h_v$ values are
forbidden. There are at least $2h_v+1$ candidates. Choose all red
targets greedily and set $B(v)=2D_v-R(v)$. Every mixed core edge
now distinguishes both colours. Split the remaining red contribution
among the cliques and realise it by the root interface theorem.
\end{proof}

The first capacity bound cannot uniformly be reduced to $L=h-1$
for $h\ge4$. Pair the integers $0,\ldots,2h-2$ within their even
and odd subsequences using adjacent values, covering a final odd
triple by two overlapping pairs. This gives exactly $h$ pairs
$(a_i,b_i)$ with $a_i=b_i+2$, distinct $a_i$, distinct $b_i$, and
union $\{0,\ldots,2h-2\}$. On a central $K_{h+1}$ with vertices
$v,u_1,\ldots,u_h$ put every weight 1. Attach a rooted $K_{h+1}$
at each $u_i$, with internal red sum $a_i$ using the theorem.
The old counts are
\[
 v:(h,h),\qquad u_i:(h+a_i,\ 3h-2-b_i).
\]
All central mixed edges are legal: both lists at the $u_i$ are
distinct and at least $h+2$. Add $K_h$ at $v$. Its contribution
$s\in[0,2h-2]$ is forbidden by $vu_i$ precisely at $s=a_i$ or
$s=b_i$. The pairs cover every candidate. Nonetheless the entire
graph is colourable by the block-tree theorem below, rooted at $v$:
each nonroot cutvertex has equal parent and child capacities. Thus
this is a sharp preserved-colour obstruction, not a graph obstruction.

\subsection{Clique block trees with full root freedom}

\begin{theorem}[A7]\label{thm:cliquetree}
Let every block of a connected graph be a clique of order at least
four. Root its block-vertex incidence tree at any vertex $z$.
At every nonroot cutvertex require the sum of its child-block degrees
$\sum(n_i-1)$ to be at least its parent-block degree $n_0-1$.
For all physical external counts $A,B$ at $z$ and all
$S\in[0,2d_G(z)]$, the graph has a legal internal colouring with
root internal red sum $S$. In particular all equal-order clique
block trees satisfy this conclusion.
\end{theorem}
\begin{proof}
First colour the bouquet of root blocks, distributing the specified
sum $S$ between them and using the root interface with the contributions
of the other root blocks included as external counts. Process the
other cutvertices outwards. Before its children are added a cutvertex
belongs to its one parent block and has current degree $h=n_0-1$.
The child bouquet has capacity at least $h$, so
Theorem~\ref{thm:bouquet} extends the current colouring without changing
any old edge weight. Actual external root counts remain included in
all old comparisons. Adding blocks away from $z$ does not change
its internal red sum. The finite incidence tree is eventually exhausted.
\end{proof}

In a vertex-minimal counterexample, any bouquet of leaf clique blocks
of order at least four at a vertex must have capacity $L<h$, where
$h$ is its degree in the remainder. Otherwise delete the leaf blocks,
colour the smaller remainder, and apply Theorem~\ref{thm:bouquet}.
If the remainder is $K_1$, use the root interface directly. If it
is $K_2$, give its edge weight 1 and the bouquet total red contribution
1; the endpoints have counts $(2,2L)$ and $(1,1)$, so the old edge
is legal and all cliques can be filled. Hence no hidden $K_2$
assumption is used in this minimality argument.

\section{Gluing coloured pieces along a tree}\label{sec:pieceglue}

Let connected pieces $P_i$ partition the edges of $G$. Let $U$ be
the vertices belonging to more than one piece; all other vertices
are called private. Suppose the bipartite incidence graph of pieces
and shared vertices is a tree. In the given local weighting of $P_i$,
write $d_i,r_i,b_i$ for the local degrees and set
\[
 a_i(v)=r_i(v)-d_i(v),\quad
 M_v=\max_{i:v\in P_i}|a_i(v)|,\quad
 I_v=[d_G(v)-M_v,d_G(v)+M_v]\quad(v\in U).
\]

\begin{theorem}[A9]\label{thm:pieceglue}
One may obtain a legal colouring solely by whole-piece colour flips
provided that: private--private edges are locally legal; each edge
from shared $v$ to private $u$ in $P_i$ satisfies
$d_G(v)-M_v>\max(r_i(u),b_i(u))$; and each shared--shared edge $uv$
satisfies $|d_G(u)-d_G(v)|>M_u+M_v$.
\end{theorem}
\begin{proof}
Root the incidence tree at a piece. On reaching a shared vertex its
parent piece's flip is already fixed. Starting with its contribution
$s=\pm a_i(v)$, choose the flip of each child piece to minimise the
new absolute sum. The identity
\[
 \min\{|s+a|,|s-a|\}\le\max\{|s|,|a|\}
\]
keeps $|s|\le M_v$. Each child piece has only one parent shared
vertex, so no choice is made twice. At the end the global red and
blue degrees at $v$ are $d_G(v)\pm s$, both in $I_v$.
Private counts are their local counts up to a swap. The three edge
types in the hypotheses are respectively preserved by local legality,
strict separation above private counts, or disjoint shared intervals.
Every edge is therefore legal.
\end{proof}

If the given pieces have maximum degree at most $q$, a sufficient
uniform condition is $d_G(v)\ge3q+1$ for shared vertices and degree
gap at least $2q+1$ on shared--shared edges. Indeed $M_v\le q$
and private colour degrees are at most $2q$. Taking known colourings
of cycle or cubic blocks gives, for example, the conditions shared
degree at least 10 and adjacent shared-degree gap at least 7. This
allows unbounded global maximum degree and pieces that are not cliques.

\subsection{Why whole-piece flips do not suffice in general}\label{sec:piecebarrier}

For each $r\ge2$, identify the distinguished left vertex $z$ of
two copies of $K_{r,r}$. In each copy put weight 2 in column 1,
weight zero in column 2, weight 1 in all remaining entries of row
$z$, and zero elsewhere. The local red sums are $r$ at $z$, 2 at
other rows, $2r$ at column 1, zero at column 2, and 1 at other
columns. As the local original degrees are equal, these values
verify local legality on every edge.

The local root counts are $(r,r)$ even after a whole-piece flip.
After identification the root counts are $(2r,2r)$, and in each
flipped or unflipped copy it still has a same-colour neighbour of
that colour degree $2r$. Thus all four choices of whole-piece flips
fail, although the incidence graph is a tree. For $r\ge3$, changing
one entry of row $z$ in a remaining column from 1 to zero in one
copy repairs the whole graph: the root becomes $(2r-1,2r+1)$,
the changed column's red sum becomes zero, and all other local
counts persist. At $r=2$ replace one copy's matrix by rows
$(0,0),(2,2)$. These direct repairs show why the obstruction concerns
given-piece gluing and not the conjecture.

\section{Fixed targets, factors and exchanges}\label{sec:targets}

This section gives applications of classical matching and factor
theory. These applications distinguish three tasks: independently
selecting mutually distinguishing targets, realising one target vector
by shared edge weights, and finding a target for which the actual
colour constraints hold. Only the last task for every graph would
prove the general conjecture.

\subsection{Independent target lists are not edge realisations}

\begin{lemma}\label{lem:complementlists}
Give each vertex $v$ of a simple graph $H$ a real number $D_v$ and a
list $L_v$ of at least $d_H(v)+1$ distinct real numbers. One can
choose $z_v\in L_v$ such that both $z_v$ and $D_v-z_v$ distinguish
the endpoints of every edge.
\end{lemma}
\begin{proof}
Orient the edges arbitrarily and write
$Q(z)=\prod_{uv}(z_u-z_v)=\sum_\alpha q_\alpha z^\alpha$.
For $d=(d_H(v))_v$ and $m=|E(H)|$,
$z^d Q(z^{-1})=(-1)^m Q(z)$, so
$q_{d-\alpha}=(-1)^m q_\alpha$. The coefficient of $z^d$ in
the highest homogeneous part of $Q(z)Q(D-z)$ is therefore
$\sum_\alpha q_\alpha^2>0$. Polynomial interpolation on the lists,
equivalently the Combinatorial Nullstellensatz~\cite{Alon99}, gives
a nonzero evaluation.
\end{proof}

This lemma says nothing about incidence realisation. On a triangle,
the lists $\{0,2,4\}$ with $D_v=4$ allow mutually distinct
complementary targets, but the targets $0,2,4$ cannot be realised
by weights $0,1,2$: the zero vertex forces its two weights zero,
preventing its degree-four neighbour from receiving four.

\subsection{A precise fixed-target criterion}

Consider a simple subgraph $H$ whose vertices receive fixed external
counts $a(v),b(v)$. Specify total red targets $R(v)$ and put
\[
 q(v)=R(v)-a(v),\qquad B(v)=b(v)+2d_H(v)-q(v).
\]
Define $E_0=\{uv:R(u)=R(v)\}$ and
$E_2=\{uv:B(u)=B(v)\}$. If the two sets overlap, no legal realisation
of these targets exists. Otherwise set
\[
 K=H-(E_0\cup E_2),\qquad t(v)=q(v)-2d_{E_2}(v).
\]

\begin{proposition}\label{prop:fixedfactor}
The targets admit a legal internal colouring exactly when
$q(v)\in[0,2d_H(v)]\cap\mathbb Z$, $E_0\cap E_2=\varnothing$,
and $\doub K$ has a spanning subgraph of degrees $t(v)$.
\end{proposition}
\begin{proof}
On $E_0$ all weights are forced to zero, on $E_2$ to 2. Every
remaining edge distinguishes both total colour degrees and is free
to have any weight. Its selected red copies must supply exactly $t$.
This correspondence is reversible.
\end{proof}

Here is an explicit ordinary matching reduction. For each copy of
$\doub K$ create its two endpoint ports and join them. At $v$ create
$2d_K(v)-t(v)$ auxiliary vertices adjacent to all ports at $v$.
A perfect matching uses that many local ports, leaving exactly
$t(v)$ ports matched across original copies. Select those copies
for the factor. Conversely any such factor completes to a perfect
matching by matching unused ports locally. The auxiliary graph has
at most $8|E(K)|$ vertices. Range failures for $t$ are rejected first.

\subsection{Capacity and odd-component certificates}

For disjoint $S,T\subseteq V(K)$ put
\[
 g_t(S,T)=t(S)+2\sum_{v\in T}d_{K-S}(v)-t(T),
\]
and let $o_t(S,T)$ be the number of components $C$ of $K-S-T$ for
which $t(C)$ is odd.

\begin{proposition}\label{prop:capacity}
Real weights $0\le x_e\le2$ realising $t$ exist if and only if
$g_t(S,T)\ge0$ for all disjoint $S,T$. Integer weights exist if and
only if $g_t(S,T)\ge o_t(S,T)$ for all such pairs.
\end{proposition}
\begin{proof}
Necessity of capacity follows by subtracting the sums over $T$ and
$S$ and bounding the remaining terms by
$4|E_K(T)|+2|E_K(T,V\setminus(S\cup T))|$.
For sufficiency, let $\mathcal Z=\{Bx:x\in[0,2]^{E(K)}\}$, where
$B$ is the unsigned incidence matrix. Its support function at $y$ is
$2\sum_{uv}\max(0,y_u+y_v)$. If $t\notin\mathcal Z$, a separating
vector $y$ violates this bound. Write $y$ as the integral, over
$h\ge0$, of the vector taking value 1 where $y_v>h$, value $-1$
where $y_v<-h$, and zero elsewhere. For each edge the same integral
identity holds for the positive part of $y_u+y_v$; it is checked
separately when the two signs agree or disagree. Some three-valued
vector must therefore violate the support bound, which is precisely
a violated $g_t(S,T)\ge0$.

For integers apply Tutte's factor theorem~\cite{Tutte54} to $\doub K$.
In its odd-component term the number of edges from a component to
$T$ is even because every original edge is doubled. The odd term
is consequently exactly $o_t(S,T)$, and its degree term is $g_t$.
This proves the stated equivalence.
\end{proof}

For even total target sum, $g_t-o_t$ is even. Colour complementation
sends $t$ to $2d_K-t$ and interchanges $S,T$, preserving both
quantities. If every $t(v)$ is even, every odd-component term is
zero, so continuous feasibility already gives an integer realisation.
One can also see this directly from a vertex of the capacity polytope.
Strictly fractional edges have independent incidence columns, hence
each of their components is a tree or has one odd cycle. Leaves
are impossible with even residual demands and all other weights at
0 or 2. The fractional support is thus a disjoint union of odd cycles;
each adjacent pair sums to 2, forcing all their weights to 1.
The resulting extreme point is integer. This is a fixed even-target
integrality fact, not permission to require all red degrees even in
the original conjecture.

Prescribed parity intervals also reduce to matching. Tighten
$[\ell(v),u(v)]\cap[0,2d(v)]$ to its first and last values
$g(v),f(v)$ of the desired parity. Add $(f(v)-g(v))/2$ virtual loops
at $v$ to $\doub G$. An exact factor of degree $f$ becomes the
required interval assignment after virtual loops are deleted, and
every such assignment can be completed by loops. For a matching
gadget with loops, create two ports per edge or loop, join each port
pair, and create $f(v)$ demand vertices joined to all local ports.
Ports are either matched together (edge unselected) or both to
demands (selected); exactly $f(v)$ incidences are selected. This
polynomial reduction includes degree feasibility only.

\subsection{Minimum conflicts and fixed-degree exchanges}

Let $\mathcal D(G)$ be all realisable integer red target vectors, and
\[
 \Psi(x)=\sum_{uv}\left(x_{uv}\mathbb{1}_{R(u)=R(v)}
 +(2-x_{uv})\mathbb{1}_{B(u)=B(v)}\right),\qquad
 F(q)=\min_{Bx=q}\Psi(x).
\]
For a fixed $q$ put $I_R(uv)=\mathbb{1}_{q(u)=q(v)}$,
$I_B(uv)=\mathbb{1}_{2d(u)-q(u)=2d(v)-q(v)}$. Then
\begin{equation}\label{eq:fixedcost}
 F(q)=2\sum_e I_B(e)+\min_{Bx=q}\sum_e (I_R(e)-I_B(e))x_e.
\end{equation}
The matching gadget above, with weight $-(I_R-I_B)$ on each
cross-port pair and zero on local edges, computes the minimum by
a maximum-weight perfect matching. This is an exact application of
weighted matching; the costs change when $q$ changes.

If two assignments realise the same target, mark their differing
copies by $+$ or $-$ and pair opposite signs at each vertex. They
decompose into even alternating closed walks. Altering one walk
moves toward the other assignment, never leaves $\{0,1,2\}$, and
preserves all degrees. Repeated vertices and repeated original edges
must be allowed. At fixed targets a legal realisation, if it exists,
can be reached without increasing $\Psi$, since all forced edges
move monotonically to their required extremes and all free edges
permit both colours.

More generally, for $q,r\in\mathcal D(G)$ and any unit step $s$
toward $r$, there is a second unit step $t$ toward $r$ such that
\[
 q+s+t,\ r-s-t\in\mathcal D(G),\qquad
 \|q+s+t-r\|_1=\|q-r\|_1-2.
\]
To prove this, pair opposite marked incidences as far as possible.
Unpaired incidences have exactly the signs of $r(v)-q(v)$. Start at
one representing $s$, follow its alternating paired trail, and stop
at the next unpaired incidence, representing $t$. Transfer the marked
copies in opposite directions in the two assignments. Every edge
stays between its two initial values. This is the standard constant
parity jump-system exchange phenomenon; see~\cite{BC95}. In particular
the two-unit target graph is connected, with distance
$\|q-r\|_1/2$. For a \emph{fixed} linear edge-cost vector $c$ the
same transfer gives
\[
 f_c(q+s+t)+f_c(r-s-t)\le f_c(q)+f_c(r),\quad
 f_c(q)=\min_{Bx=q}c\cdot x,
\]
because the two assignments have conserved total cost. The function
$F$ in~\eqref{eq:fixedcost} does not share this conclusion in general;
Section~\ref{sec:targetbarriers} gives explicit failures.

\subsection{A parity-defect lower bound}

For $E_0\cap E_2=\varnothing$, contract every component of the free
graph $K$. Let $\mathcal O$ be the resulting vertices with odd $q$
sum. On a forced edge write $z_e=x_e$ for $e\in E_0$ and
$z_e=2-x_e$ for $e\in E_2$. Then $\Psi=\sum z_e$, and each free
component $C$ satisfies
\[
 q(C)\equiv\sum_{e\in\delta(C)}z_e\pmod2.
\]
After duplicate copies cancel modulo two, the bad forced copies
contain an $\mathcal O$-join in the quotient graph. Therefore
$F(q)$ is at least its minimum join cardinality. Pairing incidences
decomposes any join into paths between paired odd vertices and
closed parts; conversely shortest paths of a minimum perfect matching
on the odd vertices, cancelled modulo two, give a join. This proves
the familiar shortest-path matching formula for this lower bound.
It detects parity only; internal factor capacities can require a
larger defect.

\section{Normal forms and conditional structural reductions}\label{sec:reductions}

\subsection{Sparse mixed-edge normal form}

\begin{proposition}\label{prop:normalform}
Every assignment can be changed, preserving all colour degrees and
not increasing $\Psi$, until the weight-one edges form a cactus
whose cycles are vertex-disjoint odd cycles.
\end{proposition}
\begin{proof}
On an even closed trail of mixed edges alternately increase and
decrease weights by one. Both signs of this change are feasible and
preserve degrees. At fixed degrees $\Psi$ is linear, so one choice
does not increase it. Every traversed edge ceases to be mixed.
Repeat until no even closed trail remains. There is no even cycle.
Two cycles sharing at least two vertices give a theta subgraph, two
of whose three paths have the same parity, giving an even cycle.
Two odd cycles sharing exactly one vertex give an even closed trail.
Thus cycles are vertex-disjoint and odd, and all other edges of each
mixed component are bridges.
\end{proof}

If there are $n$ vertices, $c$ mixed components including isolates,
and $k$ mixed cycles, their edge count is $n-c+k\le n-c+\lfloor n/3\rfloor$.
This sparsification neither asserts nor needs $\Psi=0$.

\subsection{Energy cleanup and an exact one-edge identity}

Define $\mathcal E(x)=\sum_v(R(v)-d_G(v))^2$. If a present colour
has equal degree $k>\max(d(u),d(v))$ at both ends of $uv$, changing
one copy of that colour to the other changes the energy by
\[
 -4k+2(d(u)+d(v))+2\le-2.
\]
If a present colour has equal degree one, its edge is mixed. Changing
the other copy to that colour changes the energy by
$6-2(d(u)+d(v))\le0$, provided the graph has no isolated $K_2$
component. Both new endpoint colour degrees are even, so no new
degree-one conflict appears. Hence the lexicographic pair
$(\mathcal E,h_1)$, where $h_1$ counts bad degree-one copies, strictly
decreases under these two rules. Termination takes at most
$(\lfloor\mathcal E(x)/2\rfloor+1)(2m+1)-1$ steps. The remaining
bad colour degrees are between two and the larger endpoint original
degree. The mixed-edge normal form can subsequently be imposed,
preserving these statements because it fixes all degrees and changes
only edges that already had both colours.

The energy is not $\Psi$. On a four-vertex path, changing the word
$(0,1,0)$ to $(0,2,0)$ lowers $\mathcal E$ from two to zero while
increasing $\Psi$ from two to six. The all-mixed triangle has energy
zero and six bad copies and is already stopped by both rules.

For completeness there is an exact local accounting identity. For a
permitted change $x_e\mapsto x_e+\eta$, $e=uv$, $\eta=\pm1$, put
\begin{align*}
 H_z^\eta(e)=\sum_{zw\ne e}\bigl[&x_{zw}
 (\mathbb{1}_{R(w)=R(z)+\eta}-\mathbb{1}_{R(w)=R(z)})\\
 &+(2-x_{zw})(\mathbb{1}_{B(w)=B(z)-\eta}-\mathbb{1}_{B(w)=B(z)})\bigr].
\end{align*}
Then
\[
 \Psi(x')-\Psi(x)=\eta(\mathbb{1}_{R(u)=R(v)}-
 \mathbb{1}_{B(u)=B(v)})+H_u^\eta(e)+H_v^\eta(e).
\]
On $e$ the degree differences stay fixed and only multiplicities
change; the remaining affected edges are exactly those in the two
sums. Any claimed improvement must account for these neighbour terms.

\subsection{Exact elimination of peripheral trees}

If $G$ contains a cycle, its 2-core $C$ is nonempty and connected.
Every component outside it is a tree attached by exactly one edge:
a cycle outside or two attachments would enlarge the minimum-degree-two
subgraph. At a core vertex put $h=d_C(v)\ge2$ and let
$\ell,s,b$ count attached trees of orders one, two and at least three,
respectively. Let $k=\ell+s+b,D=h+k$. At a proposed final red degree
$q$ put $B=2D-q$ and define
\begin{align*}
 L_1(q)&=\{c\in\{0,1,2\}:c=0\Rightarrow B\ne2,
 \ c=1\Rightarrow q,B\ne1,\ c=2\Rightarrow q\ne2\},\\
 L_2(q)&=\{c\in\{0,1,2\}:c=0\Rightarrow B\ne4,
 \ c=2\Rightarrow q\ne4\},\\
 S_v(q)&=\ell L_1(q)+s L_2(q)+b\{0,1,2\},\qquad
 J_v(q)=S_v(q)\cap[q-2h,q].
\end{align*}
The sums are Minkowski sums, with $0L=\{0\}$. A graph colouring is
equivalent to core weights $y_e$ and targets $q_v$ satisfying
\[
 q_v-\sum_{e\in C,e\ni v}y_e\in J_v(q_v)
\]
and the actual red/blue comparison on every core edge with total
degree $D_v$. Necessity is restriction. For sufficiency split the
chosen sum into allowed bridge weights and apply the tree interface
lemma independently. Physical counts hold because
$q=t+\sum c_i\ge c_i$ and
$B=(2h-t)+\sum(2-c_i)\ge2-c_i$.

Each $S_v(q)$ is an integer interval or an interval of step two.
Indeed $L_1$ is an interval or $\{0,2\}$, and $L_2$ is an interval,
possibly $\{1\}$. A width of at least one in any other interval
fills the gaps of a sum of $\{0,2\}$ sets. This is an exact compact
description, not an approximation of feasible root degrees. Trees
of order at least three may therefore be replaced by any rooted
three-vertex path, preserving colourability in both directions.

\subsection{A sharp tree capacity extension}\label{sec:treecapacity}

\begin{theorem}\label{thm:treecapacity}
At a core vertex of core degree $h\ge2$, $k\ge h$ attached rooted
trees can be restored to a given legal colouring of the graph with
those trees removed, preserving every old edge colour.
\end{theorem}
\begin{proof}
Fix its old internal red degree $t\in[0,2h]$. The feasible final
targets are $\mathcal Q(t)=\{q\in[t,t+2k]:q-t\in J_v(q)\}$.
For $k\ge3$ this is the whole interval, except for: value 1 if all
branches are singletons and $t=0$; its colour-complement at $t=2h$;
value 4 if there are exactly three branches, all of order two, and
$t=0$; its complement at $t=2h$. To check this list only the low
targets $1,2,4$ and their complements need inspection. At $q=1$
only singleton bridge sets become $\{0,2\}$, causing a gap exactly
when all branches are singletons and $t=0$. At $q=2$ their total
upper bound is $2k-\ell\ge k\ge3$, sufficient for any $2-t$.
At $q=4$ the upper bound is $2k-s$, insufficient only when
$k=s=3,t=0$. All other targets allow all weights. Low and high
cases are disjoint since $h+k\ge5$.

If $j$ old incident edges are mixed, old neighbours forbid at most
$h+j$ new targets. When $0<t<2h$, there are $2k+1>h+j$ candidates.
At $t=0$ or $2h$ no old edge is mixed and at least $2k>h$ remain.
The small case $h=k=2$ has at least four candidates, and all five
at $t=2$: directly inspect $q=1,2,4,6,7$ in the bridge sets.
At $t=1,3$ exactly one old edge is mixed, giving at most three
forbidden targets; at $t=0,4$ there are at most two. Thus an allowed
target always exists. Split its bridge contribution and fill the trees.
\end{proof}

The bound cannot uniformly be reduced to $k=h-1$ for any $h\ge2$.
Create roots $v,u_1,\ldots,u_h,z$, edges $vu_i$, and, for each
$i$, $d_i-1$ internally disjoint two-edge paths from $u_i$ to $z$,
where $d_i=2h-2+i$. The old degrees are $h,d_i,5h(h-1)/2$ at
$v,u_i,z$ and two at the path interiors. All-mixed old edges are
legal. Attach $h-1$ three-vertex rooted paths at $v$. A preserved
old colouring has $q\in[h,3h-2]$ and total colour sum $4h-2$ there.
The old edges forbid $q=d_i$ and $q=4h-2-d_i$, covering that entire
interval. Yet a full recolouring exists: change only $vu_2$ to zero,
give new bridges weight 1 and each new internal path word $(0,0)$.
Now $v$ has $(2h-2,2h)$, $u_2$ has $(2h-1,2h+1)$, other
$u_i$ keep $(d_i,d_i)$, and each tree has successive counts
$(1,3),(0,4),(0,2)$. These values verify every edge.

\subsection{Universal vertices and minimality}

\begin{theorem}\label{thm:universal}
If $z$ is universal in $G\ne K_2$, and every component of $G-z$
of order at least three has a given legal colouring, then the
colouring extends to $G$, preserving those given edge weights.
Isolated vertices and $K_2$ components require no given colouring.
\end{theorem}
\begin{proof}
Let $H=G-z$ have $N$ vertices, $k$ components $K_2$, $\ell$
isolates and larger components of orders $c_i$. Join every larger
component and isolate to $z$ with blue edges. For each $K_2=uv$
put $uv=1,zu=0,zv=2$. Its endpoints become $(1,3),(3,1)$ and
contribute $(2,2)$ to $z$. Thus $z$ has
$(2k,2(N-k))$. The $K_2$ interfaces compare even counts to three;
an isolated blue neighbour has count two, equal only in the excluded
case $G=K_2$. In a larger component, all old blue degrees increase
by two, preserving every old comparison, and are at most $2c_i$.
If $H$ has any other component, $2(N-k)>2c_i$, so the new
interfaces are legal.

It remains that $H$ is a single connected component of order $N\ge3$.
All-blue new edges fail only if an old vertex is universal in $H$
and has every old copy blue. If no such vertex exists, retain them.
Otherwise make every new edge red. This preserves old differences
and can fail only at an old all-red universal vertex. An all-blue
and an all-red universal vertex cannot coexist, since their common
positive-multiplicity edge would need both pure colours. Hence the
second choice works.
\end{proof}

Consequently a vertex-minimal counterexample has no universal vertex,
so $\Delta\le n-2$. It has a nonempty 2-core, every attached tree
has at most three vertices, and the number $k_v$ of trees at a
core vertex is smaller than its core degree $h_v$. These follow by
replacement and Theorem~\ref{thm:treecapacity}; the smaller remainders
contain a cycle and are not $K_2$. Hence, for core order and size
$n_C,m_C$,
\[
 |V(G)|\le n_C+3\sum_v(h_v-1)=6m_C-2n_C.
\]
Every maximal degree-two core thread has length at most five,
including return threads. At length at least six replace its interior
and attached trees by a new bare two-edge path, or a triangle for a
return. The smaller graph is simple and connected and has a legal
colouring by minimality. Its fixed endpoint counts and end weights
extend back by Theorem~\ref{thm:decorated}. The condition $k_v<h_v$
ensures that each internal core-degree-two vertex has at most one
tree. If the core is a cycle, its length is also at most five.
For a noncycle core compressed to a skeleton with $n_H$ vertices,
$m_H$ edges and minimum degree at least three, there are at most
$4m_H$ degree-two vertices. Combining the branch bound gives
\[
 |V(G)|\le22m_H-2n_H.
\]
The skeleton size and degrees remain unbounded; none of these
necessary conditions resolves an arbitrary minimum-degree-three core.

\section{Exact polynomial certificates and parity identities}\label{sec:algebra}

Orient each original edge $e=uv$ arbitrarily and put
$\Delta_e=d_G(u)-d_G(v)$ and $D_e=R(u)-R(v)$. The linear form
$D_e$ does not involve $x_e$, whose contributions cancel.

\begin{proposition}\label{prop:exactpoly}
The quadratic factor
\[
 f_e=D_e^2-2\Delta_e(2-x_e)D_e
       +2\Delta_e^2(x_e-1)(x_e-2)
\]
is nonzero exactly when the actual colour condition holds on $e$,
for $x_e\in\{0,1,2\}$. Consequently $\prod_e f_e$ has a nonzero
point on that grid if and only if $\doub G$ has a legal two-colouring.
\end{proposition}
\begin{proof}
Its evaluations at $x_e=0,1,2$ are respectively
$(D_e-2\Delta_e)^2$, $D_e(D_e-2\Delta_e)$, and $D_e^2$.
These express exactly the blue-only, mixed and red-only comparisons.
\end{proof}

For arbitrary real parameters $\tau_e$ define
\[
 f_{e,\tau}=D_e(D_e-2\Delta_e)+
 \tau_e(x_e-1)(D_e+\Delta_e(x_e-2)),
\]
and the multilinear parameter polynomial
\begin{equation}\label{eq:freeC}
 C_G(\tau)=\left[\prod_e x_e^2\right]
 \prod_e(D_e^2+\tau_e x_eD_e+\tau_e\Delta_e x_e^2).
\end{equation}
The three evaluations of $f_{e,\tau}$ are
$(D_e-2\Delta_e)(D_e-\tau_e)$,
$D_e(D_e-2\Delta_e)$, and
$D_e(D_e+\tau_e-2\Delta_e)$. Nonzero therefore always implies
actual legality, although additional legal assignments may be excluded.

\begin{proposition}\label{prop:freecertificate}
If $C_G$ is not the zero polynomial, $\lir(\doub G)\le2$.
Its nonvanishing is equivalent to nonvanishing at some
$\tau\in\{0,1\}^{E(G)}$. The fixed coefficient
$c_G=C_G(2\Delta)$ is a sufficient special case.
\end{proposition}
\begin{proof}
For every polynomial of total degree at most $2m$,
\begin{equation}\label{eq:gridinterp}
 \left[\prod_e x_e^2\right]P
 =\sum_{a\in\{0,1,2\}^{E(G)}}
 \frac{P(a)}{\prod_e w(a_e)},\qquad
 w(0)=2,\ w(1)=-1,\ w(2)=2.
\end{equation}
In one variable the weighted sums of powers of degrees zero and one
vanish and that of degree two equals one. A total-degree-at-most-$2m$
monomial with any exponent above two has another below two, so only
the all-two monomial contributes. This proves the identity. Apply it
to $\prod f_{e,\tau}$: its leading coefficient is~\eqref{eq:freeC}.
A nonzero coefficient forces a nonzero grid value. The assertion
about binary parameters follows from unique interpolation of
multilinear polynomials. This is an application of~\cite{Alon99};
no general nonzero-coefficient theorem is supplied here.
\end{proof}

Write $D=Mx$. On the binary-copy variables $z_{e,1},z_{e,2}$ put
$x_e=z_{e,1}+z_{e,2}$. The copy comparison is the nonzero linear
expression $D_e-2\Delta_e+2\Delta_e z_{e,i}$. Its coefficient matrix is
\[
 \mathcal A_G=M\otimes J_2+2\operatorname{diag}(\Delta)\otimes I_2.
\]
The coefficient of the full squarefree monomial is its permanent,
and
\begin{equation}\label{eq:permanent}
 \per(\mathcal A_G)=2^m c_G.
\end{equation}
To check this equality, the quadratic part of $f_e(a+b)$ differs
from the product of the two copy linear parts only by
$2\Delta_e^2(a^2+b^2)$, which cannot contribute to a squarefree
monomial. Each substitution $x_e^2=(a+b)^2$ contributes a factor
two to the desired coefficient. Reversing an edge reverses its two
copy rows, leaving the permanent unchanged.

\subsection{A principal-minor formula modulo two}

Let $A$ and $L$ be the adjacency and Laplacian matrices of the line
graph of $G$, indexed by original edges; all determinants in the
next statement are over $\mathbb F_2$.

\begin{theorem}\label{thm:modtwo}
For every $S\subseteq E(G)$,
\[
 [\prod_{e\in S}\tau_e]C_G(\tau)
 \equiv\det L[S,S]\,\det A[\overline S,\overline S]\pmod2.
\]
Moreover, $\det A[F,F]$ is odd exactly when $(V(G),F)$ is a
forest and each of its components has odd order, including isolates.
\end{theorem}
\begin{proof}
Modulo two, $M=A$, and
$\Delta_e\equiv d_{L(G)}(e)$, so $M+\operatorname{diag}(\Delta)=L$.
After selecting the indicated parameters the product is
\[
 \prod_{e\in S}x_e(D_e+\Delta_e x_e)
 \prod_{e\notin S}D_e^2.
\]
In characteristic two the squared factors contain only squared
variables. They cannot use an $S$ variable, already present once.
Each $S$ variable still needs one odd contribution, so all the
$|S|$ linear factors must use precisely the $S$ variables, one each.
The coefficient separates into the permanents of the two principal
submatrices; modulo two these equal the stated determinants.

Let $B$ be the unsigned vertex-edge incidence matrix modulo two.
Then $A=B^TB$. If $F$ contains a cycle its incidence columns are
dependent and so is its Gram matrix. For a tree component with $p$
vertices, every $(p-1)$-row incidence minor has determinant one,
by leaf elimination. Cauchy--Binet gives its Gram determinant
$p\pmod2$. Multiply over components. Isolates contribute the
empty determinant one. This proves the forest characterisation.
\end{proof}

For an odd-order tree the constant coefficient is odd. For an
even-order tree with an even-degree vertex there is an edge $uv$
with odd $\Delta$ whose deletion gives two odd-order components:
take a leaf even-degree vertex of the smallest subtree spanning all
even-degree vertices, and an edge leaving that subtree at that
vertex. The component on the other side has only original odd degrees,
and its original degree sum is $2(|U|-1)+1$, forcing $|U|$ odd.
Taking $S=\{uv\}$ in the formula gives an odd coefficient; the
constant coefficient is even. Similarly an odd-order unicyclic graph
has an odd single-parameter coefficient if a cycle edge has
endpoints of different degree parity. These are certificates on
already known colourable graph families, not new graph-family claims.

\subsection{Offset extensions on odd forests}\label{sec:oddforest}

\begin{corollary}\label{cor:oddforest}
If each component of a forest $F$ has odd order, arbitrary real
offsets $a(v),b(v)$ admit weights $0,1,2$ such that
\[
 a(v)+\sum_{e\ni v}x_e,\qquad
 b(v)+\sum_{e\ni v}(2-x_e)
\]
both distinguish the ends of every forest edge.
\end{corollary}
\begin{proof}
The product of both colour differences has leading homogeneous part
$(-1)^m\prod_e D_e^2$. Its all-two coefficient is odd by
Theorem~\ref{thm:modtwo}, independent of offsets. Apply
\eqref{eq:gridinterp} over the reals.
\end{proof}

For physical nonnegative integer offsets with $a(v)+b(v)$ even,
this extends constructively to any odd-order connected bipartite
graph $H$. Choose the bipartition labelling $p$ or its complement
so that $\sum p\equiv\sum a\pmod2$; odd order makes the two
choices have opposite parity. Solve
$\sum_{e\ni v}x_e\equiv p(v)-a(v)\pmod2$ by a spanning tree
with non-tree weights zero. The parity sum condition makes the root
equation hold. All weights may be 0 or 1, and the final red and
blue degrees both have parity $p(v)$. Every internal edge compares
opposite parities. Neither this result nor the corollary automatically
keeps outside edges legal when boundary degrees change.

\subsection{A fixed forbidden-value application}

If each vertex of a simple graph $H$ has real offset $a(v)$ and a
fixed forbidden set $Z(v)$ with $|Z(v)|\le d_H(v)$, there is a
weighting with $a(v)+\sum x_e\notin Z(v)$ at every vertex. Enlarge
each set to size $d_H(v)$ and consider the product of all forbidden
linear factors. Its highest homogeneous part is
$\prod_v(\sum_{e\ni v}x_e)^{d_H(v)}$, with nonnegative coefficients.
Selecting each incident variable once at every vertex shows that its
all-two coefficient is at least $\prod_v d_H(v)!>0$. Interpolation
proves the assertion. More exactly that coefficient equals
\[
 \frac{\prod_v d_H(v)!}{2^{|E(H)|}}
 \#\{\text{orientations of $\doub H$ with indegree $d_H(v)$ at each $v$}\}.
\]
Indeed allocating the two occurrences of each variable to its two
endpoint factors gives the multinomial expansion; distinguishing
the two copies contributes $2!$ per edge. This also witnesses
positivity by orienting one copy each way. With blue offsets too, combine red forbidden
values with the reflected blue values; the union must still have size
at most $d_H(v)$. Final neighbour degrees that themselves depend on
these edge weights are not fixed forbidden values, so this does not
prove the conjecture.

\subsection{Limitations of the parity certificate}\label{sec:mod2limits}

If all original degrees have the same parity, $L=A$ modulo two.
A symmetric zero-diagonal matrix has every odd-order principal
determinant zero, by pairing a permutation cycle with its reversal.
Thus when $m$ is odd every coefficient of $C_G$ is even. This does
not imply zero over the integers: the triangle has $c_G=-6$.

More generally for a connected $n$-vertex graph the line adjacency
matrix has rank $n-1$ at odd $n$ and $n-2$ at even $n$. Indeed the
incidence image is the even-coordinate-sum subspace of dimension
$n-1$; the kernel of $B^T$ consists of constants, intersecting that
subspace in dimension zero or one respectively. If $m$ exceeds
twice this rank and degrees have common parity, either $S$ or its
complement is larger than the rank in every coefficient formula.
All coefficients vanish modulo two. In particular this happens for
every connected $d$-regular simple graph with $d\ge4$. Such graphs
are known to be colourable~\cite{GOW26}; the obstruction is to this
parity certificate.

Even integer nonvanishing is a substantial demand. On a regular
graph $\Delta=0$, write $Q=\prod_eD_e$. If $C_G$ is nonzero then
$Q$ has a nonzero monomial $x^\alpha$ with all $\alpha_e\le2$,
since otherwise no product with $Q$ can supply the all-two monomial.
For arbitrary real vertex offsets the same coefficient survives in
$\prod_{uv}(a(u)-a(v)+D_{uv})$. Interpolating on an arbitrary
three-element list separately assigned to every edge gives a proper
neighbour-sum weighting with those offsets. Thus this nonvanishing
would imply the corresponding $(1,3)$ list-weighting property.
Ordinary shared-list 1--2--3 weightability does not establish it.
We leave the general nonzero-polynomial question open.

\section{A chromatic high-minimum-degree condition}\label{sec:highdegree}

We use the following published degree-constrained subgraph theorem
of Addario-Berry, Dalal and Reed~\cite{ADR08}, Theorem~5 in the
author version. For a simple host graph with degrees $d(v)$, integers
$a_v,b_v$ satisfying
\begin{equation}\label{eq:ADR}
 0\le a_v\le\lfloor d(v)/2\rfloor\le b_v<d(v),\quad
 b_v\le(d(v)+a_v)/2+1,\quad b_v\le2a_v+3
\end{equation}
admit a spanning subgraph $F$ with
$d_F(v)\in\{a_v,a_v+1,b_v,b_v+1\}$. We quote this theorem and
prove the residue and boundary consequences below.

\begin{theorem}[D1]\label{thm:chromatic}
If a simple graph has a proper $c$-vertex-colouring, $c\ge2$, and
\[
 \delta(G)\ge T(c):=8c(2c+1)-7,
\]
then it admits doubled weights $x_e\in\{0,1\}$ for which both
colour degrees distinguish the ends of every original edge. The same
internal conclusion holds with arbitrary prescribed physical external
counts at every vertex.
\end{theorem}
\begin{proof}
Put $p=2c,q=2c+1,M=pq$, and let $j(v)\in\{0,\ldots,c-1\}$
be the proper colours. With external counts $A_v,B_v$ and internal
degree $d=d_G(v)$ choose a residue $\phi_v$ by
\[
 \phi_v\equiv2j(v)-A_v\pmod p,\qquad
 \phi_v\equiv B_v+2d-2j(v)\pmod q.
\]
The moduli are coprime. Set $t=\lfloor d/2\rfloor$,
$r=(t-\phi_v)\bmod M$, $a=t-r$, and $b=a$ if $r=0$, otherwise
$b=a+M$. Both $a,b$ have residue $\phi_v$ and bracket $t$.
The assumption is $d\ge4M-7$, so
$\lfloor d/2\rfloor\ge2M-4$ and $\lceil d/2\rceil\ge2M-3$.
For $r>0$,
\[
 a\ge M-3,\quad b=a+M\le2a+3,\quad
 2b\le d+a+2\ \Longleftrightarrow\
 \lceil d/2\rceil\ge2M-r-2,
\]
and the last inequality holds since $r\ge1$. Also
$b\le t+M-1\le d-1$. For $r=0$, all conditions are immediate.
Thus~\eqref{eq:ADR} holds in the original simple graph.

Give $F$ edges weight 1 and all others zero. For some
$\epsilon_v\in\{0,1\}$ the final red and blue residues are
\[
 R(v)\bmod p\in P_{j(v)}:=\{2j(v),2j(v)+1\},\qquad
 B(v)\bmod q\in Q_{j(v)}:=\{2j(v),2j(v)-1\}.
\]
The $P_j$ are pairwise disjoint modulo $2c$. The $Q_j$ are also
pairwise disjoint modulo $2c+1$: $Q_0=\{0,2c\}$, and the others
are $\{1,2\},\{3,4\},\ldots,\{2c-3,2c-2\}$. Different
proper colours thus give unequal actual degrees in both colours.
The choice of external offsets changes no degree inequality, proving
the interface version too.
\end{proof}

For $c=2,3,4,5$ the thresholds are $73,161,281,433$. They are
sufficient bounds, with no optimality claim. The theorem has been
applied to a simple host only: on two vertices with 100 parallel
edges, the parameter pairs $(25,50)$ and $(30,55)$ satisfy the
numerical inequalities~\eqref{eq:ADR}, but the allowed factor-degree
sets at the two ends are disjoint, so no such multigraph factor
exists. Substituting the doubled host into the quoted theorem is
therefore not justified by these numerical conditions.

\subsection{One internal boundary vertex}

Suppose a connected induced region $H$ has all its outside edges
at one vertex $w$. Fix every edge outside the interior of $H$ and
assume that outside--outside edges are legal using their final counts.
If $h$ boundary edges include $m$ mixed ones, a pure boundary edge
forbids at most one residue label and a mixed edge at most two,
since the $P_j$ and $Q_j$ are disjoint. Let
$C=\max\{2,\chi(H),h+m+1\}$. At least one of the $C$ labels
avoids all outside colour counts. Map the old colour class of $w$
to that label and inject the other proper classes into remaining
labels. If $\delta(H)\ge T(C)$, Theorem~\ref{thm:chromatic} fills
$H$ and validates the boundary without changing any given edge.

In particular, a vertex-minimal counterexample cannot have a bridge
with one side $H$ of minimum internal degree at least
$T(\max\{\chi(H),3\})$. Replace $H$ by a leaf and use minimality,
then the preceding absorption with $h=1,m\le1$. If the smaller
graph is $K_2$, set the bridge mixed directly; its outside endpoint
has counts $(1,1)$ and no other outside edge to check. This treats
the exceptional smaller graph without assuming it colourable.

\subsection{Several CRT regions and low-side ports}\label{sec:crtregions}

Let disjoint induced regions $H_i=G[W_i]$ cover $W$ and put
$U=V(G)\setminus W$. Fix all edges other than those internal to
a single region, including edges between different regions. Assume
all $U$--$U$ edges are legal with their actual fixed endpoint counts.
Given a proper $c$-colouring $j$ of the \emph{entire} graph $G[W]$,
let $h_v$ be the number of edges from $v\in W$ to $U$, and $m_v$
the number of mixed ones. For every nonempty colour class define
\[
 L_a=1+\max_{j(v)=a}(h_v+m_v),\qquad C=\sum_a L_a,
\]
using $c\ge2$; unused classes may be omitted.

\begin{theorem}[D2: CRT region absorption]\label{thm:crtregions}
If every $\delta(H_i)\ge T(C)$, the fixed colouring extends over
all regions. If each high vertex has at most $t$ neighbours in $U$,
the uniform sufficient condition is $\delta(H_i)\ge T(c(2t+1))$.
\end{theorem}
\begin{proof}
Reserve disjoint blocks of $L_a$ labels for the proper colour classes.
At $v$, outside counts prohibit at most $h_v+m_v<L_{j(v)}$ labels.
Choose an allowed label in its reserved block, independently at each
vertex. All high--high edges, including inter-region edges, still
have different labels. Feed all fixed incident contributions as the
external counts into Theorem~\ref{thm:chromatic} in each region.
Internal and inter-region edges distinguish both residues; high--low
edges avoid their fixed prohibited residues. Low--low edges and
all their endpoint counts persist. The estimate $C\le c(2t+1)$
gives the uniform bound, using additional unused labels if necessary.
\end{proof}

There is a convenient way to supply the low precolouring. For every
component $D$ of $G[U]$, attach a private placeholder leaf for each
actual high--low edge, at its low endpoint. Different ports use
different leaves even if their actual high endpoints agree. A legal
colouring of this placeholder graph transfers the placeholder edge
weights to the actual boundary edges, preserving each low vertex
degree in both colours. If $D$ is bipartite or a cactus, its
placeholder graph remains in a class covered by~\cite{GW23,GW24}.
The only $K_2$ placeholder case is a single low vertex with one
port; make that port mixed, with low counts $(1,1)$ and no low--low
edge. Thus no $K_2$ colouring is invoked. Inter-region edges can be
fixed arbitrarily before absorption. This gives complete infinite
families with global minimum degree one or two, rather than only
globally high-degree graphs.

\subsection{Trees attached to an ordinarily irregular core}\label{sec:irregularcore}

\begin{proposition}[D2: matched tree attachment]\label{prop:irregularcore}
Suppose $H$ is connected, has an edge, and its ordinary degrees
differ at the ends of every edge. Let $S\subseteq V(H)$ have a
matching to $V(H)\setminus S$ covering $S$. At each vertex of $S$
attach one disjoint arbitrary rooted tree by a new bridge. The
doubled expanded graph has a legal two-colouring.
\end{proposition}
\begin{proof}
Give matching edges and the new bridges weight 1, all other core
edges zero. Core counts are
\[
 (R,B)=\begin{cases}
 (2,2d_H(v)),&v\in S,\\
 (1,2d_H(v)-1),&v\text{ is a matching helper},\\
 (0,2d_H(v)),&\text{otherwise}.
 \end{cases}
\]
Matching red edges join degrees two and one. Blue degrees at helpers
are odd, elsewhere even; within either parity group ordinary degree
irregularity makes them unequal. Each tree bridge meets final core
counts $(2,2d_H(v))$. A singleton tree has counts $(1,1)$; a larger
tree has three mixed-port candidates to avoid the two bridge
exclusions by Lemma~\ref{lem:treeports}. Fill every tree independently.
\end{proof}

A single specified root always has such a matching, by choosing
one neighbour. Arbitrary many roots need not, and the resulting
expanded graph need not remain ordinarily locally irregular, so
unconditional repeated attachment does not follow. Moreover a
\emph{given} old colouring can be rigid. Take edges $vu,vw$, two
two-edge $u$--$z$ paths and three two-edge $w$--$z$ paths. Its
degrees at $v,u,w,z$ are $2,3,4,5$ and the five interiors have degree
two. All-mixed edges give a legal old colouring. A leaf at $v$
has no possible new weight preserving that colouring: weights zero,
one and two cause respectively a blue equality with $w$, both
equalities with $u$, and a red equality with $w$. Changing old
$vw$ to zero and giving the new edge weight 1 repairs the graph.

\section{Dense graphs with antitone colour degrees}\label{sec:dense}

The second global degree condition uses Tutte's factor theorem directly
on a doubled graph, where parallel edges are allowed by that theorem.
This is separate from the simple-host theorem used in
Section~\ref{sec:highdegree}.

\begin{theorem}[D3]\label{thm:dense}
Let $G$ be simple with $n$ vertices and
\[
 8\delta(G)\ge7n+6.
\]
Then $\doub G$ has a two-colouring in which all vertex red degrees
are pairwise distinct and all vertex blue degrees are pairwise distinct.
\end{theorem}
\begin{proof}
Put $q=n-1-\delta(G)$, so the condition is $n\ge8q+14$.
Order $v_0,\ldots,v_{n-1}$ by nonincreasing ordinary degree and put
$L=\lfloor(n-1)/2\rfloor$. Set $f(v_i)=L+i$, except that if the
sum is odd, increase only $f(v_{n-1})$ by one. The increment occurs
exactly when $n\equiv2\pmod4$; only evenness is needed below.
Their sum is even and the targets are strictly increasing. Their
physical range follows explicitly from
\[
 0\le f(v_i)\le L+n\le2(n-1-q)\le2d_G(v_i),
\]
where $L\le(n-1)/2$ and $n\ge8q+14$ imply the middle bound.
If $\doub G$ has an $f$-factor, use it for red.
Then $B(v_i)=2d_G(v_i)-f(v_i)$ is strictly decreasing, proving the
asserted strong conclusion.

We verify the factor inequalities. For disjoint $S,T$, put
$s=|S|,t=|T|,h=s+t,r=n-h$. Their Tutte left side is
\[
 D=f(S)+2\sum_{v\in T}d_{G-S}(v)-f(T).
\]
The odd-component term is the number $o_f(S,T)$ of components of
$G-S-T$ with odd $f$ sum, since all boundary edges in $\doub G$
have even multiplicity. Independently of their positions in the order,
\begin{align*}
 f(S)&\ge sL+s(s-1)/2,\\
 f(T)&\le t(L+n-1)-t(t-1)/2+1,\\
 d_{G-S}(v)&\ge n-1-q-s.
\end{align*}
The last lower bound remains valid when negative. Writing
$L=(n-1)/2-\theta$, $0\le\theta\le1/2$, gives
\begin{align*}
 D&\ge (s-t)L+\frac{s(s-1)+t(t-1)}2
       +t(n-1-2q-2s)-1\\
  &\ge h\left(\frac{n-6}{4}-2q\right)-1
   \ge2h-1.
\end{align*}
For the middle inequality use
$s^2+t^2-4st\ge-h^2/2\ge-nh/2$,
$-\theta(s-t)\ge-h/2$, and $t\le h$.

If $h=0$, the graph is connected: each component would contain at
least $n-q$ vertices and $n>2q$. Its total $f$ sum is even, so
$o_f=0=D$. If $h\ge1$ and $r>2q$, the remaining graph is connected
by the same argument, now with each component of size at least
$r-q>r/2$. Thus $o_f\le1\le2h-1$. If $r\le2q$, then $o_f\le r$
and $h=n-r\ge n-2q\ge6q+14\ge r$, so again $D\ge o_f$.
All factor inequalities hold. Tutte's theorem~\cite{Tutte54} supplies
the red factor.
\end{proof}

No optimality of the constant is claimed. The construction may use
all three weights $0,1,2$, unlike Theorem~\ref{thm:chromatic}.
Searching for a particular factor is an implementation task; a search
timeout is neither a failed existence proof nor a counterexample.

\section{Dense forbidden targets and sequential absorption}\label{sec:denselists}

Let $H$ have order $n$ and $q=n-1-\delta(H)$. At each vertex fix
physical external contributions $A_v,B_v$ and a finite set $F_v$ of
forbidden \emph{final red degrees}. Put
\[
 D_v=d_H(v)+(A_v+B_v)/2,\quad
 \alpha=\min_v A_v,\quad a=\max_v A_v-\alpha.
\]
Partition the vertices into equal-$D_v$ classes $C_1,\ldots,C_k$ in
strictly decreasing $D$ order and put
\[
 b_i=\max_{v\in C_i}|F_v|,\quad
 E=\sum_i b_i+\min_i b_i+2,\quad W=\max\{a,E\}.
\]

\begin{theorem}[D5: forbidden-target dense interface]\label{thm:denselists}
If $n\ge8q+4W+10$, one can colour $\doub H$ so that every final
red degree avoids $F_v$, all final red degrees are pairwise distinct,
and all final blue degrees are pairwise distinct. The external
contributions are preserved.
\end{theorem}
\begin{proof}
Let $L=\lfloor(n-1)/2\rfloor$. Starting at $\alpha+L$, allocate
pairwise disjoint consecutive integer blocks in the specified class
order. The block for $C_i$ has length $|C_i|+b_i$, except that a
class of minimum $b_i$ receives an additional $b_i+2$ places.
Together these blocks form one interval of length $n+E$.
In each ordinary class assign distinct targets avoiding the
individual forbidden sets by greedy choice; the length exceeds the
number of assigned predecessors plus any forbidden set.

In the enlarged minimum class, designate a vertex $z$. Its interval
has at least $2b_i+3$ places, containing at least $b_i+1$ integers
of each parity. Reserve one allowed even and one allowed odd value
for $z$. Assign the other vertices distinct allowed values while
avoiding both reserved places. Writing $n_i=|C_i|$, at the last
such step at least
$(n_i+2b_i+2)-2-b_i-(n_i-2)=b_i+2$ places remain after all
exclusions; if $n_i=1$ there is no such step. Choose the parity at $z$ that makes
$\sum_v(R_v-A_v)$ even. The resulting final red targets are distinct.
Order vertices inside each degree class by increasing chosen target.
In the resulting global order $v_0,\ldots,v_{n-1}$, $D_v$ is
nonincreasing, $R_v$ is strictly increasing, and therefore
$2D_v-R_v$ is strictly decreasing.

The $n$ selected places in an interval of length $n+E$ have ordered
indices $j_i$ with $i\le j_i\le i+E$. For the internal target
$f(v_i)=R_{v_i}-A_{v_i}$ we obtain
\begin{equation}\label{eq:densetargetband}
 L+i-a\le f(v_i)\le L+i+E.
\end{equation}
The assumption gives $L-a\ge0$ and
$L+n-1+E\le2(n-1-q)\le2d_H(v)$, so these are physical internal
targets. Their total is even by the reserved parity choice.

For disjoint $S,T$ with $s,t,h,r$ as before,
\begin{align*}
 f(S)&\ge sL+s(s-1)/2-as,\\
 f(T)&\le t(L+n-1)-t(t-1)/2+Et.
\end{align*}
Using $as+Et\le Wh$ and the same degree estimate as in the dense
proof yields, without its extra final parity increment,
\[
 f(S)+2\sum_{v\in T}d_{H-S}(v)-f(T)
 \ge h\left(\frac{n-6}{4}-2q-W\right)\ge h.
\]
For $h=0$ the connected graph has even total target, so its odd term
is zero. For $h\ge1,r>2q$ the remaining graph is connected and has
odd term at most one. For $r\le2q$ it has at most $r$ odd
components, while $h\ge n-2q\ge6q+4W+10\ge r$.
Thus every factor inequality holds on $\doub H$. The red factor
realises $f$, and adding the fixed external counts realises the
chosen final $R$ and $2D-R$. Their order and forbidden-value choices
prove the theorem.
\end{proof}

Theorem~\ref{thm:dense} remains useful: with no external data or
forbidden values it needs $n\ge8q+14$, whereas the general theorem
has $E=2$ and needs $n\ge8q+18$. It is therefore not removed as
a supposedly weaker duplicate.

\subsection{Preserving arbitrary fixed boundary colours}

Let $H=G[W]$ be induced, all outside and boundary weights fixed,
and all outside--outside edges legal using their actual counts.
For each fixed edge $vu$, $v\in W,u\notin W$, put in $F_v$
the value $R(u)$ if $x_{vu}>0$, and the value
$2d_G(v)-B(u)$ if $x_{vu}<2$. Let $A_v,B_v$ be its total fixed
incident contributions. A final red degree avoiding $F_v$ makes
every boundary comparison legal, while $D_v=d_G(v)$.
Theorem~\ref{thm:denselists} therefore gives a preserved-colour
extension under its stated bound.

If each region vertex has at most $t$ outside edges, then
$a\le2t$, $b_i\le2t$, and the possible total degrees lie between
$n-1-q$ and $n-1+t$. Thus $k\le q+t+1$ and
$W\le2t(q+t+2)+2$. A uniform sufficient condition is
\begin{equation}\label{eq:denseuniform}
 n\ge8q+8t(q+t+2)+18.
\end{equation}
The bound uses the per-vertex number of boundary edges, not their
total number or the degrees of their outside endpoints. For $q=0,t=1$
it is 42, and for $q=3,t=1$ it is 90. The actual class-dependent
bound can be smaller, for example in the supplied 34-vertex clique
witness with fixed mixed ports.

\begin{corollary}[Sequential dense absorption]\label{cor:sequential}
Let $W_1,\ldots,W_s$ be disjoint induced regions. Fix every edge
not internal to one region, including all inter-region edges. Assume
the low--low edges outside the regions are legal with their fixed
final counts. If each region satisfies~\eqref{eq:denseuniform}
with its internal complement degree $q_i$ and per-vertex outside
edge bound $t_i$, all region interiors can be filled in any order,
preserving all fixed edge weights. The sharper actual $W_i$
condition of Theorem~\ref{thm:denselists} may be used instead at
each step.
\end{corollary}
\begin{proof}
At a region include all fixed outside incidences in its offsets,
including those toward unprocessed regions. Their contributions are
known even if the other endpoints' final degrees are not yet known.
Forbid conflicts with low vertices and already processed regions,
whose degrees are final. Defer comparisons toward future regions.
The dense theorem fills the present interior, verifying all its
internal and earlier boundary edges. Later fillings cannot alter
these degrees or the fixed boundary weights; they enforce the
deferred comparisons at their own steps. Induction over the chosen
order verifies every edge. The uniform bound allows the worst case
of at most $2t_i$ forbidden values per vertex at every step.
\end{proof}

The placeholder-leaf preprocessing from
Section~\ref{sec:crtregions} supplies low-side colours when its
components are bipartite or cacti, with the singleton-port exception
handled explicitly. Unlike the CRT region theorem, sequential dense
absorption requires no proper colouring of the entire union of high
regions. Neither theorem asserts that an arbitrary graph admits a
partition meeting its degree hypotheses.

\section{Four-vertex complement lifts}\label{sec:complements}

\begin{theorem}[D4: complement-path lift]\label{thm:complementlift}
Let a simple graph $H$ of order $n\ge3$ have a given legal doubled
colouring and $\Delta(H)\le n-3$. Let $uv$ be a nonedge whose
endpoints have different ordered colour-degree pairs. In the
complement replace $uv$ by the five-edge path $uabcdv$ and take
the complement again, obtaining $H^+$. Its colouring can be extended
preserving all old edge weights and increasing every old red and
blue degree by four. Moreover $\Delta(H^+)\le |H^+|-3$, and
all five edges of the new complement path have endpoints with
different ordered colour-degree pairs.
\end{theorem}
\begin{proof}
First suppose $B_H(u)\ne B_H(v)$. Give the new edge $uv$ weight
zero. For every old vertex $w$ give $wb,wc$ weight 2 and $wa,wd$
weight zero whenever those edges exist; $ua$ and $vd$ are absent.
Among new vertices the actual edges are $ac,bd,ad$; give them
weights $1,2,0$. Old vertices different from $u,v$ gain $(4,4)$
from their four new edges. At $u,v$, the extra old--old edge $uv$
replaces the missing blue contribution at $ua,vd$, so they also
gain $(4,4)$. All old comparisons persist and $uv$ is legal.

The new counts are
\[
 a:(1,2n+1),\quad b:(2n+2,0),\quad
 c:(2n+1,1),\quad d:(2,2n).
\]
Every old final colour count is at most $2(n-3)+4=2n-2$.
Consequently pure red old--$b,c$ edges and pure blue old--$a,d$
edges are legal. The three actual new--new edges are legal by their
displayed counts. The five missing path edges have different ordered
pairs as well, by the same strict bounds and the distinct new pairs.
Each old degree increases by four, and every new degree is $n+1$,
which proves the maximum-degree assertion. If only the red counts
of $u,v$ differ, complement the old colours temporarily, apply this
construction, and complement the entire result back. Old weights
are again preserved and all stated properties remain true.
\end{proof}

Different boundary pairs are a real hypothesis. On $C_5$ the cyclic
word $(0,0,1,1,2)$ is legal, but two nonadjacent vertices have the
same pair $(2,2)$. A construction adding the same $(4,4)$ to each
and adding an actual edge between them cannot make that edge legal
with any weight. This restricts this lift interface, not the graph.

\subsection{An explicit nonregular infinite family}

For odd $k\ge3$ let
$F_k=C_{2k+1}+02$, using cyclic labels $0,\ldots,2k$, and
$G_k=\overline{F_k}$. Its two chord endpoints have degree $2k-3$
and the other vertices degree $2k-2$; its size is $2k^2-k-2$.
For $k=3$, in the edge order
\[
 03,04,05,13,14,15,16,24,25,26,35,36,46
\]
use weights
\[
 2,2,2,0,0,0,0,0,2,0,1,0,1.
\]
The vertex red degrees are $(6,0,2,3,3,5,1)$ and blue degrees
$(0,8,4,5,5,3,7)$, verifying actual legality. The complement edge
$60$ has different blue endpoint counts. Replacing it by a
five-edge path produces $G_5$ with new labels extending the cyclic
order, keeping chord $02$. Repeat on the last new edge incident
with 0. Its new endpoint blue degree is $2n$, greater than the old
endpoint final bound $2n-2$, so the interface stays available.
The lift proves colourability of every $G_k$, $k$ odd, by induction.

This family also limits a parity-refined defective partition. Suppose
one requires a partition $(A,B)$ with
\[
 \Delta(G_k[A])\le k-1,\quad \Delta(G_k[B])\le k-2,
 \quad R|_A\text{ even},\quad R|_B\text{ odd}.
\]
The degree conditions imply minimum induced complement degrees
$|A|-k$ on $A$ and $|B|-k+1$ on $B$.
Only two vertices of $F_k$ have degree three; hence these bounds
force $|A|\le k+2$, $|B|\le k+1$, and
$|B|\in\{k-1,k,k+1\}$. A minimum-degree-two induced subset of
$F_k$ with fewer than $2k$ vertices can only be its triangle $012$:
if it contains any vertex on the other arc from 2 to 0, its two
neighbours are forced, propagating along that entire arc of $2k$
vertices. If $|B|=k-1$, then $A$ has size $k+2\ge5$ and induced
minimum degree two, impossible; if $|B|=k+1$, the same holds for
$B$ of size at least four. Therefore $|B|=k$ is odd, contradicting
the even red handshake sum. The degree partition itself exists:
$B=\{0,\ldots,k-1\}$ and its complement both have no isolates in
$F_k$, which verifies the required complement degree bounds. It is
the appended parity prescription that fails. The legal colourings
above show that this is not a conjecture counterexample.

\subsection{Lifting complements of cubic graphs}

Let $F\ne K_4$ be a connected edge-transitive cubic simple graph
of order at least six and $H=\overline F$. Each component of $H$
is regular of degree $n-4\ge2$, so it has a legal doubled colouring
by~\cite{GOW26}. Some edge of $F$ has distinct endpoint pairs:
otherwise connectedness of $F$ would make the pair constant on all
vertices, contradicting legality of any actual edge of $H$.
An automorphism carries that complement interface to any selected
edge of $F$. Repeated use of Theorem~\ref{thm:complementlift}
therefore handles replacing that selected edge of $F$ by a path
of length $4t+1$ and then complementing, for any $t\ge0$. At
$t>0$ these graphs are nonregular. Edge-transitivity is used to
control the selected initial edge and is not silently omitted.

\subsection{Restrictions in complement degree two or three}

Suppose $G$ is a vertex-minimal counterexample to the general
conjecture and every vertex of $F=\overline G$ has degree two or
three. If $n\ge11$, $F$ cannot contain a path $uabcdv$ whose four
internal vertices have degree two, whose endpoint degrees differ,
and with $uv\notin E(F)$. Delete the four internal vertices and
replace that path by $uv$. The resulting complement $H$ has order
$N=n-4\ge7$, degree between $N-4$ and $N-3$, and is connected
because $2(N-3)>N$. Minimality gives a legal colouring of $H$.
The two endpoint ordinary degrees differ, so their ordered colour
pairs differ. The lift reconstructs $G$, a contradiction.
It follows that no path starting at an $F$-degree-three vertex can
continue through six consecutive distinct degree-two vertices:
use its first five-edge subpath, whose endpoint degrees are three
and two; the sixth vertex ensures the missing closing edge.

There is also no cycle component of $F$ of length at least nine.
Shorten that cycle by four. Its length remains at least five, so
its complement is connected (and is completely joined to any other
complement components). Minimality supplies a colouring of the
smaller actual graph. Some edge of the shortened complement cycle
has distinct endpoint pairs; otherwise every vertex of that cycle
would have the same pair, contradicting an actual edge within its
complement. Rotate that cycle component to place the interface at
the edge to be lengthened and apply the lift. Finally
Theorem~\ref{thm:dense} covers this entire complement-degree class
when $n\ge38$ since $q\le3$. Any hypothetical counterexample in
this class therefore has $n<38$. No exhaustion of that residual
finite range, or reduction of all graphs to this class, is claimed.

\section{Cut, parity and modular templates}\label{sec:templates}

Every obstruction in this and the next sections concerns an additional
restriction on a colouring or on a proposed algorithm. We supply
actual legal colourings or cite an already established family result
where needed to keep that distinction explicit.

\subsection{A maximum-cut template with no exceptions}

The template chooses a maximum cut $(S,T)$, requires all edges in
$G[S]$ to have weight zero, and requires $R(v)\equiv\mathbb{1}_T(v)$
at every vertex. For any vertex set $U$ the actual handshake identity is
\begin{equation}\label{eq:boundaryhandshake}
 \sum_{v\in U}R(v)=2\sum_{e\in E(U)}x_e+
 \sum_{e\in\delta(U)}x_e.
\end{equation}
With parity-exception set $X$, the last sum has parity
$|U\cap T|+|U\cap X|$. A single bridge of odd required parity
must therefore be mixed. This is a necessary boundary constraint,
not a guarantee of an extendible interface.

\begin{proposition}\label{prop:cutobstruction}
For every odd $r\ge3$, take a triangle $abc$ and attach one copy
of $K_{r,r}$ to each triangle vertex by a bridge to a distinguished
module root. The resulting graph $G_r$ has minimum degree three
and a legal doubled two-colouring, but no maximum cut admits the
zero-exception template above. Exactly one exception suffices.
\end{proposition}
\begin{proof}
The graph has $3+6r$ vertices and $3r^2+6$ edges. A maximum cut
has size $3r^2+5$, since each module contributes at most $r^2$,
each bridge at most one, and the triangle at most two; orienting
the module bipartitions independently attains all these bounds.
Every maximum cut therefore cuts every module edge and every bridge.
Each module has $r$ vertices in $T$, odd. With no exceptions,
\eqref{eq:boundaryhandshake} forces each module bridge to have
weight 1. The total red sum is even; the three modules contribute
an odd number of $T$ vertices, so precisely one triangle vertex is
in $T$. The other two triangle vertices lie in $S$, their common
edge has weight zero, and their edges to the $T$ vertex must both
have weight 1 to give them even red sums after the mixed bridges.
They both have counts $(2,4)$, contradicting their common blue edge.

Here is a full legal colouring. Give $ab,bc,ca$ weights $0,1,2$
and all bridges weight 1. Then
$a:(3,3), b:(2,4), c:(4,2)$. In each module use rows and columns
$0,\ldots,r-1$, with root at row zero. In the $a$ module give the
root row all weights 1 and all other entries zero; the root is
$(r+1,r+1)$, other rows $(0,2r)$ and columns $(1,2r-1)$.
For the $b$ module start with root entry $(0,0)$ of weight 1,
row 1 entries in all other columns of weight 1, all remaining
entries zero, and complement its internal weights. The root is
then $(2r,2)$; row counts are even and column counts odd. For the
$c$ module give the root row zero, every other row weight 1 only
in column zero, and all other entries zero. Its root is
$(1,2r+1)$; other rows $(1,2r-1)$, column zero $(r-1,r+1)$,
other columns $(0,2r)$. Row--column parity verifies every module
edge and the displayed root pairs verify the bridges. The triangle
edges are checked directly.

Choose the maximum cut with $a,b\in T,c\in S$, module roots on
opposite sides to their triangle neighbours. The module patterns
have precisely the cut parities, as do $a,c$. Only $b$ is an
exception. Thus the minimum exception count in this template is one.
\end{proof}

The same graph also excludes the different prescription
$R(v)\equiv d_G(v)\pmod2$ at \emph{every} vertex, without mentioning
a cut. Each module has odd sum of original degrees, so its bridge
would again be mixed. The three triangle vertices would then have
distinct odd red degrees, necessarily $1,3,5$, because they have
equal original degree three. Subtracting the mixed bridge contribution
would require internal triangle red sums $0,2,4$, impossible at the
edge between the zero and four vertices. This new restriction also
fails despite the explicit legal colouring in the proof.

\subsection{Small boundary modules and parity transfer}

There is a seven-vertex module consisting of triangle $abc$,
two-edge pendant paths at $b,c$, and a fixed mixed outside edge
at $a$. If its internal cut is maximum, all four path edges cross
and two triangle edges cross. Under the zero-exception template,
the first edge of a two-edge pendant must be mixed: its leaf and
triangle vertex lie on the same side, and the middle parity forces
that first weight odd. The module has an odd number of $T$ vertices
by its mixed boundary, hence one triangle vertex in $T$. Each of
the two $S$ vertices has one mixed outside contribution, so its
edge to the $T$ vertex is mixed too. Their common blue edge again
joins counts $(2,4)$. Thus fixed mixed boundaries need not extend
even when total parity is compatible.

Put two-edge pendant paths at all three triangle vertices to obtain
a nine-vertex graph $O$. Every maximum cut and zero-exception
template fails by the same argument. For clarity, in edge order
$ab,ac,bc,aa_1,a_1a_2,bb_1,b_1b_2,cc_1,c_1c_2$ a legal word is
$(1,2,0,1,0,0,0,1,1)$. With
$T=\{b,c,a_1,b_2,c_2\}$ it has exactly the one exception $b_2$.
This small existence statement is subsumed by the stronger
minimum-degree-three family above, but its boundary properties differ.

In $O$ choose root $a_2$ and give it arbitrary outside total red
contribution $A$. Internal maximum-cut template feasibility with no
exceptions requires $A$ odd. If the pendant at $a$ instead has
length four and its endpoint is the root, feasibility requires $A$
even. To see both assertions, for a pendant of length $2t$ apply
\eqref{eq:boundaryhandshake} to its $2t$ nontriangle vertices.
Its first weight $q$ satisfies $q+A\equiv t\pmod2$. The wrong
boundary parity in either case forces $q=1$ and one triangle vertex
in $T$, repeating the $(2,4)$ conflict. These are necessary
conditions only; sufficiency of every boundary with that parity is
not claimed.

\subsection{A six-vertex modulo-three obstruction}

Take triangles $uab$ and $vcd$ joined by bridge $uv$. Even though
this cactus is three-colourable, one cannot require every present
red and blue comparison to be nonzero modulo three. Write the left
weights $ua,ub,ab,uv$ as $p,q,t,z$ and put $h=z-t\pmod3$.
The degree-two edge $ab$ forces $p\ne q$. On $ua$ the restrictions
are respectively $q+h\ne2$, $q+h=1$, $q+h\ne0$ when
$p=0,1,2$. The $ub$ conditions interchange $p,q$. Unordered pairs
$\{0,1\}$ and $\{1,2\}$ immediately contradict these restrictions;
only $\{p,q\}=\{0,2\}$ remains. The right triangle obeys the
same rule, so $R(u)\equiv R(v)\equiv2+z$ and
$B(u)\equiv B(v)\equiv1-z$. Every bridge weight then fails.
In contrast the actual legal word in order
$ua,ub,uv,ab,vc,vd,cd$ is $(1,2,2,0,2,0,1)$, giving counts
$u:(5,1),a:(1,3),b:(2,2),v:(4,2),c:(3,1),d:(1,3)$.
For instance a legitimate red degree difference of three need not
survive modulo three. The modulo-two certificate has its own
independent restrictions in Section~\ref{sec:mod2limits}.

\section{Barriers to nonincreasing local descent}\label{sec:descent}

\subsection{A closed two-state single-edge plateau}

On the wheel with centre 0 and rim $123451$, use edge order
\[
 01,02,03,04,05,12,15,23,34,45.
\]
The words
$x=(1,2,1,1,0,0,2,2,0,0)$ and
$y=(1,2,1,1,0,0,2,2,2,0)$ both have $\Psi=1$.
All single-edge changes, allowing even direct $0\leftrightarrow2$
changes, have the following conflict values. Each triple gives the
result for new weights $0,1,2$.
\begin{center}\small
\begin{tabular}{ccc}\toprule Edge&From $x$&From $y$\\\midrule
$01$&$(4,1,4)$&$(4,1,4)$\\
$02$&$(2,4,1)$&$(2,2,1)$\\
$03$&$(2,1,4)$&$(4,1,4)$\\
$04$&$(4,1,4)$&$(4,1,2)$\\
$05$&$(1,2,2)$&$(1,4,2)$\\
$12$&$(1,3,2)$&$(1,5,2)$\\
$15$&$(2,5,1)$&$(2,3,1)$\\
$23$&$(4,3,1)$&$(2,2,1)$\\
$34$&$(1,4,1)$&$(1,4,1)$\\
$45$&$(1,2,2)$&$(1,3,4)$\\\bottomrule
\end{tabular}
\end{center}
The only nontrivial nonincreasing steps exchange $x,y$. By induction
every such process remains in this two-state set. A simultaneous
change of $01$ from 1 to zero and $12$ from zero to 1 gives
$(0,2,1,1,0,1,2,2,0,0)$ with $\Psi=0$. The complete finite table
proves a barrier for single-edge monotonicity, including neutral
steps; it does not exclude larger simultaneous repairs.

\subsection{Strict positive minima after fixed-target optimisation}

On a $d$-regular graph every assignment realising $q$ has
\[
 F(q)=2|\{uv:q(u)=q(v)\}|.
\]
Suppose equal-target edges are a perfect matching, every available
target $q(v)\pm1$ occurs on at least three neighbours, and every
available $q(v)\pm2$ on at least two. Then $F(q)=n>0$ and
every realisable target at $\ell_1$ distance two has larger $F$.
Changing one coordinate by two creates at least two equalities and
removes at most one. Changing two coordinates by one creates at
least two equalities to unchanged neighbours at each coordinate,
removing at most the two old matching edges. Thus no neutral first
step is available, even after completely optimising each fibre.

Here is an explicit infinite realisable family for that criterion.
For any even $N\ge4$, take vertices $(i,t)$ with $0\le i\le26$
and $t\in\mathbb Z_N$, target $q(i,t)=i$. Put a perfect matching
within each layer by pairing consecutive even/odd $t$ labels. Between
layers start with $a_{ij}=3$ when $|i-j|=1$, 2 when $|i-j|=2$,
and zero otherwise. Add each increment in the first table below and
its reflection $(26-j,26-i)$, then make $a$ symmetric.
\begin{center}\small
\begin{tabular}{rrr@{\qquad}rrr}\toprule
$i$&$j$&increment&$i$&$j$&increment\\\midrule
0&1&1&0&2&1\\0&4&2&0&11&2\\0&24&1&1&3&2\\
1&20&1&5&14&2\\6&13&1&7&9&2\\8&10&2&&&\\\bottomrule
\end{tabular}
\end{center}
Its rows sum to 12 and $0\le a_{ij}\le4$. For $i<j$, add the
$a_{ij}$ matchings $(i,t)(j,t+s)$, $s=0,\ldots,a_{ij}-1$.
They are distinct for $N\ge4$. Together with the within-layer edge
the graph is 13-regular. Shifts zero and one between consecutive
layers make it connected. Equal targets occur exactly on the
within-layer perfect matching and the required neighbour counts hold.

To realise $q$, for $i<j,i+j\le26$ set $b_{ij}=0$ except
\[
\begin{array}{c|rrrrrrrrr}
(i,j)&(3,4)&(5,7)&(6,7)&(8,9)&(9,11)&(10,11)&(10,12)&(11,13)&(12,14)\\
b_{ij}&2&3&4&6&1&6&4&4&2
\end{array}
\]
and define $b_{ij}=2a_{ij}-b_{26-j,26-i}$ for $i+j>26$;
symmetrise. Within-layer weights at $i=0,\ldots,13$ are
\[
 (c_0,\ldots,c_{13})=(0,1,0,1,2,2,0,0,2,2,0,0,2,1),
\]
with $c_{26-i}=2-c_i$. Direct row addition verifies
$0\le b_{ij}\le2a_{ij}$ and $c_i+\sum_{j\ne i}b_{ij}=i$.
Give shift-$s$ edges between layers weight
$\min\{2,\max\{0,b_{ij}-2s\}\}$ and within-layer edges weight
$c_i$. Each vertex receives total $i$, as required. The tables are
a fixed integer certificate, independently checked in the artifact.
Thus for this connected infinite family of order $27N$,
$F(q)=27N$ and every two-unit neighbour has $F\ge27N+2$.
These regular graphs have other legal targets by~\cite{GOW26}.
This bad initial target cannot itself arise by subtracting degrees
from a proper 1--2--3 neighbour-sum weighting, since its equal-target
matching edges have equal original degrees.

\subsection{A proper 1--2--3 initialisation can also be trapped}

On the five-vertex path use $x=(1,0,0,1)$. Then $w=x+1$ has
proper neighbouring sums $(2,3,2,3,2)$, but the doubled target
$q=(1,1,0,1,1)$ has $F(q)=2$. A path target determines its weights
uniquely by leaf recursion. Its alternating bipartition sum is zero,
so a two-unit target change must be two opposite-balanced unit
changes, corresponding to an alternating contiguous path of weight
changes. The two adjacent internal zero weights preclude all such
paths except the following eight:
\begin{center}\small
\begin{tabular}{cc@{\qquad}cc}\toprule
New word&$F$&New word&$F$\\\midrule
$(0,0,0,1)$&3&$(2,0,0,1)$&3\\
$(1,1,0,1)$&3&$(1,0,1,1)$&3\\
$(1,0,0,0)$&3&$(1,0,0,2)$&3\\
$(0,1,0,1)$&5&$(1,0,1,0)$&5\\\bottomrule
\end{tabular}
\end{center}
No nonincreasing two-unit first step exists, even with an initially
proper 1--2--3 weighting. The word $(0,0,1,1)$ is legal and has
target distance four from $q$, attaining the least possible
nonincreasing first-step distance. This does not prove that distance
four is sufficient for general graphs or specially selected initialisations.

\section{Fixed-target and radius obstructions}\label{sec:targetbarriers}

\subsection{Bridge locking, capacity and pure odd defects}

For a bridge $e=uv$ whose $u$ side $U$ is bipartite with signs
$\varepsilon$, fixed internal target demands give
\[
 x_e=\varepsilon(u)\sum_{v\in U}\varepsilon(v)q(v).
\]
Internal contributions cancel. Thus a bridge may be fixed uniquely;
on a forest the entire weighting is unique for a fixed target.
Take a triangle $uab$ with path $u w l$, in edge order
$ua,ub,uw,ab,wl$. The word $(0,1,1,1,1)$ is realisable and
comes from a proper 1--2--3 weighting, but the bridge $uw$ joins
equal red degree two, so legality would require $x_{uw}=0$.
Its bipartite side forces $x_{uw}=q(w)-q(l)=1$. No fixed-target
repair exists. The same graph has the legal word $(1,0,0,2,0)$.

For a purely integer obstruction, take triangles $012,345$, bridge
$03$, and leaves $26,57,38$. In edge order
\[
 01,02,12,34,35,45,03,26,57,38
\]
use $(0,1,1,0,1,1,1,2,2,0)$. Its targets and blue degrees are
\[
 q=(2,1,4,2,1,4,2,2,0),\quad B=(4,3,2,6,3,2,0,0,2).
\]
The sums $q+d=(5,3,7,6,3,7,3,3,1)$ distinguish all neighbours.
Only bridge $03$ is forced blue, and there are no forced red edges.
After deleting that bridge, the two free components each have target
sum nine. At $S=T=\varnothing$ the factor certificate is
$g=0,o=2$. However giving each triangle the word
$(1/2,3/2,1/2)$, setting the bridge zero and retaining the leaf
weights is a continuous legal realisation of exactly $q$. Hence
continuous feasibility at integer targets need not permit integer
rounding. A legal integer repair changes $02,35$ to 2 and the
bridge to zero, changing the targets at 2 and 5 by one. The actual
new red degrees are $(2,1,5,2,1,5,2,2,0)$.

The parity lower bound alone can also be strict. On edge set
$01,02,03,12,34$ the word $(1,0,0,0,0)$ has all free-component
target sums even, but vertex 1 needs red degree one through its
only free neighbour 2 of red target zero. The capacity pair
$S=\{2\},T=\{1\}$ has $g=-1$. Thus its quotient join bound is
zero while $F(q)=1$.

\subsection{Failure of exchange convexity}

For a triangle the targets uniquely determine weights by
$x_{ij}=(q_i+q_j-q_k)/2$. The targets $(1,2,3)$ and $(3,2,1)$
are legal. Starting from the first toward the second with step
$e_0$, the possible second steps are $e_0$ and $-e_2$.
The first gives target pair $(3,2,3),(1,2,1)$, each with $F=2$;
the second gives two copies of $(2,2,2)$, each with $F=6$.
Neither satisfies the fixed-cost exchange inequality with initial
total zero. In particular the midpoint of two legal targets can
be realisable with positive conflict. The favourable classical
exchange properties in Section~\ref{sec:targets} cannot be assigned
to the changing conflict-cost function $F$.

\subsection{An unbounded fixed-parity repair radius}\label{sec:radius}

\begin{proposition}\label{prop:radius}
Let even $c\ge2$ and $k\ge1$. Take $c$ disjoint cycles of length
$L=4k+2$ and two nonadjacent hubs $a,b$, each joined to every cycle
vertex. Put $N=cL$. Give the hubs target zero and every cycle
vertex target four, realised by cycle weights 2 and hub weights
zero. Fix the maximum cut with both hubs in $S$ and alternate
$S,T$ around each cycle, and require final red parity zero on $S$,
one on $T$. The least maximum coordinate change from those targets is
\[
 \rho_{\min}=2\lceil c/4\rceil,
\]
and this bound is achieved by legal doubled colourings. With $E$
parity exceptions and maximum change at most $M$, necessarily
$2M+E\ge c$.
\end{proposition}
\begin{proof}
The cut has size $2N$ and is maximum. If hubs are on opposite sides,
the hub edges contribute $N$ and cycle edges at most $N$. If hubs
share a side and $t$ cycle vertices share it, hub edges contribute
$2(N-t)$ and cycle edges at most $2t$. The alternating cut reaches
$2N$. Each cycle has $2k+1$ odd-labelled vertices. Its red boundary
sum is therefore a positive odd integer by
\eqref{eq:boundaryhandshake}. Hence $R(a)+R(b)\ge c$. Both hub
targets have required even parity, forcing the larger to be at
least $2\lceil c/4\rceil$.

Write each cycle $s_0,t_0,\ldots,s_{2k},t_{2k},s_0$, starting
with all cycle weights 2. Change one hub edge to $t_0$ to weight
1. For $i=1,\ldots,k$ change both edges incident with $s_{2i}$
along its two adjacent $t$ vertices to 1. The counts are
$t_0:(5,3)$, other $t$ vertices $(3,5)$, even positive-index
$s$ vertices $(2,6)$, and other $s$ vertices $(4,4)$.
Assign $2\lfloor c/4\rfloor$ mixed hub edges to $a$ and the rest
to $b$. Hub red degrees are those two even numbers, with blue
degrees at least $2N-c\ge22$. All cycle and mixed hub comparisons
are parity-separated. The pure blue hub edges to same-side $s$
vertices compare a degree at least 22 to one at most six. The
maximum change is exactly the asserted radius.

With $E$ exceptions, at least $c-E$ cycles have no exception and
each still needs at least one boundary red copy. The two hub red
degrees are each at most $M$, giving $2M\ge c-E$ independently
of any local irregularity requirement.
\end{proof}

This is a minimum-degree-four, 2-connected family with no peripheral
trees or degree-two chains. The starting target is feasible but is
not a legal colouring. The radius conclusion fixes the stated cut
parities. Without fixing them, keep hub weights zero and use the
cyclic word $(0,1,2,0,1,2)(0,1,1,0)^{k-1}$ on each cycle.
Positions two apart differ, so these words are legal; all changes
from the initial target are at most four, with unrestricted parity
exceptions. Thus the lower bound is not a distance statement about
all legal colourings.

Even an everywhere-one centre window is too narrow on $C_5$:
some red degree in any legal colouring is zero or four. If all red
degrees were in $\{1,2,3\}$, the number of degree-two vertices
would be odd by the handshake sum, and at most two by independence,
hence exactly one. The other four form an alternating 1--3 path.
An internal degree-one vertex between two degree-three vertices
would need two incident weights each at least one, impossible.
The word $(0,0,1,1,2)$ supplies a legal colouring. Translating all
degrees by fixed equal external offsets gives the corresponding
fixed-interface small-window obstruction.

\section{Sharp interface limits and multiple marks}\label{sec:interfacebarriers}

The path-length-six threshold in Theorem~\ref{thm:decorated} is
sharp for arbitrary fixed physical endpoint data. On a five-edge
bare path with end weights 1 and both endpoint counts $(3,3)$,
write the word $(1,a,b,c,1)$. Each mixed end edge forces $a=c=1$.
The middle edge then joins identical counts $(1+b,3-b)$ for any
$b$, and fails. The data are physically realised by adding two
mixed leaf edges at each endpoint. That whole ten-vertex tree is
colourable by Lemma~\ref{lem:treeports}; only the fixed interface
is impossible. The tree and clique capacity obstructions of
Sections~\ref{sec:treecapacity} and~\ref{sec:cliquecapacity} independently
show sharpness of their preserved-colour thresholds.

An arbitrary-root-sum interface also fails for general graphs. On
triangle $vabv$ with leaf $vz$, requiring the internal root red
sum zero forces all three incident weights zero. The endpoints of
$ab$ then have identical red and blue degrees whatever its weight.
Sum six is impossible by complementation. Nonetheless the word
$(x_{va},x_{vb},x_{ab},x_{vz})=(0,2,1,0)$ is legal, with
$(R(v),R(a),R(b),R(z))=(2,1,3,0)$ and blue degrees $(4,3,1,2)$.
The full clique interface is stronger than ordinary colourability.

Multiple fixed bridges need not preserve the one-bridge candidate
counts. On rooted path $r-u-v$, give $r$ external counts $(2,2)$
from two mixed bridges. If the internal word is $(a,b)$, the red
degrees are $(2+a,a+b,b)$ and blue degrees $(4-a,4-a-b,2-b)$.
At $a=0$ the edge $uv$ forces $b=0$, then $ru$ fails in blue;
at $a=2$ it forces $b=2$, then $ru$ fails in red. At $a=1$
the first edge excludes $b=0,2$, leaving exactly $(1,1)$.
Thus the internal root candidate set is the singleton $\{1\}$.

Finally consider a cubic-root safe-core requirement that every root
red degree belong to $\safe=\{0,1,3,5,6\}$ and every marked
mixed port have root red degree in $\{1,5\}$. It cannot accept
arbitrary many marks. For a two-vertex core with one mixed mark
and one pure port at each vertex, a central edge of weight 0 forces
both red sums 1, weight 1 allows neither strong marked sum, and
weight 2 forces both red sums 5. No choice is legal.

The failure persists with four independent marked roots: take a
double-star core with centres $u,v$, two leaves at each centre,
and one mixed and one pure port at every leaf. A strong marked leaf
forces its core edge and its pure port to be both zero (red sum 1)
or both 2 (red sum 5). If the central edge has weight zero, safety
forces both centre red sums zero; if it has weight 2, safety forces
both centre sums six. Both cases conflict on the central edge.
If its weight is 1, each centre with two zero leaf edges has sum
1, conflicting with those leaves, and with two 2-edges has sum 5,
also conflicting. Its only possible choice is one of each, giving
sum three at both centres, again a central conflict. This is a
complete parameter argument; its finite assignment domain has
$3^5\cdot2^4=3888$ members. The weak two-mark lemma in A1 imposes
different hypotheses and permits certain unexposed roots to have
red sums 2 or 4. These strong-mark obstructions therefore do not
contradict the subcubic proof.

\section{Signed gaps and paired partitions (E1)}\label{sec:signedgap}

The constructions in this section use actual integer colour degrees.
In particular, a nonzero difference of four is useful, although it is
invisible modulo four. No assumption about an unknown initial colouring
is made. The results of Sections~\ref{sec:signedgap}--\ref{sec:rootports}
are independent of the subcubic proof.

\begin{lemma}[A signed gap with at most two occurrences]\label{lem:signedgap}
Let $L_i(y)=\sum_{e\in E_i}a_{ie}y_e$, where $a_{ie}\in\{-1,1\}$,
every $E_i$ contains at least four distinct variables, and every variable
belongs to at most two different sets $E_i$. There is a choice
$y_e\in\{-1,1\}$ for which $|L_i(y)|\ge2$ for every $i$.
Such a choice can be found in polynomial time.
\end{lemma}
\begin{proof}
Make an auxiliary vertex for each form. A variable occurring twice is an
edge between its two forms, with label $\tau_e=a_{ie}a_{je}$; parallel
edges are allowed, but there are no loops. A variable occurring once is
a boundary variable at its form.

Choose signs $\sigma_i\in\{-1,1\}$ and call an auxiliary edge good when
$\tau_e=\sigma_i\sigma_j$. Start with any signs and flip a vertex whenever
it has more bad than good internal edges. Each flip strictly increases
the total number of good edges, so there are at most as many flips as
internal edges. At termination, the number of good internal edges at
each vertex is at least its number $b_i$ of bad edges. Count boundary
variables as good and write $g_i$ for this enlarged good count. Then
\[
 m_i=g_i-b_i=|E_i|-2b_i\ge0.
\]
For a boundary variable put $y_e=\sigma_i a_{ie}$; for a good edge put
$y_e=\sigma_i a_{ie}=\sigma_j a_{je}$. A bad edge contributes opposite
values at its ends after multiplying the two rows by $\sigma_i$ and
$\sigma_j$. We may therefore select which end receives $+1$. If $o_i$
bad edges select $i$ as their positive end, then
\begin{equation}\label{eq:signedmargin}
 \sigma_i L_i=g_i-b_i+2o_i=m_i+2o_i.
\end{equation}
Call a vertex critical if $m_i\in\{0,1\}$. Since $|E_i|\ge4$, every
critical vertex has bad-edge degree at least two. In a component of the
bad-edge graph with a noncritical vertex, root a spanning tree at such
a vertex and select each nonroot vertex as the positive end of its
parent edge. If all vertices of a component are critical, its minimum
degree is at least two, so it contains a cycle, possibly a pair of
parallel edges. Select positive ends cyclically around that cycle and
use attached rooted trees to give every remaining vertex a positive
parent edge. Set the positive ends of unused edges arbitrarily.
Every critical vertex now has $o_i\ge1$, and every noncritical vertex
has $m_i\ge2$. Equation~\eqref{eq:signedmargin} proves the claim.
The local flips and all tree and cycle operations have polynomial bounds.
\end{proof}

Odd-length rows are permitted; their sums have odd parity, so their
absolute values are actually at least three. This distinction will
remove a boundary obstruction on an even-order path.

\begin{theorem}[Paired vertex partition]\label{thm:integerpairs}
Let $V(G)=S\sqcup T$ be a supplied partition. Suppose that both induced
graphs have maximum degree at most one, that $G[T]$ has no isolated
vertex, and that every edge $uv$ internal to either side satisfies
\[
 \min(d(u),d(v))\ge3,\qquad |d(u)-d(v)|\le1.
\]
Then $\lir(\doub G)\le2$. The mixed original edges can be exactly
$E(G[T])$. No degree bound is imposed on an internally isolated vertex
of $S$, and there is no upper bound on the maximum degree of $G$.
\end{theorem}
\begin{proof}
Give edges of $G[S]$ weight zero, edges of $G[T]$ weight one, and each
crossing edge $e$ weight $1+y_e$, where $y_e\in\{-1,1\}$. For each
internal pair $uv$ form
\[
 L_{uv}=\sum_{\substack{e\text{ crossing}\\e\ni u}}y_e
       -\sum_{\substack{e\text{ crossing}\\e\ni v}}y_e.
\]
The variables in a row are distinct and number $d(u)+d(v)-2\ge4$.
Each crossing edge occurs in at most one row at each end, hence at most
twice altogether. Apply Lemma~\ref{lem:signedgap}. The equal internal
contributions at $u,v$ give
\begin{equation}\label{eq:integerpairdifference}
 R(u)-R(v)=d(u)-d(v)+L_{uv},\qquad
 B(u)-B(v)=d(u)-d(v)-L_{uv}.
\end{equation}
Both differences are nonzero. On crossing edges, vertices of $S$ have
even red and blue degrees, whereas vertices of $T$ have odd red and
blue degrees. Thus all crossing edges are legal as well.
\end{proof}

The assumption on degree differences concerns the endpoints of the
internal pairs, not the degree range of the whole graph. The artifact
includes a connected nonbipartite example with a leaf and degrees
ranging from one to six.

\subsection{An integer repair with a fixed mixed set}

To see why retaining integer values matters, take the cube on
$0,\ldots,7$, with adjacency given by exclusive-or differences $1,2,4$,
and add $03$. Put $S=\{0,1,6,7\}$, $T=\{2,3,4,5\}$ and
$M=\{23,45\}$. The following weights and degrees are a complete witness.
\begin{center}\small
\begin{tabular}{c|rrrrrrrrrrrrr}
$e$&01&02&04&13&15&23&26&37&45&46&57&67&03\\\hline
$x_e$&0&2&2&0&0&1&2&0&1&2&0&0&0
\end{tabular}\par\medskip
\begin{tabular}{c|rrrrrrrr}
$v$&0&1&2&3&4&5&6&7\\\hline
$R(v)$&4&0&5&1&5&1&4&0\\
$B(v)$&4&6&1&7&1&5&2&6
\end{tabular}
\end{center}
With this same $M$, the stronger demand that both differences on every
mixed edge be nonzero modulo four is impossible: on $23$ the red degrees
are odd, while the ordinary degrees differ by one. A nonzero red
difference modulo four would have to be two, making the blue difference
zero modulo four. In the displayed legal colouring the actual red
difference is four and the blue difference is minus six.

There is also an infinite global obstruction to that stronger demand.
For $t\ge1$, take $H_t=C_{4t}$ with its antipodal matching. Subdivide
every edge once, let $S$ be the $4t$ old vertices and $T$ the $6t$ edge
vertices, and add a perfect matching inside each side. The resulting
graph $F_t$ has $10t$ vertices and $17t$ edges; its $S$ degrees are four
and its $T$ degrees are three. It is connected and nonbipartite: an
internal $S$ matching edge closes an even path in the subdivision into
an odd cycle. Theorem~\ref{thm:integerpairs} colours every $F_t$ with
exactly the $T$ matching mixed.

For the same mixed set, the stronger modulo-four condition is feasible
if and only if $t$ is even. Indeed, write
$R(v)\equiv p_v+2q_v\pmod4$, with $p=0$ on $S$ and $p=1$ on $T$.
Every internal pair must have opposite $q$ values. The $5t$ pairs force
$\sum_vq_v\equiv5t\pmod2$. On the other hand, writing every nonmixed
weight as $2z_e$ gives $q_v\equiv\sum_{e\ni v,\,e\notin M}z_e\pmod2$,
so the same sum is even. Conversely, when $t$ is even, assign opposite
$q$ values on each pair. Their even total can be realised as vertex
parities by a spanning-tree leaf elimination in the connected graph
$F_t-M$. The resulting weights satisfy the stronger condition.
Thus every odd $t$ supplies an integer repair on the same graph and
with the same mixed set, rather than a repair obtained by changing
the instance or the interface.

\section{Path and cycle sides in degree three and four (E2)}\label{sec:pathcyclepart}

Here and in Section~\ref{sec:isolatepart}, every ordinary degree is three
or four. An internal degree always means degree in an induced side of
the specified partition.

\begin{lemma}[A useful maximum-cut shape]\label{lem:biasedcut}
Every connected graph with $3\le d(v)\le4$ has a partition $S\sqcup T$
whose crossing graph is connected, such that $G[S]$ has maximum degree
two, every internally degree-two vertex has ordinary degree four,
and $G[T]$ is a matching together with isolated vertices.
\end{lemma}
\begin{proof}
Among maximum cuts choose one maximising the number of ordinary
degree-four vertices in $S$. A vertex cannot have more than half of
its neighbours on its own side. Thus ordinary degree-three vertices
have internal degree at most one and degree-four vertices at most two.
If a degree-four vertex in $T$ had internal degree two, flipping it
would preserve the cut size and improve the secondary objective.
Consequently $G[T]$ has maximum degree one.
If the crossing graph were disconnected, some original same-side edge
would join two of its components, since $G$ is connected. Flipping all
vertices of one such component loses no crossing edge and gains all its
boundary edges, contrary to maximum cut size.
\end{proof}

This existence argument does not supply a polynomial-time maximum-cut
algorithm. The colouring theorems below assume their partitions are
given. They do not prove that the isolated vertices permitted by
Lemma~\ref{lem:biasedcut} can always be handled.

\begin{theorem}[Path/cycle side against a matching]\label{thm:pathcyclepart}
Suppose $d(v)\in\{3,4\}$ and a supplied partition $S\sqcup T$ satisfies:
$G[T]$ is a matching covering $T$, $\Delta(G[S])\le2$, and every
internally degree-two vertex of $S$ has ordinary degree four.
Then $\lir(\doub G)\le2$, with mixed set exactly $E(G[T])$.
The colouring can be constructed in polynomial time.
\end{theorem}

We prove the local completion facts first. Give every crossing edge
weight $1+y_e$ with $y_e\in\{-1,1\}$, every $T$ matching edge weight one,
and each $S$ internal edge weight $2z_e$ with $z_e\in\{0,1\}$.
The crossing comparisons are automatically separated by parity.
For every $T$ pair impose the row condition $|L_{uv}|\ge2$ from
\eqref{eq:integerpairdifference}. It suffices to add a matching of
selected $S$ edges to these rows, then complete the internal pure edges.

\subsection{Binary relations}

A relation on two bits has full projections if either variable, when
fixed, has some compatible value for the other. Equality, inequality,
the exclusion of one ordered pair, and the unrestricted relation all
have this property. On a variable path, full projections permit an
extension from any one prescribed bit. On a variable cycle consisting
only of equality and inequality relations, consistency is equivalent
to an even number of inequalities. A variable cycle of full-projection
relations is also consistent if at least one relation allows three or
four pairs. To prove the last claim, compose the relations around the
cycle. Full projections are preserved under composition; composing a
relation with at least three pairs with any full-projection relation
still gives at least three pairs. The resulting relation on the first
variable contains a diagonal pair. These statements include a two-edge
cycle of parallel relations.

\subsection{Completing a cycle}

Let an $S$ cycle be $v_0\ldots v_{n-1}v_0$ and put $e_i=v_iv_{i+1}$,
with indices modulo $n$. Each vertex has two crossing edges. Let $s_i$
be its number of wholly red crossing edges. Then
\[
 R(v_i)/2=s_i+z_{i-1}+z_i,
 \qquad B(v_i)/2=4-s_i-z_{i-1}-z_i.
\]
Whatever the colour of $e_i$, its legality is exactly
\begin{equation}\label{eq:cyclebits}
 s_i-s_{i+1}+z_{i-1}-z_{i+1}\ne0.
\end{equation}
When the first difference is zero this is inequality of the two bits;
when it is $1$ or $-1$ it excludes one ordered pair; when its absolute
value is at least two it imposes no restriction. The variable graph
uses step two around the cycle. For odd $n$ it is one odd cycle; for
even $n$ it consists of two cycles of length $n/2$, including two-edge
cycles when $n=4$. The binary relation facts give the exact alternatives:
\begin{itemize}
\item If $n$ is odd, completion fails exactly when all $s_i$ are equal.
\item If $n\equiv0\pmod4$, completion is always possible.
\item If $n\equiv2\pmod4$, each of the two variable cycles must contain
at least one constraint with $s_i\ne s_{i+1}$.
\end{itemize}
It is therefore enough to select $e_0$ on an odd cycle, $e_0,e_3$ on a
cycle of length two modulo four, and no edge otherwise. In the second
case $n\ge6$, the two selected edges are disjoint and control different
variable cycles. For a selected edge, the crossing row is
$L_i=2(s_i-s_{i+1})$, with four distinct variables. The gap condition
ensures the required unequal values.

\subsection{Completing a path}

Let an $S$ path have vertices $v_0,\ldots,v_{n-1}$ and $m=n-1$ edges,
and write $z_{-1}=z_m=0$. Then
\[
 R(v_i)/2=s_i+z_{i-1}+z_i,\qquad
 B(v_i)/2=d(v_i)-s_i-z_{i-1}-z_i.
\]
An isolated vertex needs no completion. For a two-vertex path, if the
ordinary endpoint degrees differ, at least one of the two pure colours
is legal: the two failures would require simultaneously
$s_0-s_1=0$ and $s_0-s_1=d(v_0)-d(v_1)$. If the degrees are equal,
select its only edge and require $s_0\ne s_1$.

For $n\ge3$, every nonendpoint has ordinary degree four. Edges between
such vertices again give~\eqref{eq:cyclebits}, producing two variable
paths, one on each parity of edge index. At a left endpoint of ordinary
degree four, put $\delta=s_0-s_1$. Its condition is
$\delta-z_1\ne0$: it fixes $z_1=1$ for $\delta=0$, fixes $z_1=0$ for
$\delta=1$, and otherwise imposes no condition. At a degree-three left
endpoint, its own bit $z_0$ must also be considered:
\begin{equation}\label{eq:threeendbits}
 z_0=1\Longrightarrow\delta-z_1\ne0,
 \qquad z_0=0\Longrightarrow\delta-z_1\ne-1.
\end{equation}
For $\delta=0$ this is equality of $z_0,z_1$; for $\delta=-1$ it
excludes $(0,0)$; for $\delta=1$ it excludes $(1,1)$; otherwise it is
unrestricted. In each case it is a full-projection relation joining the
two variable paths. The right endpoint has the same description after
reversing the indices. Thus:
\begin{center}\small
\begin{tabular}{ccl}
\toprule Endpoint degrees&$n$&Sufficient treatment\\\midrule
$3,4$&either parity&One variable path, at most one prescribed bit.\\
$4,4$&odd&Prescribed bits lie on different variable paths.\\
$4,4$&even&Require left $\delta\notin\{0,1\}$.\\
$3,3$&odd&A relation cycle with an even number $m-2$ of inequalities.\\
$3,3$&even&Require left $\delta\ne0$ to make a relation nonbijective.\\
\bottomrule
\end{tabular}
\end{center}
In the odd-order $3,3$ case, either a relation is nonbijective, or both
endpoint relations are equality and all $m-2$ others are inequality.
This includes the parallel endpoint relations when $n=3$.

Select the first internal edge only when $n$ is even and the endpoint
degrees are equal. For endpoints of degree three and $n\ge4$, the row
has four variables and is $L=2\delta$. For endpoints of degree four,
it has five variables and is $L=2\delta-1$, since the first vertex has
three crossing edges and the second has two. The bound $|L|\ge2$ in
the latter case actually gives $|L|\ge3$ and excludes $\delta=0,1$.
The two-vertex case was handled separately. Each path selects at most
one internal edge.

\begin{proof}[Completion of Theorem~\ref{thm:pathcyclepart}]
All selected $S$ edges form a matching. Combine their crossing rows
with the rows for the $T$ pairs. Each row has at least four distinct
variables. Each crossing variable appears at most once on each side,
because the selected edges and the $T$ pairs are matchings.
Lemma~\ref{lem:signedgap} supplies all crossing signs simultaneously.
The preceding completion arguments then colour each $S$ component.
The $T$ comparisons follow from~\eqref{eq:integerpairdifference}, and
all crossing edges are separated by parity. Local binary completion is
equivalently a two-literal clause system, solvable by strongly connected
components, so the whole construction is polynomial.
\end{proof}

\subsection{Why internal recolouring is a real extension}

Take $2t$ disjoint triangles as $S$ and a perfect matching $Q$ between
their vertices, with no pair inside a triangle, connecting the triangles
as units. Such a matching exists for every $t\ge1$: allocate the three
ports of each triangle according to a connected cubic loopless
multigraph on $2t$ units. For each pair $uv\in Q$ introduce $a,b\in T$
and edges $ua,ub,va,vb,ab$. This gives a connected nonbipartite graph
with $12t$ vertices and $21t$ edges, of degrees four on $S$ and three
on $T$. Theorem~\ref{thm:pathcyclepart} applies with all $ab$ mixed.

If all $S$ edges were blue while this mixed set stayed fixed, each
$K_{2,2}$ crossing block would require a vertex of $S$ with exactly one
red crossing edge. Otherwise both $S$ vertices send equal contributions
to the two ends of the mixed edge $ab$. The $3t$ blocks use disjoint
$S$ pairs and require at least $3t$ such vertices. In a wholly blue
triangle there can be at most one: two would both have blue degree six
on their common blue edge. There are only $2t$ triangles, a contradiction.
The new theorem restores legality by changing internal pure colours,
while preserving the mixed set and the original graph.

The partition requirement itself is not universal. In $K_{3,3}$ plus
one chord on one side, an induced matching covering $T$ can contain at
most two vertices. Choosing the chord leaves an $S$ vertex of internal
degree three; choosing a crossing edge leaves two ordinary degree-three
vertices of internal degree two. Taking $T=\varnothing$ also fails.
This six-vertex graph has a stored legal colouring; it is an applicability
obstruction, not a counterexample to the conjecture.

\section{Isolated vertices of the matching side (E3)}\label{sec:isolatepart}

Keep $d(v)\in\{3,4\}$ and the same conditions on $G[S]$, but now allow
$G[T]$ to consist of a matching and an internally isolated set $I$.
Let $C$ be the crossing graph. Let $F$ be the selected $S$ matching from
the path and cycle proof, and let $P$ be its set of endpoints.

\begin{theorem}[A protected parity join]\label{thm:protectedjoin}
Put $A=C[(S\setminus P)\cup I]$. If every component of $A$ contains
an even number of vertices of $I$, then $\lir(\doub G)\le2$.
The original $T$ matching remains mixed, and the only additional mixed
edges belong to $A$. The component condition is necessary and sufficient
for an additional edge set in $A$ with odd degrees on $I$ and even
degrees on $S\setminus P$.
\end{theorem}
\begin{proof}
The last assertion is the elementary parity-subgraph criterion. Necessity
follows by summing degrees in each component. For sufficiency, root a
spanning tree in each component and determine parent edges from leaves
upwards; the root equation follows from the even number of odd targets.
Let $N$ be the resulting edge set and make exactly $N$ and the original
$T$ matching mixed. Every $T$ colour degree is odd and every $S$ colour
degree is even, so crossing edges are automatically legal.

No $N$ edge touches a $T$ pair or a protected $S$ endpoint. Consequently
the rows for $T$ pairs and selected $S$ edges keep all their original
pure crossing variables, with no new constants. They still have at
least four terms and at most two occurrences per variable. Apply
Lemma~\ref{lem:signedgap} and then the preceding component completion.
At an unprotected vertex the even number of mixed crossing edges still
makes its crossing red contribution even; the local completion only
needs the protected row conditions. This proves all required comparisons.
\end{proof}

A component of $A$ with an odd number of $I$ vertices is a certificate
against this protected-join restriction. It does not certify that the
graph has no locally irregular two-colouring.

For example, choose odd $t\ge1$ and $r\ge0$. On $2t$ auxiliary vertices
take a connected bridgeless cubic loopless multigraph: three parallel
edges for $t=1$, or $C_{2t}$ with an antipodal matching for $t\ge3$.
Subdivide one edge by $2r$ new vertices. Use these auxiliary vertices as
$T$, replace each auxiliary edge by a new $S$ vertex joined to its ends,
and join the $3t+2r$ vertices of $S$ into one path. Pair the $2r$ new
$T$ vertices consecutively by added internal edges. This is a simple
graph with $5t+4r$ vertices and $9t+7r-1$ edges. Its $S$ path has odd
order, so $F=\varnothing$; the auxiliary graph with the designated edge
removed remains connected. Thus the main component of $A$ contains all
$2t$ isolated $T$ vertices and the theorem applies. For $r>0$, internally
paired and isolated $T$ vertices coexist with unbounded counts.

\subsection{An exact condition with an independent side}

Suppose now that $T$ is independent. Fix a crossing mixed set $N$ with
\begin{equation}\label{eq:fixedjoinparity}
 d_N(v)\equiv0\pmod2\ (v\in S),\qquad
 d_N(v)\equiv1\pmod2\ (v\in T),
\end{equation}
and require all other original edges to be pure. Call an $S$ component
critical if it is a cycle of length not divisible by four, or an
even-order path with both ordinary endpoint degrees three. All its
vertices have crossing degree two. Call it saturated when every one
of its vertices has $N$ degree two.

\begin{theorem}[Exact independent-side completion]\label{thm:exactindependent}
For the fixed $N$ satisfying~\eqref{eq:fixedjoinparity}, a pure completion
of all remaining edges exists if and only if no critical component is
saturated.
\end{theorem}
\begin{proof}
Let $c_v$ be crossing degree and $m_v=d_N(v)$. The crossing red
contribution divided by two can be selected independently at each $S$
vertex as any integer in
\begin{equation}\label{eq:crossinterval}
 s_v\in[m_v/2,c_v-m_v/2]\cap\mathbb Z.
\end{equation}
Independence holds because every crossing edge has only one $S$ end,
and $T$ has no internal comparison. At a nonisolated $S$ vertex the
possibilities are $(c,m)=(2,0),(2,2),(3,0),(3,2)$, giving respectively
$s\in\{0,1,2\},\{1\},\{0,1,2,3\},\{1,2\}$.

Saturation of a critical component forces $s_v=1$ throughout. On a
critical cycle, the step-two constraint graph has an odd cycle of
inequality relations. On an even-order path with degree-three ends,
the two endpoint equalities and the odd number of internal inequalities
are inconsistent. This proves necessity, including a two-vertex path.

For each unsaturated critical component choose one vertex with $m_v=0$,
set its $s$ to zero and all other $s$ values to one. On an odd cycle this
creates a nonzero neighbouring difference. On a cycle of length two
modulo four the two differences around the special vertex occur in
different variable cycles. On a path with degree-three ends the change
makes an endpoint or internal relation nonbijective (and directly
separates the two vertices when the path has order two). The earlier
relation argument gives a completion.

For a cycle of length zero modulo four, an odd-order path, or a path
with different endpoint degrees, take all $s_v=1$. An isolated vertex
can take any allowed value. The only remaining case is an even-order
path with both endpoint degrees four. Its endpoints can each take
values one and two. For order two choose these different values. For
order $n\ge4$, set the left endpoint and all internal values to one.
Set the right endpoint to one if $(n-4)/2$ is even, and to two otherwise.
The two unary constraints on one variable path are separated by exactly
$(n-4)/2$ inequalities, so this choice is consistent; the other variable
path has no conflicting unary constraints. Realise the chosen crossing
degrees using~\eqref{eq:crossinterval}, then complete the internal
components. All crossing edges satisfy parity separation.
\end{proof}

\subsection{A bridge obstruction and its repair}

The opposite parity configuration, developed more generally in
Section~\ref{sec:rootports}, gives the following useful special case.

\begin{corollary}[Even components with an independent opposite side]
\label{cor:evencomponents}
If $T$ is independent and every $S$ component is an even-order path or
an even cycle, under the degree-three/four assumptions above,
$\lir(\doub G)\le2$. There is no parity restriction on $|T|$.
\end{corollary}
\begin{proof}
Make an internal perfect matching of every $S$ component mixed, and
all other internal edges blue. Bipartition each component, and give its
two classes respectively zero and two wholly red crossing edges at each
vertex. There are at least two crossing edges at every vertex. Red
degrees are one and five; every internal red difference has absolute
value four, and its blue difference is nonzero because ordinary
degrees differ by at most one. All crossing edges are pure, joining
odd colour degrees in $S$ to even colour degrees in $T$.
\end{proof}

An infinite family shows why parity solvability alone does not guarantee
the first configuration. For $t\ge1$, subdivide one edge of
$C_{4t}$ plus its antipodal matching, obtaining a graph $Q_t$ of order
$4t+1$ with one degree-two vertex. Join three disjoint copies of $Q_t$
to a new centre through these vertices. The resulting cubic graph $H_t$
has three bridges. Let $T=V(H_t)$ and introduce an $S$ vertex on each
edge of $H_t$, giving the subdivision as the crossing graph. Denote the
three bridge vertices in $S$ by $w_1,w_2,w_3$. Add $w_1w_2$, make six
nonbridge $S$ vertices into a six-cycle, and pair all remaining $S$
vertices. There is an even number remaining, and there are enough
nonbridge vertices for every $t\ge1$.

This simple connected nonbipartite graph has $30t+10$ vertices and
$45t+18$ edges; exactly six vertices have degree four, and all others
degree three. Every $S$ component has even order, so
Corollary~\ref{cor:evencomponents} colours it. But if all $S$ internal
edges are pure, $S$ has even colour degrees, and $T$ odd colour degrees,
each half of a subdivided bridge is forced mixed. Indeed, its cut side
contains $4t+1$ odd targets in $T$, so summing the mixed degrees modulo
two forces the sole allowed crossing boundary edge into $N$. Thus the
crossing red contribution is two at both $w_1,w_2$. If their common
internal edge has weight $x\in\{0,2\}$, then
\[
 R(w_1)=R(w_2)=2+x,\qquad B(w_1)=B(w_2)=4-x.
\]
Its actual colour is illegal. The obstruction excludes that entire
parity configuration, not just a particular choice of $N$, while the
opposite configuration colours the same infinite family.

\section{Odd-component root ports at arbitrary degrees (E4)}\label{sec:rootports}

The following theorem removes the global degree restriction from
Corollary~\ref{cor:evencomponents} and permits odd paths, odd cycles and
internally isolated vertices. The needed port certificate is explicit.

\begin{theorem}[Root-port construction]\label{thm:rootports}
Let a finite simple graph have a supplied partition $S\sqcup T$ with
$T$ independent and $\Delta(G[S])\le2$. Suppose every nonisolated
vertex of $G[S]$ has at least two neighbours in $T$, and every $S$
internal edge $uv$ satisfies $|d(u)-d(v)|\le1$.

Select one root in each odd-order $S$ component and no root in an
even-order component. In an odd path, number vertices from an end and
select a root with even index; in a singleton, select its only vertex.
In an odd cycle, choose a root whose two neighbours are not both of
degree $d(r)+1$ and not both of degree $d(r)-1$. Let $\mathcal U$ be
the selected roots. Suppose a crossing edge set $N$ exists with:
\begin{itemize}
\item every edge of $N$ joins a root to $T$;
\item every $T$ vertex has even $N$ degree;
\item every internally nonisolated root has $N$ degree one or three;
\item every internally isolated root has positive odd $N$ degree.
\end{itemize}
Then $\lir(\doub G)\le2$, constructively. There is no bound on the
maximum degree of $G$ or on degrees in $T$; internally isolated
vertices of $S$ may include original leaves.
\end{theorem}
\begin{proof}
An eligible root exists on every odd cycle. Otherwise each vertex
would be a strict local minimum or maximum of the degree sequence,
with these two types alternating, which cannot close around an odd
cycle. Deleting the chosen root from an odd cycle leaves an even-order
path; deleting an even-index vertex from an odd path leaves two
even-order paths. Give each remaining path or even cycle an internal
perfect matching. Make all its matching edges and all edges of $N$
mixed; make every other $S$ internal edge blue.

Each nonroot $S$ vertex has one internal mixed edge and no mixed crossing
edge. Every root has an odd number of mixed crossing edges and no
internal mixed edge. Thus all $S$ colour degrees are odd, all $T$
colour degrees are even, and all crossing edges are automatically legal.

On each remaining even path or cycle, assign red degrees one and five
to the two bipartition classes. A nonroot has at least two pure
crossing edges, so these targets are realised by making zero or two
of them wholly red. Internal edges between nonroots have red difference
of absolute value four; their blue differences are nonzero because
ordinary degree differences are at most one.

Give each nonisolated root red degree three. If it has one port, use
one other crossing edge wholly red; if it has three ports, use no other
red crossing edge. The assumed two crossing neighbours ensure capacity
in the first case. For a root neighbour $u$, the internal edge $ru$ is
blue, and its blue-degree difference is respectively
\begin{equation}\label{eq:rootbluedifference}
 B(r)-B(u)=2(d(r)-d(u))-2\quad\text{if }R(u)=1,
 \qquad
 B(r)-B(u)=2(d(r)-d(u))+2\quad\text{if }R(u)=5.
\end{equation}
If $d(r)-d(u)=1$, label five is forced; if it is $-1$, label one is
forced; if it is zero, either works. On an odd path, the components
on the two sides of the root can be phased independently. On an odd
cycle, the two endpoints after root deletion must receive opposite
labels. The only inconsistent requests are that both must be one or
both must be five, exactly the two cases excluded by root eligibility.
Hence a compatible phase exists.

An internally isolated root needs no internal comparison: give all
its remaining crossing edges blue. Its colour degrees retain odd
parity, with no need for red degree three. Every crossing edge has
only one $S$ end, so all these choices are independent. The preceding
arguments cover every original edge and produce the colouring by
finitely many polynomial-time traversals.
\end{proof}

\begin{corollary}[An exact root parity test within capacity]\label{cor:rootparity}
Fix eligible roots $\mathcal U$, and let
$A=C[\mathcal U\cup T]$ in the crossing graph. A subgraph with odd
degrees on $\mathcal U$ and even degrees on $T$ exists if and only if
each component of $A$ contains an even number of roots. If every
internally nonisolated root has crossing degree at most four, this
condition guarantees the port set of Theorem~\ref{thm:rootports}.
\end{corollary}
\begin{proof}
Apply degree summation and spanning-tree parity recovery as in
Theorem~\ref{thm:protectedjoin}. An odd port count at a nonisolated
root of crossing degree at most four is necessarily one or three.
Internally isolated roots permit any positive odd count.
\end{proof}

For roots of larger crossing degree the parity condition is still exact
for odd/even edge selection, but it alone does not guarantee the
one-or-three capacity for nonisolated roots. A failed tree recovery with
five ports must not be called infeasibility. A supplied qualified port
set remains sufficient at arbitrary degrees.

Pairing roots with common neighbours in $T$ is a special case: give
each pair its two crossing ports, making all $T$ counts even. It is not
necessary. For odd $b\ge3$, take a vertex $s_0\in S$ and $b$ groups of
three $S$ vertices; in each group add a private triple of $T$ vertices,
join the two triples completely, and join $s_0$ to one chosen private
$T$ vertex. All $S$ components are singletons. The root adjacency graph
formed by common $T$ neighbours is $b$ copies of $K_4$ sharing $s_0$.
It has no perfect matching: deleting $s_0$ leaves $b$ odd components,
which cannot all be matched through a single vertex. Nevertheless the
root-port graph is connected and has $1+3b$ roots, an even number, so
Corollary~\ref{cor:rootparity} applies. Explicitly, in each group select
the edge from $s_0$, one other edge to that same private $T$ vertex,
and two edges from the other two roots to a common private $T$ vertex.
This also shows why internally isolated roots need arbitrary odd port
capacity. These particular examples are bipartite, whose ordinary
existence theorem is already known~\cite{GW23}; their purpose is to
separate two port criteria, not to claim a new bipartite result.

Finally, fixing the roots can still lose a legal colouring. Take
$S=\{0,1,2,3\}$, $T=\{4,5,6\}$ with ordered edges
\[
 04,15,24,26,35,36,23.
\]
The odd components give roots $0,1$; the permitted root-port graph
has two components each containing one root, so the root-only scheme
fails. Yet the weights $(1,1,1,2,1,0,0)$ give
\[
 R=(1,1,3,1,2,2,2),\qquad B=(1,1,3,5,2,2,2),
\]
which are legal edge by edge. Here the even component $2\! -\!3$
also receives crossing mixed ports. Thus the exact root parity test
does not solve general even-component transport or produce the required
partition for every graph.

\section{Evidence and limits of the integer partition results}\label{sec:partitionevidence}

The four new groups are not four independent solutions of the mother
conjecture. E1 contains a general signed tool and an unbounded-degree
paired-partition application. E2 enlarges the $S$ side to paths and
cycles within degrees three and four. E3 addresses some, but not all,
isolated vertices on the matching side and gives an exact fixed-interface
criterion when that side is independent. E4 handles odd components
under an arbitrary-degree root-port certificate; the even-component
existence corollary is retained as its simpler special case rather than
counted as an additional maximal theorem. Fixed mixed sets, port counts
and parity patterns are part of the statements, not interchangeable
notions of extension.

The implementation checkers test 647 signed systems, 15,757 local
path/cycle profiles against full binary enumeration, 52 independent-side
fixed-join domains, and 727 eligible fixed-root atlas domains. They also
check parameter families, the actual integer repairs and the explicitly
failed restricted strategies. Separate standard-library auditors,
without importing the constructors or solvers, rebuild the saved signed
sums, local truth tables, cuts and actual copy degrees. Across the four
auditors, the stored whole-graph witnesses account for 972 actual
copies: 100, 190, 258 and 424 in the respective archived source audits.
The release report supplies machine-readable per-audit counts
so that repeated or overlapping examples are not counted as distinct
new graph families. These are implementation and certificate audits;
the displayed proofs justify the unbounded parameter ranges.

The finite counterexamples to weakened signed-gap hypotheses are also
retained. With only three terms, antisymmetric incidence forms of $K_4$
would require every vertex to be a source or a sink to reach absolute
value at least two, forcing a bipartition. Allowing three occurrences
already fails for
\[
 L_1=y_0+y_1+y_2+y_3+y_4,\quad
 L_2=-y_0-y_1+y_2+y_3+y_4,\quad
 L_3=y_0+y_1-y_2-y_3+y_4.
\]
After a global sign change $L_1\ge3$, at most one variable is negative.
With none negative $L_2=1$; if the sole negative is in positions zero or
one then $L_3=-1$, otherwise $L_2=-1$. Arbitrary affine shifts fail too:
for antisymmetric incidence forms of $K_5$, $L_v+1=2d^+(v)-3$.
The condition $|L_v+1|\ge2$ allows only outdegrees $0,3,4$; at most one
vertex has outdegree zero, so their sum is at least twelve rather than
the ten edges of $K_5$. These are obstructions to proposed stronger
tools, not to locally irregular two-colourability of the doubled graphs.

No claim of exhaustive historical priority is made for these tools.
The known regular and bipartite cases remain attributed to their
original sources. General graphs still need a suitable partition,
compatible mixed-edge choices, and control of the remaining component
configurations. The unrestricted conjecture is not proved.

\section{Artifacts, verification status and remaining problem}\label{sec:release}

\subsection{Three distinct roles of computation}

The artifact separates complete proof kernels, explicit graph witnesses,
and bounded implementation tests. The first role is logically part of
the computer-assisted proofs. A1 uses the complete safe-core relation
tables, 163 sparse-core representatives, and a 4,256-case weak two-mark
kernel, with fixed-point and path-saturation arguments for unbounded
length. The clique interface uses 836 complete small-order offset bases
and its two-vertex induction. The decorated-path interface uses all
93 endpoint comparison types, all 96 transition rows, and the displayed
lower-bound and saturation identities. These domains are finite because
their equivalence reductions are proved, not because experiments stopped
at a chosen graph order.

The second role is independently checking submitted weight assignments.
The release verifier rebuilds red and blue degrees from the actual
edge-copy counts, including physical offsets when present, verifies all
required comparisons, checks any recorded final degrees, and confirms
preservation of supplied boundary weights. Large stored region witnesses
can therefore be audited without repeating their discovery search.
The 13-regular descent barrier is supported by a fixed integer table
whose row sums and neighbour counts are checked directly; its arbitrary
fibre sizes are justified by the explicit matching construction.

The third role is testing the constructive implementations on graphs,
offsets and lengths beyond the finite kernels. These checks include
small graph atlases, random compressed graphs, long paths and return
paths, prescribed-root cliques, and several complete dense-region
examples. They audit code, rather than replace any infinite proof.
Original research-run reports are retained separately from reports of
the new release run. The provenance ledger identifies frozen original
notes and the script dependencies. Optional factor searches use
NetworkX and sometimes Z3. Their timeouts are reported as unfinished
searches; a timeout is never treated as an infeasibility certificate.

\subsection{Reproducing the release}

The default complete-certificate command, from an extracted archive, is
\begin{center}\code{python -I -B verify.py}\end{center}
It uses the Python standard library. It invokes the A1 certificate
verifier and its independent local recalculation, then the additional
clique and decorated certificates, fixed integer tables, stored graph
witnesses and independent arithmetic checks. It also runs all four
portable E1--E4 audits without importing their constructors or solvers.
Those audits recount 972 actual copies in the saved whole-graph witnesses
and rederive the submitted rows, local relations, parity cuts and
strategy obstructions. Do not run Python with
optimisation enabled: the certificate checkers deliberately use
assertions. A JSON report can be written using the option
\texttt{-{}-output validation/certificates.json}.

For constructive implementation tests, install the pinned dependencies
in \code{requirements.txt} and run
\begin{center}\code{python -B run_checks.py}\end{center}
This runs the released construction checkers and saves a bounded report.
The A1 command-line constructor remains available as
\code{construct.py}, with examples and an independent copy-degree
checker in the root directory. Other constructors and their scope
checks are in \code{artifact/}; the coverage ledger points to each
public implementation. A method-not-applicable response is not a
counterexample to the mathematical conjecture.

The editable English source is \code{manuscript/main.tex} with included
section files and an explicit bibliography. It can be compiled by a
standard LaTeX installation or Tectonic. No user-specific absolute path
is required by the released proof data or verifier. The manifest records
SHA-256 digests of the released files. The final validation report records
fresh extraction, certificate execution, constructor checks, PDF rebuild
and source-reference diagnostics. Checks of the archive are local;
no repository deposit, Zenodo upload, DOI allocation, or journal
submission is implied.

\subsection{Use of AI and research status}

This manuscript and its research artifacts were developed with extensive
assistance from OpenAI Codex, including mathematical exploration, proof
drafting, programming, calculations, internal review and English editing.
Automated checks and internal reviews are documented as such; they
are not independent expert peer review. No Lean or other proof-assistant
verification is claimed.

The artifact attributes the pre-existing graph-family results and the
classical degree-list, factor, matching and interpolation tools used in
the proofs. Its grouping of results records this research programme's
contributions and boundaries; it is not an assertion that every grouped
lemma first appears here. A literature-screening record is supplied as
provenance, with no guarantee of exhaustive historical priority.
The manuscript, documentation and finite certificate data are licensed
under Creative Commons Attribution 4.0 International (CC BY 4.0).
The accompanying Python source code is licensed under the MIT licence.
The licence-scope file and full licence texts accompany the release.

The remaining mathematical problem is the original arbitrary-graph
conjecture. None of the positive results forces a general graph to have
a sparse, homogeneous, clique, parity-separated, high-degree or dense
region structure. Target optimisation is exact only after a target is
fixed, and the algebraic route lacks a general nonvanishing theorem.
The E1--E4 partition conditions and mixed-port certificates have not been
proved to exist for every graph. In particular, a parity-subgraph
solution can saturate a forbidden local component, and root-only joins
can exclude valid colourings using even components.
Minimal-counterexample restrictions still permit arbitrary large,
nonregular, minimum-degree-three cores. The explicit obstructions show
why several universal parity and monotone-repair assertions cannot
fill that gap. A1 covers the full subcubic range; the later conditions
extend beyond that range but do not combine into a proof for every graph.

\appendix
\section{Finite domains and certificate specification}\label{app:cert}

The proof certificates describe exact relations, not probabilistic
tests.  The five data files are listed in Table~\ref{tab:files}.
The marked-theta table supports a more general internal helper;
Corollary~\ref{cor:cactus} makes that helper unnecessary for the
single-mark step of Theorem~\ref{thm:main}.  It is included so that
the implementation's dependencies remain self-contained.

\begin{table}[H]\centering\small
\begin{tabular}{p{0.48\linewidth}p{0.43\linewidth}}
\toprule
Data file in \code{artifact/}&Contents\\
\midrule
\code{safe_cactus_interfaces.json}&Six message types, 43 node cases,
letter matrices and complete even product layers\\
\code{safe_all_cactus_interfaces.json}&The same data, odd product layers
and even-cycle rooted profiles\\
\code{sparse_bip_core_interfaces.json}&Nine states, six-step walks,
163 sparse-core representative witnesses\\
\code{weak_two_marked_interfaces.json}&Arm layers, 4,256 two-mark
witnesses and the embedded chain certificate\\
\code{marked_theta_interfaces.json}&608 marked-theta witnesses and the
chain certificate for the optional helper\\
\bottomrule
\end{tabular}
\caption{Finite certificate data.  Equality of entire domains and relation
sets is checked, rather than equality of recorded totals alone.}
\label{tab:files}
\end{table}

\subsection{Sparse-core representatives}

Put $I=\{1,\ldots,8\}$ and $J=\{4,6,8\}$.  Length parameters in
Table~\ref{tab:sparse} record only chains between forks or around
cycles; leaf-arm lengths are handled by the separate induction in
Lemma~\ref{lem:free}.  A fork is a vertex of core degree three.
An even cycle represented by two fork-to-fork paths has lengths $a,b$.

\begin{table}[H]\centering\small
\begin{tabular}{p{0.30\linewidth}p{0.51\linewidth}r}
\toprule
Shape&Complete parameter domain&Cases\\
\midrule
Three-arm tree&No chain parameter&1\\
Cycle with an arm&Cycle length $a\in J$&3\\
Double-star tree&Fork-to-fork length $b\in I$&8\\
Cycle attached to a fork&Cycle length $a\in J$, joining length $b\in I$&24\\
Cycle with two arms&$a,b\in I$, $a\le b$, $a+b$ even, $a+b\ge4$&19\\
Theta&$a\le b\le c$ in $I$, all of the same parity, at most one equals 1&36\\
Two cycles joined by a chain&Cycle lengths $a,b\in J$, joining length $c\in I$&72\\
\midrule
Total&&163\\
\bottomrule
\end{tabular}
\caption{Complete finite domain for the sparse bipartite core lemma.}
\label{tab:sparse}
\end{table}

Each record supplies one or two fork triples in $\{0,1,2\}^3$, with
their sums in $\{1,3,5\}$.  For every compressed chain the checker
starts at its given end state, applies \eqref{eq:T} once per internal
vertex and checks its other end state and final unequal-degree
comparison.  Each leaf arm is checked against the relation of possible
pure-port leaf completions.  The shape data specify how triples are
shared by chains, so consistency of all three incidences is explicit.

\subsection{The two-mark kernel}

Label the direct unmarked position \code{u0}.  Label arms of lengths
1,2 and at least 3 by \code{u1}, \code{u2}, \code{u3}.
Use \code{m0} for a direct marked position and
\code{m1},\ldots,\code{m6} for marked arm lengths 1 through 5 and
at least 6.  Write $U$ and $M$ for these four- and seven-element
sets of labels.  The complete domain is
\begin{equation}\label{eq:kernel}
 \left(\operatorname{Sym}^2(M)\times\operatorname{Sym}^2(U)
       \;\cup\; (M\times U)^2\right)\times\{1,3,5,7\}.
\end{equation}
The two pieces are kept as distinct canonical patterns.  Each record
contains the two side labels, the middle length, the two end states
$(x,r)$ and the actual choices for the two positions at each side.

For a direct position, its contribution is 0 or 2 if unmarked and
1 if marked.  A direct marked position requires its root red degree
not to equal 3.  For an arm position, its pair $(y,s)$ must belong
to the message set for that label and satisfy $s\ne r$ at the
distinguished root.  At either side, the root degree must equal
the middle end weight plus its two position contributions.  The
middle path must satisfy the recurrence \eqref{eq:T}, allowing actual
end red degrees 2 and 4.  These are all the checks imposed by
\code{weak_two_marked_certificate.py}.

The independent local verifier does not read these 4,256 witnesses.
It reconstructs the arm messages, forms every possible side state,
computes all middle-path relations and checks that each member of
\eqref{eq:kernel} has a witness.  This provides a separate calculation
of the finite existence claim, while the main checker validates the
portable witness data.

\subsection{Execution and trust boundary}

From a fresh extraction of the archive run
\begin{center}\code{python -B verify.py}\end{center}
The program rejects optimized Python execution, since the supplied
checkers use assertions.  It verifies the six-state closure, the
163-case sparse-core certificate, the 4,256-case two-mark certificate,
the optional 608-case marked-theta certificate and the fresh local
calculation.  It also rejects deliberately damaged copies of selected
certificate data in memory.  No source data is changed.

For a full graph example, install the optional dependency and run
\begin{center}
\code{python -m pip install -r requirements.txt}\\
\code{python -B construct.py examples/triangle.json}
\end{center}
The input gives an integer vertex list and unordered pairs of distinct
vertices.  Duplicate edges, disconnected inputs, degree above three
and $K_2$ are rejected.  The output lists red and blue multiplicities
and actual vertex degrees, which are recounted before output.
The optional argument \texttt{-{}-output} saves the resulting JSON to a
chosen path.

All calculations rely on ordinary exact Python execution.  They are
reproducible computation, not proof-assistant kernel checking.
Conversely, the claim of arbitrary size rests on the explicitly
proved classifications, finite-state inductions and global
compatibility arguments in the body, which cannot be replaced by
running a finite suite of graph examples.

\section{Additional finite certificates}\label{app:additional}

\subsection{Decorated path tables}\label{app:decorated}

Order the 24 states by degree $D=2$ then 3, then by outgoing
weight $p=0,1,2$ and increasing red degree. Encode a state subset
by the sum of powers $2^i$ at its indices, starting at zero. The
complete lower-bound families used in Theorem~\ref{thm:decorated}
are as follows:
\begin{flushleft}\small\ttfamily
$F_1$: 9, 16, 36, 65, 72, 130, 260, 288, 52736, 473088,\newline
1132032, 1149440, 1166336, 1688064, 2100224, 2671104,\newline
3217408, 3228160, 3768832, 3784704, 5313536, 6293504,\newline
6754304, 8599552, 8722432, 8792064, 8859648, 10686464,\newline
12914688, 15073280.\medskip\newline
$F_2$: 27, 45, 54, 130, 216, 325, 360, 432, 3264000,\newline
5313536, 7425536, 7588352, 10344960, 10984448, 11463168,\newline
12022272, 15014912, 15126016, 15144960.\medskip\newline
$F_3$: 219, 325, 438, 7458304, 11463168, 15126016, 15145984.\medskip\newline
$F_4$: 365, 12581376, 16244224.
\end{flushleft}
The three final masks are the signatures in the table below; each
cell records possible red degrees at the given outgoing weight.
\begin{center}\small
\begin{tabular}{ccccc}\toprule
Signature&$D$&$p=0$&$p=1$&$p=2$\\\midrule
$M_A$&2&$\{0,2\}$&$\{1,3\}$&$\{2,4\}$\\
$M_L$&3&$\{0,2,3,4\}$&$\{1,2,3,4,5\}$&$\{2,3,4,6\}$\\
$M_S$&3&$\{0,1,2,3\}$&$\{1,2,3,4,5\}$&$\{3,4,5,6\}$\\\bottomrule
\end{tabular}
\end{center}
The local full-state signatures are
\begin{align*}
 \Omega_A&=(\{0,1,2\},\{1,2,3\},\{2,3,4\}),&&D=2,\\
 \Omega_L=\Omega_F&=(\{0,1,2,3,4\},\{1,2,3,4,5\},\{2,3,4,5,6\}),&&D=3,\\
 \Omega_S&=(\{0,1,2,3\},\{1,2,3,4,5\},\{3,4,5,6\}),&&D=3.
\end{align*}
For every $M\in\{M_A,M_L,M_S\}$ and every letter $X$,
$T_X(M)=\Omega_X$. Every $\Omega_X$ contains one of the final
three masks. The machine-readable table contains all 372 initial
cases, 96 transition rows, 224 successive-family containment checks
and 12 saturation checks, with the exact state and endpoint domains.
The independent checker reconstructs the local relations rather
than trusting the recorded counts.

\subsection{Clique bases and additional witnesses}

The root clique table is indexed by $(n,A,B,S)$ for $n=4,5$,
the clipped physical domain in Theorem~\ref{thm:cliqueport}, root
label zero and lexicographic edge order. Every entry records a full
word in $\{0,1,2\}^{\binom n2}$ and the truth of the nonroot
interior bound. The independent checker verifies each bound itself
and permits weakness only at the two specified $K_4$ inputs. The
domain has 836 entries and 14,424 actual coloured copies.

The portable datasets include the regular-layer integer table, the
decorated-path transfer table, and the complete clique root bases.
Their file names and schemas are identified in \code{COVERAGE.md}.
The witness files for dense, offset and sequential regions instead
record complete actual graphs, weights, final colour degrees and,
where applicable, the fixed boundary edges and offsets. Their
verification checks a concrete construction and does not replace
the factor-theorem arguments proving arbitrary-size existence.

The standard-library release verifier additionally recalculates the
wheel's entire single-edge neighbourhood, the five-vertex path's
entire two-unit target neighbourhood, the strong-mark domain, and
the nine-vertex continuous/integer distinction using rational
arithmetic. It checks the frozen regular-layer table against the
explicit formula, and tests coefficient/permanent and modular
identities on small exact inputs. Its report records actual domain
sizes computed in that run. These are finite identities or witness
audits, with their roles kept separate from the unbounded proofs.

\subsection{Coverage and dependency files}

\code{COVERAGE.md} maps every retained inventory item and completed
post-17:00 note to its English proof, implementation and certificates.
\code{PROOF_DEPENDENCIES.md} lists the exact external statements used
and the difference between complete proof kernels and sampled graph
tests. \code{provenance/source_manifest.json} hashes the original
top-level source snapshot, and \code{provenance/source_scope.json}
records the scientific scope and snapshot date without internal chat
identifiers. Chinese source notes are
included for provenance; the English proofs do not require reading
them. No unfinished current-thread experiment is promoted to a theorem.

\begin{thebibliography}{99}
\raggedright\small
\bibitem{GW23}
I.~Grzelec and M.~Wo\'zniak,
On decomposing multigraphs into locally irregular submultigraphs,
\emph{Applied Mathematics and Computation} \textbf{452} (2023), 128049.
\href{https://doi.org/10.1016/j.amc.2023.128049}{doi:10.1016/j.amc.2023.128049}.
Preprint: \href{https://arxiv.org/abs/2208.08809}{arXiv:2208.08809}.

\bibitem{GW24}
I.~Grzelec and M.~Wo\'zniak,
Local irregularity conjecture for 2-multigraphs versus cacti,
\emph{Opuscula Mathematica} \textbf{44}(1) (2024), 49--65.
\href{https://doi.org/10.7494/OpMath.2024.44.1.49}{doi:10.7494/OpMath.2024.44.1.49}.
Preprint: \href{https://arxiv.org/abs/2211.08270}{arXiv:2211.08270}.

\bibitem{GOW26}
I.~Grzelec, A.~Onderko and M.~Wo\'zniak,
On local irregularity conjecture for 2-multigraphs,
\emph{Applied Mathematics and Computation} \textbf{514} (2026), 129832.
\href{https://doi.org/10.1016/j.amc.2025.129832}{doi:10.1016/j.amc.2025.129832}.
Preprint: \href{https://arxiv.org/abs/2412.04200}{arXiv:2412.04200}.

\bibitem{CR17}
D.~W.~Cranston and L.~Rabern,
Beyond degree choosability,
\emph{Electronic Journal of Combinatorics} \textbf{24}(3) (2017), P3.29.
\href{https://doi.org/10.37236/6179}{doi:10.37236/6179}.
Preprint: \href{https://arxiv.org/abs/1511.00350}{arXiv:1511.00350}.

\bibitem{ADR08}
L.~Addario-Berry, K.~Dalal and B.~A.~Reed,
Degree constrained subgraphs,
\emph{Discrete Applied Mathematics} \textbf{156}(7) (2008), 1168--1174.
\href{https://doi.org/10.1016/j.dam.2007.05.059}{doi:10.1016/j.dam.2007.05.059}.
\href{https://problab.ca/louigi/papers/dcsubgraphs.pdf}{Author version,
Theorem~5}.

\bibitem{Tutte54}
W.~T.~Tutte,
A short proof of the factor theorem for finite graphs,
\emph{Canadian Journal of Mathematics} \textbf{6} (1954), 347--352.
\href{https://doi.org/10.4153/CJM-1954-033-3}{doi:10.4153/CJM-1954-033-3}.

\bibitem{Alon99}
N.~Alon,
Combinatorial Nullstellensatz,
\emph{Combinatorics, Probability and Computing} \textbf{8}(1--2)
(1999), 7--29.
\href{https://doi.org/10.1017/S0963548398003411}{doi:10.1017/S0963548398003411}.
\href{https://web.math.princeton.edu/~nalon/PDFS/null2.pdf}{Author version}.

\bibitem{BC95}
A.~Bouchet and W.~H.~Cunningham,
Delta-matroids, jump systems, and bisubmodular polyhedra,
\emph{SIAM Journal on Discrete Mathematics} \textbf{8}(1) (1995), 17--32.
\href{https://doi.org/10.1137/S0895480191222926}{doi:10.1137/S0895480191222926}.
\end{thebibliography}

\end{document}
